\documentclass[lang=cn,11pt,scheme=chinese]{elegantbook}
\usepackage{longtable}
\usepackage{graphicx}
\addbibresource{tex/common/references.bib}
% Reader-facing categories for statements that are not theorems.
\clubpenalty=10000
\widowpenalty=10000
\displaywidowpenalty=10000

% Short transitions stay inside the prose rather than creating another heading level.
\providecommand{\MFQLead}[1]{\par\noindent\textbf{#1。}\enspace}
\providecommand{\MFQInterviewMeta}[5]{%
  \par\smallskip\noindent\textcolor{blue!55!black}{%
  \textbf{#1｜#2｜#3}；能力：#4；来源映射：#5。}\par\smallskip}

\newenvironment{heuristic}
  {\begin{tcolorbox}[breakable,colback=yellow!5,colframe=yellow!45!black,title=经验近似]}
  {\end{tcolorbox}}

\newenvironment{diagnostic}
  {\begin{tcolorbox}[breakable,colback=orange!5,colframe=orange!55!black,title=诊断式]}
  {\end{tcolorbox}}

\newenvironment{implementationnote}
  {\begin{tcolorbox}[breakable,colback=blue!4,colframe=blue!45!black,title=实现说明]}
  {\end{tcolorbox}}

% Unified visual language for both volumes.
\definecolor{MFQBlue}{HTML}{2B6CB0}
\definecolor{MFQGreen}{HTML}{2F855A}
\definecolor{MFQOrange}{HTML}{C05621}
\definecolor{MFQGray}{HTML}{718096}
\definecolor{MFQLightBlue}{HTML}{EBF8FF}
\definecolor{MFQLightGreen}{HTML}{F0FFF4}
\definecolor{MFQLightOrange}{HTML}{FFFAF0}
\definecolor{MFQLightGray}{HTML}{F7FAFC}

\newcommand{\MFQFigure}[4][0.94\linewidth]{%
  \begin{figure}[htbp]
    \centering
    \includegraphics[width=#1]{#2}
    \caption{#3\protect\newline{\footnotesize\color{MFQGray}图源：\texttt{#4}}}
    \label{fig:#4}
  \end{figure}%
}


\title{量化研究数学}
\subtitle{下册：方向模型与研究项目}
\author{Jason W.}
\providecommand{\MFQReleaseDate}{2026-07-31}
\ifdefined\MFQVersion\else
  \PackageError{math-for-quant}{MFQVersion is undefined}{Build with tools/publish.ps1 so VERSION is injected.}
\fi
\date{\MFQReleaseDate}
\version{\MFQVersion}
\hypersetup{pdftitle={量化研究数学：下册},pdfauthor={Jason W.},pdfsubject={方向模型与研究项目}}
\extrainfo{用数学解释模型，用模型分析问题。}
\cover{assets/cover.jpg}
% 目录只显示“章—节”两级；更细的行文提示不进入目录。
\setcounter{tocdepth}{1}

\begin{document}

\maketitle
\frontmatter

\chapter*{如何选择路线}
下册包含多因子与计量、时间序列与统计套利、机器学习 Alpha、衍生品、组合与风险、高频与执行六个学习方向。读者可以选择感兴趣的方向，借助前面的自测问题回顾所需基础，再依次学习模型、估计方法与综合项目。

\section*{如何使用本书}

这不是一本把公式按字典排列的手册。每章从一个量化研究中会真正遇到的问题出发，先说明为什么需要新的数学对象，再逐步建立定义、直觉和定理，随后用完整算例和反例检验结论的边界。读者应在每个主要推导处停下来，尝试解释每个假设的作用；如果某个假设在证明中没有出现，通常意味着命题可以变强，或者证明漏掉了关键一步。

上册的章节按照依赖关系排列，但不要求所有读者走同一条路线。希望尽快进入应用的读者，可以优先学习极限、线性代数、概率分布、统计推断、时间序列、优化、数值稳定性和 Monte Carlo；希望建立严格基础的读者，则从分析、测度和条件期望开始顺读。每章末尾的练习从概念解释逐渐过渡到证明、计算和研究判断，完整答案与常见误区收录在共享答案册中。

Notebook 展示数学对象如何一步步变成计算：先预测现象，再观察中间量与图形，最后修改参数并解释变化。正文中的小算例可用纸笔理解，复杂实验无需事先手算全部输出。解析结果帮助核对计算；改变数据后应重新推理，不必维持原来的小数。完整推导与练习反馈在共享答案册中。

本书只在这里解释一次阅读方式。项目维护者所需的文件布局、Notebook 命令、图形生成和出版流程另见仓库文档，不属于读者必须掌握的数学内容。


\tableofcontents

\part*{路线选择与共同原则}
\addcontentsline{toc}{part}{路线选择与共同原则}
\setcounter{secnumdepth}{-1}

\chapter{如何判断一项量化研究是否可信}

一个模型可以解释收益的共同变化，也可以预测未来观测、给期权定价或选择交易时机。这些问题的目标各不相同。阅读时先辨明模型中的未知量和可用信息，再考察估计与计算怎样得到答案，以及数据变化和交易成本怎样影响结论。

\section{六条路线的共同结构}

各方向的章节按概念之间的联系展开。方向开头的自测问题帮助回顾上册知识，末尾的综合项目则把推导、估计与数据分析放在同一个问题中。遇到尚未熟悉的概念，可以根据先修表返回相应章节。

\begin{note}
综合项目中的算例用于学习方法；合成数据上的结论需要与真实数据上的经验结果区分。
\end{note}

\section{按问题选择学习起点}

先尝试回答所选方向的问题。能够解释概念并完成相应计算的部分，可以继续向下学习；尚不熟悉的部分，则利用后面的先修表返回上册。部分问题会在本方向中进一步展开，首次阅读时不必全部答出。

\newcommand{\RouteScoreInterpretation}{}

各方向与经典资产定价、时间序列和统计研究方法相联系；相应参考文献也可用于深入阅读\cite{fama1973risk,brockwell2016timeseries,white2000reality,bailey2017backtest}。

\section{多因子与计量}\label{route:multifactor}

\begin{enumerate}
  \item 能否解释横截面回归的系数、残差与暴露分别表示什么？
  \item 能否区分时间序列回归与 Fama--MacBeth 两步估计？
  \item 能否说明异方差、时间相关和横截面相关分别破坏哪类标准误？
  \item 能否用投影矩阵解释行业和市值中性化？
  \item 能否说明协方差病态为何放大组合权重？
  \item 能否写出 ridge 与 lasso 的目标函数及其差异？
  \item 能否为因子搜索设计多重检验控制？
  \item 能否定义时点数据并识别一次发布日期泄漏？
  \item 能否把换手、佣金和冲击写入净收益？
  \item 能否解释样本外因子表现与交易规模之间的关系？
\end{enumerate}
\RouteScoreInterpretation

\section{时间序列与统计套利}\label{route:stat-arb}

\begin{enumerate}
  \item 能否区分弱平稳、严格平稳与遍历性？
  \item 能否判断 AR 模型的因果性和单位根风险？
  \item 能否解释白噪声、鞅差创新与 IID 噪声的强弱关系？
  \item 能否说明协整不等于高相关？
  \item 能否从误差修正模型解释短期偏离与长期约束？
  \item 能否解释 OU 半衰期估计的模型依赖？
  \item 能否写出 Kalman 滤波的预测与更新步骤？
  \item 能否设计不泄漏未来的 walk-forward 划分？
  \item 能否处理配对选择带来的多重检验？
  \item 换手、未成交和结构变化会怎样影响策略收益？
\end{enumerate}
\RouteScoreInterpretation

\section{机器学习与 AI Alpha}\label{route:ml-alpha}

\begin{enumerate}
  \item 能否从经验风险最小化解释训练损失与泛化误差的差别？
  \item 能否解释偏差--方差权衡而不把模型复杂度等同于有效性？
  \item 能否写出 $L^1$ 与 $L^2$ 正则化目标？
  \item 能否说明时间序列中随机交叉验证为何泄漏？
  \item 能否设计 purge、embargo 与嵌套验证？
  \item 能否区分概率校准、排序能力与组合收益？
  \item 能否检查特征可得时间与标签重叠？
  \item 能否解释 concept drift 并设置监测统计量？
  \item 能否评估特征重要性的跨窗口稳定性？
  \item 能否将预测转换为含成本、约束和容量的组合决策？
\end{enumerate}
\RouteScoreInterpretation

\section{衍生品定价与对冲}\label{route:derivatives}

\begin{enumerate}
  \item 能否说明条件期望、滤过与信息时点的关系？
  \item 能否定义鞅并说明可选停止需要条件？
  \item 能否说明 Brownian 运动的独立增量与二次变差？
  \item 能否从简单复制组合解释无套利定价？
  \item 能否区分真实测度与风险中性测度？
  \item 能否说明 It\^o 公式为何含二阶项？
  \item 能否把定价问题联系到 PDE 或条件期望？
  \item 能否检查 Monte Carlo 估计误差和方差缩减？
  \item 能否识别隐含波动率校准的病态性？
  \item 能否报告离散对冲、交易成本和模型错设的影响？
\end{enumerate}
\RouteScoreInterpretation

\section{组合与风险}\label{route:portfolio-risk}

\begin{enumerate}
  \item 能否从二次型解释组合方差与边际风险？
  \item 能否诊断样本协方差矩阵的条件数？
  \item 能否解释收缩估计的偏差--方差交换？
  \item 能否写出带约束的均值--方差问题和 KKT 条件？
  \item 能否区分风险平价、最小方差与均值--方差目标？
  \item 能否解释 Black--Litterman 中先验、观点和置信度？
  \item 能否定义 VaR、CVaR 与最大回撤并说明局限？
  \item 能否设计压力情景而不把历史最差日当作全部尾部风险？
  \item 能否把换手惩罚与再平衡频率写入优化？
  \item 估计误差、约束变化和交易规模会怎样改变组合？
\end{enumerate}
\RouteScoreInterpretation

\section{高频、微观结构与执行}\label{route:microstructure}

\begin{enumerate}
  \item 能否区分日历时间、事件时间与交易时间？
  \item 能否从订单簿定义 bid、ask、spread 与 midprice？
  \item 能否区分限价单、市场单、撤单与部分成交？
  \item 能否用条件强度解释 Poisson 或 Hawkes 订单流？
  \item 能否说明排队位置和延迟如何改变成交概率？
  \item 能否区分暂时冲击与永久冲击？
  \item 能否写出执行成本与库存风险之间的权衡？
  \item 能否设计事件驱动仿真并避免偷看后续订单？
  \item 能否把未成交、拒单和数据丢包纳入反例实验？
  \item 能否报告延迟、流动性、容量与制度规则的适用边界？
\end{enumerate}
\RouteScoreInterpretation

\section{各方向需要的上册知识}

下表列出各方向用到的上册章节。某一道题涉及的概念尚不熟悉时，可以先回顾相应定义与例题，无须重读整册。

% Generated from curriculum/manifest.json; do not edit by hand.
\begin{longtable}{p{0.22\textwidth}p{0.68\textwidth}}
\toprule
路线 & 需要核对的上册学习单元 \\
\midrule
\endfirsthead
\toprule
路线 & 需要核对的上册学习单元 \\
\midrule
\endhead
多因子与计量 & \texttt{upper.ch06}（线性代数、分解、投影与条件数）；\newline \texttt{upper.ch10}（估计、检验、渐近、回归与多重检验）；\newline \texttt{upper.ch13}（凸优化、KKT 与对偶）；\newline \texttt{upper.ch14}（浮点数、数值线代与病态问题）；\newline \texttt{upper.ch16}（研究有效性、样本外与交易摩擦）；\newline \texttt{upper.ch17}（可复现研究审计） \\
时间序列与统计套利 & \texttt{upper.ch09}（概率收敛、极限定理与集中不等式）；\newline \texttt{upper.ch11}（随机过程、Markov、鞅、Poisson 与 Brownian）；\newline \texttt{upper.ch12}（时间序列与状态空间基础）；\newline \texttt{upper.ch14}（浮点数、数值线代与病态问题）；\newline \texttt{upper.ch15}（Monte Carlo、bootstrap 与方差缩减）；\newline \texttt{upper.ch16}（研究有效性、样本外与交易摩擦）；\newline \texttt{upper.ch17}（可复现研究审计） \\
机器学习与 AI Alpha & \texttt{upper.ch05}（多元微积分与矩阵微分）；\newline \texttt{upper.ch06}（线性代数、分解、投影与条件数）；\newline \texttt{upper.ch09}（概率收敛、极限定理与集中不等式）；\newline \texttt{upper.ch10}（估计、检验、渐近、回归与多重检验）；\newline \texttt{upper.ch12}（时间序列与状态空间基础）；\newline \texttt{upper.ch14}（浮点数、数值线代与病态问题）；\newline \texttt{upper.ch15}（Monte Carlo、bootstrap 与方差缩减）；\newline \texttt{upper.ch16}（研究有效性、样本外与交易摩擦）；\newline \texttt{upper.ch17}（可复现研究审计） \\
衍生品定价与对冲 & \texttt{upper.ch07}（概率空间、随机变量与分布）；\newline \texttt{upper.ch08}（条件期望、独立性与概率核）；\newline \texttt{upper.ch09}（概率收敛、极限定理与集中不等式）；\newline \texttt{upper.ch11}（随机过程、Markov、鞅、Poisson 与 Brownian）；\newline \texttt{upper.ch13}（凸优化、KKT 与对偶）；\newline \texttt{upper.ch14}（浮点数、数值线代与病态问题）；\newline \texttt{upper.ch15}（Monte Carlo、bootstrap 与方差缩减）；\newline \texttt{upper.ch16}（研究有效性、样本外与交易摩擦）；\newline \texttt{upper.ch17}（可复现研究审计） \\
组合与风险 & \texttt{upper.ch06}（线性代数、分解、投影与条件数）；\newline \texttt{upper.ch10}（估计、检验、渐近、回归与多重检验）；\newline \texttt{upper.ch13}（凸优化、KKT 与对偶）；\newline \texttt{upper.ch14}（浮点数、数值线代与病态问题）；\newline \texttt{upper.ch16}（研究有效性、样本外与交易摩擦）；\newline \texttt{upper.ch17}（可复现研究审计） \\
高频、微观结构与执行 & \texttt{upper.ch07}（概率空间、随机变量与分布）；\newline \texttt{upper.ch08}（条件期望、独立性与概率核）；\newline \texttt{upper.ch09}（概率收敛、极限定理与集中不等式）；\newline \texttt{upper.ch11}（随机过程、Markov、鞅、Poisson 与 Brownian）；\newline \texttt{upper.ch12}（时间序列与状态空间基础）；\newline \texttt{upper.ch14}（浮点数、数值线代与病态问题）；\newline \texttt{upper.ch15}（Monte Carlo、bootstrap 与方差缩减）；\newline \texttt{upper.ch16}（研究有效性、样本外与交易摩擦）；\newline \texttt{upper.ch17}（可复现研究审计） \\
\bottomrule
\end{longtable}


\section{怎样开展综合练习}
综合练习把一个明确问题贯通到结果：先建立模型与假设，完成小算例，再用新的数据或情景观察结论是否延续。每条路线只需要选择一个问题，不必把所有方法放进同一项目。

例如“一个因子是否预测下一期收益”，应先解释该因子在何时可得，再计算预测误差；如果进一步讨论交易，才需要定义仓位、成交和成本。模型预测正确、组合盈利和未来持续有效，是不同强度的结论。

完成后写一段自己的解释：最关键的推导是哪一步？哪个例子改变了原来的直觉？若去掉一条假设，会出现什么反例？答案应指向具体公式和观察，不以材料数量或固定输出判断理解。


\mainmatter
\setcounter{secnumdepth}{1}
\part{通用笔面试训练}
\chapter{Brainteaser：把陌生题压缩成可证明结构}
\label{ch:lower-brainteasers}

\begin{itemize}
  \input{tex/generated/prerequisites/lower-brainteasers}
  \item 学习产出：能把陌生题压缩成状态、不变量、下界和构造，并在有限时间内给出可核验的口述或笔试答案。
  \item 研究边界：谜题表现不能替代方向模型、样本外验证与交易摩擦证据。
\end{itemize}

Brainteaser 不是猜谜比赛。量化研究笔面试真正考察的是：能否在信息不完整、时间有限时，迅速选择状态、写出不变量、给出下界，并说明答案依赖哪些假设。本章把这类能力拆成五种可复用动作；题目来自《量化金融面试实用指南》第二章的核心能力，但题面、参数和解答均重新设计。

\section{先缩小状态，再开始计算}

遇到海盗分配、动物过河或多人博弈，第一问不是“有哪些招”，而是“决定未来的最小状态是什么”。若历史只通过当前人数、轮次或余数影响以后，就把它压缩成状态 $s$；若目标是保证成功，则从最小规模向后归纳。

\begin{example}
五位研究员按资历提议分配奖金，方案通过需要至少一半票数，否决者离场。不要枚举所有分配；从一人、两人局面逆推，每一步先识别哪些人的下一轮收益最低，再用最少成本买票。
\end{example}

Python 中可以递归枚举小规模局面，但面试答案必须先说明状态、终止条件和参与者的偏好顺序；否则程序只是在替一个未定义的博弈求解。

\section{逻辑约束：把语言变成可排除的集合}

真假标签、称球和过桥题适合先声明概率空间 $(\Omega,\mathcal F,\mathbb P)$，再把可行状态写成候选集合 $\Omega$。一次观察把它切成若干分支；最坏分支的大小决定信息效率。若一次实验最多有 $b$ 种结果，$k$ 次实验区分 $N$ 个状态至少需要 $b^k\ge N$。这个下界先告诉你“几次不可能”，再指导构造。

\begin{example}
若天平一次有左重、平衡、右重三种结果，两次称量最多编码九个叶节点。因此想在十二种互斥状态中保证定位，两次称量从信息量上已不可能。
\end{example}

\section{不变量、对称与奇偶性}

硬币堆、变色球、开关和蚂蚁题经常故意诱导逐步模拟。更短的路线是寻找每次操作都保持的量：总和模 $m$、逆序数奇偶性、颜色计数差，或在交换标签后不变的轨迹集合。对称不是“看起来一样”，而是存在保持概率或可达性的双射。

\begin{example}
若每次操作同时翻转两个开关，则“亮灯数的奇偶性”保持不变。于是从偶数盏亮灯出发，不可能到达奇数盏亮灯的目标；这比搜索所有开关序列更强。
\end{example}

\section{条件概率与期望：先声明信息集}

生日、卡牌与错标袋问题的陷阱通常不是算术，而是条件事件含糊。先写样本空间和已知信息，再计算 $P(A\mid B)$；若过程有多轮，优先尝试全期望公式或指示变量线性性。不要因为随机变量相关，就误以为其指示变量不能逐项求期望。

\section{递推、博弈与停止规则}

轮流取物、囚徒策略和序列题要区分三件事：状态转移、双方信息、停止条件。先计算小规模表格寻找模式，再用归纳或模运算证明；小规模程序只是猜想生成器，不是证明。

\section{精选题：五种方法}

以下每题只有一个主归属，交叉章节只可引用，不重复计入覆盖率。四级训练覆盖\textbf{口述概念、笔试推导、数值编程、研究判断}；建议先口述建模，再在纸上给出最坏情形证明。

\subsection*{口述概念}
限时解释状态、信息与关键不变量：GB-2-01、02、05、11--14、17--19、21、23、26、33、35、36。

\subsection*{笔试推导}
给出下界、构造与完备性证明：GB-2-03、04、06--08、15、16、20、24、25、27、28、30--32、34。

\subsection*{数值编程}
将闭式解或统计量恢复写成可验证程序：GB-2-10、22。

\subsection*{研究判断}
在信息、成本或协议限制下评估结论是否成立：GB-2-09、29。

下面按逆向归纳、信息下界、不变量、条件概率和递推五种方法组织精选题。来源映射和完整覆盖台账保存在附录与项目报告中；读者首先需要看到的是题目要求掌握的思路。

\begin{enumerate}[leftmargin=*]
  \item \MFQInterviewMeta{口述}{12 分钟}{难}{逆向归纳与博弈均衡}{GB-2-01 Screwy pirates}六位合伙人依次提出整数奖金分配；提案需严格过半，否决者退出。所有人依次偏好生存、奖金、让别人退出。写出从一人局面开始的逆推算法，并说明平票规则改变时哪里必须重算。
  \item \MFQInterviewMeta{口述}{8 分钟}{中}{状态机与安全策略}{GB-2-02 Tiger and sheep}一只守卫和两只闯入者在三节点通道上轮流移动。把题抽象成有限状态图，说明怎样判定守卫是否存在必胜策略，而不是逐句描述追逐过程。
  \item \MFQInterviewMeta{笔试}{12 分钟}{中}{约束搜索与下界}{GB-2-03 River crossing}一名研究员要把模型文件、密钥和原始数据送过河。船每次只能载研究员和一件物品，且必须由研究员驾驶；研究员不在某岸时，密钥不得与原始数据同岸，模型文件也不得与密钥同岸。给出最短运输序列，并用状态图的层数证明不能再少一次渡河。
  \item \MFQInterviewMeta{笔试}{10 分钟}{中}{同余与日历建模}{GB-2-04 Birthday problem}某人的生日在连续四年中分别落在周二、周三、周四、周六。只用闰年规则判断哪一年跨过了闰日，并列出仍无法确定的信息。
  \item \MFQInterviewMeta{口述}{10 分钟}{中}{不完全信息与策略}{GB-2-05 Card game}两人各抽一张 $1$ 到 $9$ 的牌，只看自己的牌，可选择加注或弃牌。说明为什么“按牌面单调加注”是合理候选策略，并指出还需哪些收益参数才能求阈值。
  \item \MFQInterviewMeta{笔试}{8 分钟}{中}{构造与误差边界}{GB-2-06 Burning ropes}两根绳子的总燃烧时间各为一小时，但燃速不均匀。设计测量四十五分钟的方法，并说明方案依赖的唯一物理假设。
  \item \MFQInterviewMeta{笔试}{15 分钟}{难}{信息论下界与决策树}{GB-2-07 Defective ball}十二个外观相同的零件中恰有一个重量异常，且未知偏轻偏重。先用叶节点计数给出称量次数下界，再说明构造决策树时必须保留哪些状态。
  \item \MFQInterviewMeta{笔试}{10 分钟}{中}{估值与计数}{GB-2-08 Trailing zeros}不用计算阶乘本身，求 $125!$ 末尾零的个数，并推广到任意进制 $b$；指出只数 $5$ 的因子何时会失效。
  \item \MFQInterviewMeta{研究判断}{12 分钟}{难}{分组淘汰与比较下界}{GB-2-09 Horse race}二十五个策略只能五个一组回测，每次只得到组内相对名次。没有计时器时找出全局前三至少需要几轮？给出候选淘汰表。
  \item \MFQInterviewMeta{数值编程}{12 分钟}{难}{递推与闭式验证}{GB-2-10 Infinite sequence}序列满足 $a_{n+2}=3a_{n+1}-2a_n$，只给定 $a_0,a_1$。推导闭式、验证极限条件，并说明数值递推何时会放大舍入误差。
  \item \MFQInterviewMeta{口述}{10 分钟}{中}{几何装箱与反例}{GB-2-11 Box packing}判断边长为 $1$ 的立方体能否放入尺寸为 $0.9\times1.2\times1.4$ 的长方体。该长方体体积大于 $1$、空间对角线也大于 $\sqrt3$；证明它仍无法容纳单位立方体，并解释为什么体积和对角线都不是充分条件。
  \item \MFQInterviewMeta{口述}{10 分钟}{中}{位值编码与可辨识性}{GB-2-12 Calendar cubes}用两个立方体的六个面标数字，使每天都能并排显示 $01$ 到 $31$。先写出数字出现次数的必要条件，再给出一组可行标记，并说明为什么 $6$ 与 $9$ 可共用一个旋转面。
  \item \MFQInterviewMeta{口述}{8 分钟}{中}{最坏情形决策}{GB-2-13 Door to offer}三扇门后分别是录用、加面和拒绝。证明一次是/否询问不能保证识别三种结果；若允许第二个根据首答自适应选择的是/否问题，给出最少询问协议和完整决策树。
  \item \MFQInterviewMeta{口述}{10 分钟}{中}{通信协议与公共知识}{GB-2-14 Message delivery}两名研究员只能通过一名可能丢失消息但不会篡改消息的信使通信。解释为什么“我知道你收到”会产生无限确认链，并给出工程上可终止的超时、序号与重试协议及其保证边界。
  \item \MFQInterviewMeta{笔试}{10 分钟}{中}{递推与奇偶性}{GB-2-15 Last ball}袋中有红蓝球，每次取出两球：若同色则放回一个蓝球，若异色则放回一个红球。设初始蓝球数为 $B$ 且总球数为 $N$，用不变量判断最后一球的颜色；不要模拟抽取顺序。
  \item \MFQInterviewMeta{笔试}{10 分钟}{中}{奇偶不变量}{GB-2-16 Light switches}一排 $n$ 个开关初始全灭，每次只能翻转相邻两个。刻画所有可达状态，并给出到达某个目标状态的线性时间判定法。
  \item \MFQInterviewMeta{口述}{8 分钟}{中}{期望值与条件缺失}{GB-2-17 Quant salary}公司给出“首年奖金期望为固定工资的 $40\%$”。说明仅凭期望为什么不能比较两个 offer 的效用，并列出至少四项还需询问的分布、约束或相关性信息。
  \item \MFQInterviewMeta{口述}{8 分钟}{中}{对称变换}{GB-2-18 Coin piles}一堆硬币中已知恰有 $k$ 枚正面，但不能观察单枚朝向。蒙眼把它分成两堆，使两堆正面数相同；证明翻转操作为何成立。
  \item \MFQInterviewMeta{口述}{8 分钟}{易}{条件排除}{GB-2-19 Mislabeled bags}三个袋子标签全部贴错，分别装两种资产及其混合。最少抽取几次可恢复全部标签？写出第一次抽样后为何没有分支歧义。
  \item \MFQInterviewMeta{笔试}{12 分钟}{难}{公共知识与递推}{GB-2-20 Wise men}三人各戴一顶黑或白帽，只能看见别人。主持人宣布“至少一顶黑帽”。若前两人依次说不知道，第三人如何推断？明确公共信息如何逐轮更新。
  \item \MFQInterviewMeta{口述}{8 分钟}{中}{圆周切割与最坏间隔}{GB-2-21 Clock pieces}从钟面上选两条经过圆心的直线切成四块。能否让四块上的数字和相等？先用总和给出目标，再用对径数字配对证明或否定可行性。
  \item \MFQInterviewMeta{数值编程}{10 分钟}{中}{等差和与缺失量}{GB-2-22 Missing integers}从 $1,\ldots,n$ 中删去两个不同整数，只给剩余数的和与平方和。推导恢复两数的公式，并写出输入不一致的判定条件。
  \item \MFQInterviewMeta{口述}{10 分钟}{难}{编码与称量}{GB-2-23 Counterfeit coins I}若一批硬币中假币都偏轻且数量未知，只有一次精密秤读数，怎样设计不同取样权重来编码假币集合？讨论量程和误差限制。
  \item \MFQInterviewMeta{笔试}{6 分钟}{易}{鸽巢原理}{GB-2-24 Matching socks}黑、灰、蓝三色袜子混放，闭眼最少取几只才能保证一双同色？若要求保证两双且两双颜色不同，答案如何变化？
  \item \MFQInterviewMeta{笔试}{8 分钟}{中}{图的度数和}{GB-2-25 Handshakes}一场活动中每两人至多握手一次。证明握手次数等于度数和的一半，并推出奇数度人数必为偶数。
  \item \MFQInterviewMeta{口述}{8 分钟}{中}{鸽巢原理与补图}{GB-2-26 Have we met before?}六位候选人中证明必有三人两两认识，或三人两两不认识。先固定一人，再对其五条关系分类。
  \item \MFQInterviewMeta{笔试}{8 分钟}{中}{轨迹对称}{GB-2-27 Ants on a square}单位正方形的四条边是轨道，四个顶点都是出口。若干蚂蚁从边上给定位置出发，以相同速度 $v$ 按给定顺/逆时针方向移动，相撞即掉头。用每只蚂蚁沿初始方向到下一个顶点的距离表示全部离开的时间，并证明“相撞掉头”等价于“穿透而过”的标签交换。
  \item \MFQInterviewMeta{笔试}{15 分钟}{难}{三进制编码}{GB-2-28 Counterfeit coins II}一次天平称量有三种结果。给定三次称量，设计可识别的异常状态编码，并说明为何每枚物品的编码不能与另一枚互为相反数。
  \item \MFQInterviewMeta{研究判断}{15 分钟}{难}{模运算与协同策略}{GB-2-29 Prisoner problem}多人依次观察他人的二元状态 $0/1$，每人回答时只能公开一个比特。设计奇偶校验协议，让除至多一人外都能恢复自己的状态；说明牺牲者的公开比特保存了什么信息。
  \item \MFQInterviewMeta{笔试}{8 分钟}{中}{十进制与模 $9$}{GB-2-30 Division by 9}证明一个十进制整数与其各位数字和模 $9$ 同余；设计可检测单个数字抄错的校验，并说明哪些错误模式仍无法被发现。
  \item \MFQInterviewMeta{笔试}{12 分钟}{难}{模不变量与可达性}{GB-2-31 Chameleon colors}初始有红、绿、蓝变色龙各 $8,4,3$ 只。每次两只不同颜色的变色龙相遇，二者都变成第三种颜色。用两两颜色计数之差模 $3$ 的不变性，判定是否可能最终全部同色。
  \item \MFQInterviewMeta{笔试}{10 分钟}{中}{构造与归纳}{GB-2-32 Coin split problem}有 $2n$ 枚硬币，其中 $n$ 枚正面朝上但位置未知。蒙眼把它们分成两堆，使两堆正面数相同；给出构造并证明它与初始排列无关。
  \item \MFQInterviewMeta{口述}{8 分钟}{易}{操作下界}{GB-2-33 Chocolate bar problem}一块 $m\times n$ 的巧克力每次只能拿一块现有碎片掰成两块。分成 $mn$ 个小格最少需要几次？证明下界并给出达到下界的任意策略。
  \item \MFQInterviewMeta{笔试}{12 分钟}{中}{相对速度与周期}{GB-2-34 Race track}两名跑者在长度 $L$ 的环形跑道同点同向出发，速度分别为 $u>v>0$。推导第 $k$ 次追及的时刻与位置；再说明反向出发时公式如何改变。
  \item \MFQInterviewMeta{口述}{10 分钟}{中}{反证法与构造}{GB-2-35 Irrational number}证明存在两个无理数 $a,b$ 使 $a^b$ 为有理数。答案不得预先假定 $\sqrt2^{\sqrt2}$ 的有理性，而应按其是否有理分情况构造。
  \item \MFQInterviewMeta{口述}{12 分钟}{难}{编码、信息与最坏保证}{GB-2-36 Rainbow hats}一列研究员各戴有限种颜色之一的帽子，只能看见前方并依次回答自己的颜色。设计使最坏情况下错误人数有界的编码协议，并明确答案依赖颜色数、回答顺序及能否听见先前回答。
\end{enumerate}

\section{作答协议}

每道题的完整答案都应包含：一句建模、一个下界或不变量、构造步骤、边界假设和至少一个常见误区。高频题还要准备面试官追问：参数变化后哪些结论仍成立，哪些必须重算。

本章不设独立综合项目；它的学习目标是把“状态--不变量--下界--构造--限制”这套方法带入六条研究路线，并为各路线的 Capstone 提供可迁移的推理习惯。谜题表现不能替代方向模型、样本外检验与交易摩擦证据。

\part{方向模型与研究项目}
\chapter{多因子模型：预测回归与风险价格}
\label{ch:lower-multifactor-model}

\begin{itemize}
  \item 直接先修：上册第 6、10、13、14、16、17 章；路线诊断见第 \ref{route:multifactor} 节。
  \item 学习产出：能从估计对象、数据轴和识别假设区分预测型 Fama--MacBeth 与经典两遍资产定价回归。
  \item 研究边界：本章的精确线性面板是计算 oracle，不是市场规律、因果证据或实盘 alpha。
\end{itemize}

\section{先画清三根轴}

令 $i$ 表示资产，$t$ 表示决策时点，$k$ 表示特征或因子。预测研究在 $t$ 时刻只能使用合法信息集 $\mathcal I_t$ 中的 $x_{i,t}$，目标是之后实现的 $r_{i,t+1}$：
\[
r_{i,t+1}=a_t+x_{i,t}^{\mathsf T}b_t+u_{i,t+1}.
\]
这里的 $b_t$ 是当期横截面预测系数。它不自动是结构性风险价格，也不等于因果处理效应。若把财报修订值、未来指数成分或 $t+1$ 的标准化参数回填到 $t$，线性代数仍可正确运行，但研究问题已经失效。

对单一中心化信号，含截距 OLS 斜率为
\[
\hat b_t=\frac{\sum_i(x_{i,t}-\bar x_t)(r_{i,t+1}-\bar r_{t+1})}
{\sum_i(x_{i,t}-\bar x_t)^2}.
\]
冻结小面板的两期斜率是 $0.020$ 与 $0.030$，所以时间均值是 $0.025$。把第二期信号乘以 2 会把该期斜率缩为 $0.015$，但不改变排序。这说明系数的单位依赖信号尺度，而 Rank IC 不依赖正的线性缩放。

\begin{diagnostic}
信号日期必须严格早于结果日期。整体调用 \texttt{shift(-1)} 只能改变数组位置，不能证明发布日期、修订历史、成分股资格和停牌状态已按时点冻结。
\end{diagnostic}

\section{预测型 Fama--MacBeth}

第一步在每个 $t$ 做一次横截面回归，保存 $\hat b_t$；第二步对完整系数序列求均值：
\[
\bar b=\frac1T\sum_{t=1}^T\hat b_t,
\qquad
\widehat{\operatorname{se}}(\bar b)
=\sqrt{\frac{1}{T(T-1)}\sum_{t=1}^T(\hat b_t-\bar b)^2}.
\]
若滚动收益或平滑信号令 $\hat b_t$ 自相关，应使用声明带宽的 HAC 长期方差，而不能把同一期资产数当成独立时间样本\cite{fama1973risk,newey1987hac}。预测型协议回答“给定时点特征是否预测下一期收益”；它不要求特征是可交易因子的 beta。

\section{经典两遍资产定价回归}

经典两遍法的第一步沿时间轴估计每个资产对可交易或已声明因子 $f_t$ 的暴露：
\[
r_{i,t}=\alpha_i+\beta_i^{\mathsf T}f_t+\varepsilon_{i,t}.
\]
第二步沿资产轴回归平均收益：
\[
\bar r_i=\lambda_0+\beta_i^{\mathsf T}\lambda+\eta_i.
\]
$\lambda$ 才是该模型中的因子风险价格。本章精确面板设 $\beta=(0.5,1,1.5)$、资产 alpha 为 $0.004$，因子收益均值为 $0.006$；两遍回归恢复相同 beta，并得到 $\hat\lambda=0.006$。这与前一节的预测斜率 $0.025$ 是不同估计对象，数值相等也不能互换含义。

第二遍使用第一遍估计的 $\hat\beta_i$，因此解释变量含估计误差。Shanken 型修正讨论这类生成回归量对风险价格推断的影响；常见标量诊断包含
\[
\sqrt{1+\lambda^{\mathsf T}\Sigma_f^{-1}\lambda}.
\]
它依赖具体模型、同方差/条件分布、因子可交易性和样本设定，不能把这个乘数机械套在所有预测回归标准误上\cite{shanken1992beta}。

\section{相关、秩相关与中性化}

Pearson IC 保留线性距离，Rank IC 只保留次序并需要声明并列秩规则。中性化则把信号投影到暴露矩阵 $Z_t$ 的正交补：
\[
x_t^\perp=x_t-Z_t(Z_t^{\mathsf T}Z_t)^{-1}Z_t^{\mathsf T}x_t,
\qquad Z_t^{\mathsf T}x_t^\perp=0.
\]
实现应使用 QR、SVD 或最小二乘而非显式求逆，并在秩亏或条件数过高时拒绝。中性化不是免费的“清洁”：若规模同时承载预测信息，残差化会连同风险暴露一起删除。

\section{多重检验的保证边界}

每尝试一个特征、窗口、预处理、持有期或资产池，就消耗一次研究选择权。Benjamini--Hochberg（BH）把排序后的 $p_{(j)}$ 与 $j\alpha/m$ 比较，在独立或适当正依赖条件下控制 FDR\cite{benjamini1995controlling}。任意相关、层级搜索或根据结果重定义检验族时，不能只报告“BH 已通过”；应保存完整尝试账本，并考虑更保守校正、重采样或独立最终测试。

冻结零信号搜索进行 20 次尝试，朴素 $p<0.05$ 留下 2 项，BH 拒绝 0 项。这不是证明所有信号无效，而是说明当前检验族的证据不足。删除 18 个失败尝试后把 $m$ 写成 2，会直接破坏审计边界\cite{white2000reality,bailey2017backtest}。

\begin{implementationnote}
运行 \path{notebooks/lower/multifactor_model.py}。图中分别显示预测斜率序列与经典风险价格；日期相等的故障注入必须以稳定类别拒绝。
\end{implementationnote}

\section{连接到路线 Capstone}

Capstone 的模型卡必须先声明属于“特征预测”还是“beta 定价”，列出时间轴、资产轴、因子轴、参数单位和识别限制。只有模型对象冻结后，下一章的正则化与库实现才有可比较的目标。

\section{综合练习}
\begin{enumerate}[leftmargin=*]
  \item \textbf{口述概念：}为什么预测型横截面系数、经典风险价格和因果效应不是同一对象？
  \item \textbf{口述概念：}何时 Pearson IC 与 Rank IC 会给出不同结论？
  \item \textbf{笔试推导：}从两期斜率 $0.020,0.030$ 推导均值和标准误。
  \item \textbf{笔试推导：}写出经典两遍法的两根回归轴，并解释生成回归量误差。
  \item \textbf{数值编程：}复现 beta 与风险价格，再把资产 alpha 改成不同值观察第二遍截距。
  \item \textbf{数值编程：}给零信号搜索加入相关候选，比较 BH 结论并记录依赖条件。
  \item \textbf{研究判断：}审计一份把预测系数称为风险价格的摘要，列出缺失假设。
  \item \textbf{研究判断：}说明为什么最终测试期不能参与信号尺度、窗口和检验族选择。
\end{enumerate}


\section{正则化、数值实现与估计不确定性}
\label{ch:lower-multifactor-estimation}

\begin{itemize}
  % Generated from curriculum/manifest.json; do not edit by hand.
\item 直接先修：下册《多因子模型：预测回归与风险价格》。

\end{itemize}

两个几乎相同的因子可能给出稳定的预测，却给出剧烈变化的系数。正则化针对的正是这种可识别方向不足的问题；先看目标怎样变化，再讨论参数选择。


\MFQLead{从正规方程到数值求解}


以下回归使用决策时点的信息集 $\mathcal I_t$。标准化等变换由当时已知的数据估计，因此模型的预测对应一个明确的时间起点。

线性模型 $y=X\beta+\varepsilon$ 的 OLS 目标是
\[
\min_\beta \frac12\lVert y-X\beta\rVert_2^2.
\]
满列秩时正规方程给出 $(X^{\mathsf T}X)\hat\beta=X^{\mathsf T}y$，但实现不应显式计算逆矩阵。QR 保留最小二乘几何，SVD 还能揭示近零奇异值；条件数检查回答“该样本能否稳定区分这些暴露”，不能被“求解器返回成功”替代。

面板研究还要声明误差依赖：同一资产跨期相关、同一期资产共同冲击、行业聚类与重叠持有期会改变标准误。固定效应吸收常数异质性，但不会自动解决未来信息、动态偏误或错误聚类层级。

\MFQLead{Ridge：用偏差换稳定}

Ridge 目标为
\[
\min_\beta \frac12\lVert y-X\beta\rVert_2^2+\frac\lambda2\lVert\beta\rVert_2^2,
\qquad
\hat\beta_{\mathrm{ridge}}=(X^{\mathsf T}X+\lambda I)^{-1}X^{\mathsf T}y.
\]
$\lambda>0$ 抬高小特征值，缓和共线性导致的方差放大，但引入偏差\cite{hoerl1970ridge}。惩罚参数必须只在训练/选择窗口决定；用最终测试期选 $\lambda$ 会把样本外变成调参样本。

设 $X=UDV^T$，奇异值为 $d_j$。把 $\beta=V\theta$ 代入 ridge 正规方程，逐坐标得到
\[
 \hat\theta_j=\frac{d_j}{d_j^2+\lambda}u_j^Ty.
\]
OLS 在该方向的系数为 $u_j^Ty/d_j$，所以 ridge 把它乘以 $d_j^2/(d_j^2+\lambda)$。小奇异值对应的方向收缩最多；它不是把所有系数乘同一常数。$d_1=10,d_2=0.1,\lambda=1$ 时两个收缩因子分别约为 $0.9901,0.0099$。

\MFQLead{Lasso：稀疏不是稳定选择}

Lasso 目标为
\[
\min_\beta \frac12\lVert y-X\beta\rVert_2^2+\lambda\lVert\beta\rVert_1.
\]
坐标下降更新第 $j$ 维时，对去除本维后的残差计算得分 $z_j=x_j^{\mathsf T}r_{-j}$，再使用软阈值
\[
\beta_j\leftarrow\frac{\operatorname{sign}(z_j)(|z_j|-\lambda)_+}{x_j^{\mathsf T}x_j}.
\]
该软阈值来自绝对值在零点的次梯度区间\cite{tibshirani1996regression}。零范数列不影响拟合，惩罚参数为正时可直接令其系数为零；高度相关特征中“选中哪一个”可能随小扰动改变，所以稀疏不等于经济解释稳定。

\MFQLead{同一组数据比较三种估计}

一列设计 $x=(1,2)$、$y=(1,1)$，有 $x^Tx=5,x^Ty=3$。OLS 为 $3/5$；ridge 取 $\lambda=1$ 得 $3/6=1/2$；lasso 的次梯度方程给 $\operatorname{soft}(3,1)/5=2/5$。若损失改为平均平方和 $\|y-X\beta\|^2/(2n)$，保持同一解需要把惩罚参数也除以 $n$。notebook 会显示这一列算例和两个共线特征的惩罚路径，允许改变噪声和惩罚强度。

\MFQTransition{滚动、扩展与加权估计}

扩展窗口使用全部历史，方差较小但对制度变化反应慢；固定滚动窗口遗忘旧数据，方差更大。指数权重最小化
\[
\sum_{s\le t}\lambda^{t-s}(y_s-x_s^{\mathsf T}\beta)^2,
\qquad 0<\lambda<1,
\]
在两者之间连续折中。半衰期 $h_{1/2}=\log(1/2)/\log\lambda$ 应用时间单位解释，而不是只报 $\lambda$。

横截面回归常按市值或误差方差加权。权重改变估计目标：市值权重更接近大盘可投资组合，等权更强调小股票，逆方差权重依赖一个额外噪声模型。权重若由未来成交量或事后波动估计，同样泄漏。

\MFQTransition{共线性、测量误差与稳定选择}

高度相关因子使 $X^{\mathsf T}X$ 条件数升高，单个系数对样本扰动敏感。Ridge 在相关方向上收缩系数，Lasso 则可能只保留其中一个变量。对样本作重抽样或改变估计窗口，可以观察选择频率与系数分布，了解这种选择有多稳定。

若特征带测量误差 $x^*=x+u$，简单回归系数通常向零衰减。财务数据库的回填、供应商映射错误和估计型特征都可能造成测量误差。正则化并没有建立测量误差模型，因此不能由此消除衰减偏差。

\MFQTransition{面板标准误与聚类层级}

误差可能在同一日期跨资产相关，也可能在同一资产跨期相关。双向聚类协方差可概念性写为
\[
\widehat V_{\rm two-way}=\widehat V_{\rm date}
+\widehat V_{\rm asset}-\widehat V_{\rm intersection}.
\]
聚类数太少时渐近近似不可靠，应报告簇数并考虑 wild bootstrap。标准误方法不能事后根据显著性选择。

\section{组合构建与交易摩擦}
\label{ch:lower-multifactor-research}

\begin{itemize}
  % Generated from curriculum/manifest.json; do not edit by hand.
\item 直接先修：下册《多因子估计：正则化、数值实现与外部数据》。

\end{itemize}

因子分数给出排序，交易需要确定数量。下面从同一组股票的排序出发，算出目标权重、实际换手和成本，让每一步都能追溯到前一步。


\MFQTransition{从预测到仓位与收益}

信号到净收益的联系可以写为
\[
x_{i,t}\longrightarrow \tilde x_{i,t}\longrightarrow q_{i,t}
\longrightarrow w_{i,t}\longrightarrow R^{\mathrm{gross}}_{t+1}
\longrightarrow R^{\mathrm{net}}_{t+1}.
\]
$\tilde x$ 是时点预处理后的信号，$q$ 是分组，$w$ 是实际仓位。组合收益由这些仓位与之后实现的资产收益共同决定；只计算信号的相关性，还不能得到组合收益。

% Generated by tools/build_figures.py; edit figures/figure-specs.json instead.
\MFQFigure{tex/figures/generated/lower-mf-pipeline.pdf}{本例依次将信号转为中性化暴露、目标权重与持仓；持仓产生收益，调仓带来成本。}{lower-mf-pipeline}


\MFQTransition{分组、权重与重叠持有期}

先在每个调仓日用事先确定资产池排序，再把最低组做空、最高组做多。等权使每一侧绝对权重和为 1；市值权重则在组内按正市值归一化。两者回答不同问题，不能因为某一结果更好而事后选择。

若持有 $H$ 期，每个时点的有效仓位是仍存续 vintage 的平均：
\[
w_t=\frac1{|A_t|}\sum_{s\in A_t}w_t^{(s)},
\qquad A_t=\{s:t-H+1\le s\le t\}.
\]
这会平滑仓位并引入收益重叠；推断时需处理序列相关。资产退市、停牌或不可卖空不能靠删除行解决，必须进入可交易集合和退出规则。

\MFQTransition{换手、成本与容量}

双边换手定义为 $\tau_t=\sum_i|w_{i,t}-w_{i,t-1}|$。净收益可写为
\[
R^{\mathrm{net}}_{t+1}=w_t^{\mathsf T}r_{t+1}
-c\tau_t-\operatorname{impact}(w_t,\mathrm{ADV}_t,\sigma_t).
\]
两期等权例的毛收益分别为 $0.035$ 与 $0.030$，期均毛收益 $0.0325$；平均双边换手为 2，单位换手成本 $0.001$，容量冲击 $0.0005$，所以净收益为 $0.0300$。市值加权与两期持有给出不同仓位和净收益，这是一项预先声明的敏感性分析，不是择优汇报。

% Generated by tools/build_figures.py; edit figures/figure-specs.json instead.
\MFQFigure{tex/figures/generated/lower-mf-cost-chain.pdf}{预测优势只有经过权重映射、换手和成本扣除后才成为净收益；每一层都可能改变排序。}{lower-mf-cost-chain}


真实报告应拆出佣金、价差、冲击、借券、税费、参与率、成交完成率与资金规模。容量不是一个永久常数；应报告规模到净收益、参与率和成交失败的曲线。

\MFQTransition{两种容易混淆的研究结论}

\textbf{情形 A：预测与收益断开。} 程序打印 AR、回归或 ML 预测，但 \texttt{trade\_return} 直接来自 合成样本。审计步骤是追踪每一笔 $w_t$ 是否由事先确定预测生成，并用 $w_t r_{t+1}$ 独立重算。修复不是给报告补一句“模型驱动”，而是删除外生收益字段。

\textbf{情形 B：选择后重写检验族。} 研究者测试 500 个信号，只把 20 个较好结果交给 BH，并把 $m$ 写成 20。检验族应包括实际参与选择的 500 个信号。只保留较好的 20 项会改变错误率计算；这与某项结果是否来自样本外是两个需要分别判断的问题。

\MFQTransition{分组组合的逐步算例}

假设五只股票得分排序为 $A>B>C>D>E$，未来收益为 $(3\%,1\%,0\%,-1\%,-2\%)$。最简单多空组合做多 $A$、做空 $E$，毛收益为 $5\%$；若每侧取两只等权，则收益为
\[
\frac{3\%+1\%}{2}-\frac{-1\%-2\%}{2}=3.5\%.
\]
单只头尾更容易被个股异常主导，分组更稳但信号稀释。市值加权、风险加权和分数加权回答不同投资问题，不能在测试期择优。

若下一期排名变成 $E>A>C>B>D$，目标仓位变化决定换手。换手应从旧实际仓位到新目标仓位计算，部分成交时还要保留未完成订单。只按入选名单变化计数会忽略权重大小。

\MFQTransition{例：从两次排名计算换手}

前一期排名为 $A>B>C>D>E$，每侧取两只等权，权重按 $A,B,C,D,E$ 排列为
\[
 (0.5,0.5,0,-0.5,-0.5).
\]
下一期排名变为 $E>A>C>B>D$，新权重为 $(0.5,-0.5,0,-0.5,0.5)$。暂忽略持有期间权重漂移，调仓差为 $(0,-1,0,0,1)$，双边换手为 2。

若单位换手成本为 10 个基点，成本为 $0.001\times2=0.002$。下一期收益为 $(0.03,0.01,0,-0.01,-0.02)$，新仓位毛收益为 $0.015-0.005+0.005-0.01=0.005$，净收益为 $0.003$。高分股票 $E$ 实现负收益，说明排序规则可以明确而单期预测仍失败。

真实调仓应先把上期持仓按已实现价格更新到调仓前权重，再计算订单；若只成交部分订单，下一期必须从实际仓位继续。notebook 将把两个排名转成权重数组并显示每只股票的贡献，读者可改变成本率，直接找出本期盈亏平衡成本 $0.005/2=0.0025$。

\MFQTransition{归因：模型收益究竟来自哪里}

组合收益可分为市场、行业、风格和特异部分。一个“质量 Alpha”若主要来自低波动或小盘暴露，应先说明它是复合暴露，而非把中性化后的剩余全部称为选股能力。归因至少按日期、行业、市值、流动性和多空两侧展开。

回撤归因比平均收益更能暴露失效：是信号翻转、拥挤、成本上升、股票池变化还是数据缺口？对每种失败预先定义诊断统计量，避免亏损后任意讲故事。

\MFQTransition{稳健性不是堆叠更多回测}

合理稳健性包括邻近参数、不同但事前合理的股票池、子时期、成本压力和数据供应商交叉核对。尝试越多，越要把全部结果计入选择权。热力图不能只截取正收益区域。

真正的最终测试是研究设计事先确定后一次性打开的未来区间。一旦利用这一结果继续修改模型，该区间就参与了选择；下一次确认需要新的、未参与修改的数据。

\begin{diagnostic}
高 IC、单调分组和显著 t 值可以同时存在，但若信号换手过高、容量过低或依赖不可借券股票，仍不是可实施策略。统计证据和实施证据必须并列。
\end{diagnostic}

\section{本路线的综合项目}

选择一个因子问题，将前面的回归、估计与组合构建联系起来。在未参与选择的区间比较分组收益，并考察成本与交易规模变化的影响。另选一种信号失效或成交受限的情形，解释净收益为何改变。

\section{综合练习}
\begin{enumerate}[leftmargin=*]
  \item \textbf{口述概念：}为什么预测型横截面系数、经典风险价格和因果效应不是同一对象？
  \item \textbf{口述概念：}何时 Pearson IC 与 Rank IC 会给出不同结论？
  \item \textbf{笔试推导：}从两期斜率 $0.020,0.030$ 推导均值和标准误。
  \item \textbf{笔试推导：}写出经典两遍法的两根回归轴，并解释生成回归量误差。
  \item \textbf{数值编程：}复现 beta 与风险价格，再把资产 alpha 改成不同值观察第二遍截距。
  \item \textbf{数值编程：}给零信号搜索加入相关候选，比较 BH 结论并记录依赖条件。
  \item \textbf{研究判断：}审计一份把预测系数称为风险价格的摘要，列出缺失假设。
  \item \textbf{研究判断：}说明为什么最终测试期不能参与信号尺度、窗口和检验族选择。
\end{enumerate}

\MFQLead{估计、正则化与数值实现题组}
\begin{enumerate}[leftmargin=*]
  \item \textbf{口述概念：}数值稳定、统计识别和经济解释分别由什么证据支持？
  \item \textbf{口述概念：}为什么“lasso 选中”不等于变量解释稳定？
  \item \textbf{笔试推导：}从 ridge 目标函数推导线性方程。
  \item \textbf{笔试推导：}推导 lasso 单坐标的软阈值更新。
  \item \textbf{数值编程：}画出 ridge/lasso 惩罚路径并比较透明实现与成熟库。
  \item \textbf{数值编程：}注入全零列与近共线列，分别记录失败或条件数响应。
  \item \textbf{研究判断：}审计一份用最终测试期选择 $\lambda$ 的研究包。
  \item \textbf{研究判断：}解释 WDI 截面为什么不能支持股票 alpha 或因果结论。
\end{enumerate}

\MFQTransition{组合实施、应用分析与开放项目 题组}
\begin{enumerate}[leftmargin=*]
  \item \textbf{口述概念：}信号、排序、目标权重、成交权重和净收益分别是什么对象？
  \item \textbf{口述概念：}等权、市值权重与重叠持有期改变了什么研究问题？
  \item \textbf{笔试推导：}手算两期等权组合的毛收益、换手、成本和净收益。
  \item \textbf{笔试推导：}写出 $H$ 期重叠持有的有效权重及其推断后果。
  \item \textbf{数值编程：}实现五分组、等权/市值权重与 1/3/6 期持有敏感性表。
  \item \textbf{数值编程：}从各期实际持仓和下一期资产收益计算组合毛收益，再根据成交增量扣除成本。
  \item \textbf{研究判断：}审计“只保留 20 个候选再做 BH”的研究报告。
  \item \textbf{研究判断：}审计未记录停牌、融券与成交完成率的 A 股组合。
  \item \textbf{综合项目：}选择一个因子，完成从回归到分组组合的分析。比较两种持有期或权重规则，并解释成本与样本选择怎样影响结果。
\end{enumerate}

\subsection*{基础衔接：独立计算}
同一组合的总权重与绝对权重和分别是多少？能否把 0.5\% 毛收益称作未加杠杆收益？


\chapter{时间序列、协整与状态估计：从单位根到 OU}
\label{ch:lower-stat-arb-model}

\begin{itemize}
  % Generated from curriculum/manifest.json; do not edit by hand.
\item 直接先修：上册第 9 章《概率收敛、极限定理与集中不等式》、上册第 11 章《随机过程、Markov、鞅、Poisson 与 Brownian》、上册第 12 章《时间序列与状态空间基础》、上册第 14 章《浮点数、数值线代与病态问题》、上册第 15 章《Monte Carlo、bootstrap 与方差缩减》、上册第 16 章《研究有效性、样本外与交易摩擦》、上册第 17 章《可复现研究审计》。

  \item 学习产出：把单位根、协整、OU、状态空间、regime filtering 与变点检测串成可估计模型，并将预测接入仓位、成交、成本和失效监控。
  \item 研究边界：合成路径是数学 oracle；公开快照只验证数据协议，任何统计套利结论都必须经过候选选择、多重检验和 walk-forward 台账。
\end{itemize}

\section{先区分平稳、单位根与共同趋势}

AR(1) 模型 $X_t=\phi X_{t-1}+\varepsilon_t$ 在 $|\phi|<1$ 且创新满足适当矩条件时具有稳定二阶结构。若 $\phi=1$，冲击永久进入水平，通常要研究 $\Delta X_t$。Dickey--Fuller 回归
\[
\Delta X_t=a+\gamma X_{t-1}+u_t
\]
检验 $\gamma=0$，其 t 统计量不服从普通 Student t 分布，必须使用与确定性项、滞后阶相匹配的临界值。ADF 不是“平稳证明器”：低功效、结构突变和错误滞后都可能改变结论。
所有滞后、变换与预测都必须可由决策时信息集 $\mathcal F_t$ 构造；全样本选规格不属于该信息集。

\begin{diagnostic}
先画水平、差分和滚动尺度，再确定确定性项与最大滞后。看到一个小 p 值以后才选择含趋势或不含趋势的规格，会把规格搜索藏进检验。
\end{diagnostic}

\MFQTransition{单位根检验的原假设与检验力}

ADF 回归在差分方程中检验水平项系数，原假设是单位根；KPSS 的原假设则是平稳。二者结论可以组合成“平稳”“单位根”或“证据不足”。短样本、接近单位根和结构突变会显著降低检验力，不能把 $p>0.05$ 翻译成“证明存在单位根”。

确定性项必须事前决定。对带趋势序列漏掉趋势会错误拒绝或接受；滞后阶数过低残差相关，过高又损失自由度。价格通常先看对数水平与差分，利差和估值比率则要依据经济含义选择变换。

\section{Engle--Granger 与误差修正}

若 $x_t,y_t$ 都是 $I(1)$，但存在 $\beta$ 使 $z_t=y_t-a-\beta x_t$ 为 $I(0)$，则二者协整。Engle--Granger 两步法先估长期关系，再对残差做专用单位根检验\cite{engle1987cointegration}。用于核对运算的确定性算例令
\[
y_t=1.5+2x_t+0.55^t,
\]
对此数据作 OLS 得斜率约 $1.9912$，残差 DF 统计量约 $-8.3135$。斜率没有精确等于 2，是有限样本中确定性残差与趋势相关造成的差异。这里 $0.55^t$ 没有随机创新，不能用这一条序列说明平稳模型的检验力。

配套教学实验另使用 $x_t=x_{t-1}+u_t$、$z_t=\phi z_{t-1}+\eta_t$、$y_t=a+\beta x_t+z_t$，其中 $u_t\sim N(0,1)$、$\eta_t\sim N(0,\sigma_\eta^2)$ 相互独立且跨期独立，$|\phi|<1$。令 $z_0\sim N(0,\sigma_\eta^2/(1-\phi^2))$ 并独立于未来创新，则价差从一开始就平稳，而两个水平序列共享随机趋势。改变样本量并重复估计，才能看到有限样本误差。

% Generated by tools/build_figures.py; edit figures/figure-specs.json instead.
\MFQFigure{tex/figures/generated/lower-sa-cointegration.pdf}{两条非平稳价格路径可以共享共同趋势，而特定线性组合保持平稳并围绕均值波动。}{lower-sa-cointegration}





长期偏离进入误差修正模型：
\[
\Delta y_t=c+\alpha z_{t-1}+\delta\Delta x_t+u_t.
\]
合成例得到 $\hat\alpha\approx-0.4606$；负号表示 $y_t$ 对正偏离向下调整。只有在协整向量、方向和时点协议固定后，$\alpha$ 才能解释为调整速度。

\section{选读：Johansen 的多变量秩问题}

VAR 的 VECM 形式为
\[
\Delta w_t=\Pi w_{t-1}+\sum_{j=1}^{p-1}\Gamma_j\Delta w_{t-j}+u_t,
\qquad \Pi=\alpha\beta^{\mathsf T}.
\]
$\operatorname{rank}(\Pi)=r$ 是协整秩。Johansen 方法用广义特征值构造 trace 与最大特征根统计量，并按确定性项与滞后阶读取非标准临界值\cite{johansen1991estimation}。本项目的透明奇异值诊断只帮助理解“共同趋势方向”；正式 rank=1 结论由 \texttt{statsmodels} Johansen 临界值表独立计算。两者假设不同，不能把透明诊断冒充正式检验。

\section{从离散 AR 映射到 OU}

连续时间 OU 过程
\[
dZ_t=\kappa(\mu-Z_t)dt+\sigma dW_t
\]
对 $e^{\kappa t}(Z_t-\mu)$ 使用乘积法则，再在 $[t,t+\Delta]$ 积分，得到
\[
Z_{t+\Delta}=\mu+e^{-\kappa\Delta}(Z_t-\mu)
+\sigma\int_t^{t+\Delta}e^{-\kappa(t+\Delta-s)}dW_s.
\]
最后一项是独立于时点 $t$ 信息的零均值正态创新，方差为 $\sigma^2(1-e^{-2\kappa\Delta})/(2\kappa)$，所以 $\phi=e^{-\kappa\Delta}$。条件均值偏离按指数衰减，令其减半即得到半衰期。当 $0<\hat\phi<1$ 时，
\[
\hat\kappa=-\frac{\log\hat\phi}{\Delta},
\qquad t_{1/2}=\frac{\log 2}{\hat\kappa}.
\]
合成残差给出 $\hat\phi\approx0.5442$、半衰期约 $1.1393$ 个周期。$1/\hat\kappa$ 只是特征回归时间，不是首次通过时间。若从 $z_0>0$ 首次命中均值边界 0，扩散率为 $\sigma$，边值问题的解为
\[
\mathbb E_{z_0}\tau_0
=\frac{\sqrt\pi}{\kappa}
\int_0^{z_0\sqrt\kappa/\sigma}e^{u^2}\operatorname{erfc}(u)\,du.
\]
实现从离散创新方差映射 $\sigma$，并数值积分该式；起点、边界或扩散率改变时必须重算，不能用半衰期替代。

% Generated by tools/build_figures.py; edit figures/figure-specs.json instead.
\MFQFigure{tex/figures/generated/lower-sa-ou.pdf}{OU 漂移强度决定条件均值回归速度；半衰期描述偏离缩小一半所需的时间尺度。}{lower-sa-ou}





\begin{note}[Python 边界]
\texttt{adfuller}、\texttt{coint} 与 \texttt{coint\_johansen} 的默认确定性项、自动滞后和临界值并不相同。调用函数前先把统计问题写成完整规格，再在报告中保存库版本与参数。
\end{note}

\paragraph{实验阅读。}先比较确定性衰减与随机平稳价差，再计算 OLS 残差、协整检验和误差修正系数。公开宏观快照用于理解观测频率与真实数据背景，不作为可交易配对。

\section{ARMA、VAR 与 GARCH 的职责}

ARMA 描述条件均值的线性动态，VAR 处理多变量相互滞后，GARCH 描述条件方差：
\[
r_t=\mu_t+\varepsilon_t,\qquad
\varepsilon_t=\sigma_tz_t,\qquad
\sigma_t^2=\omega+\alpha\varepsilon_{t-1}^2+\beta\sigma_{t-1}^2.
\]
$\alpha+\beta$ 接近 1 表示波动冲击衰减慢；它不是价格均值回归速度。把 GARCH 波动预测直接当方向信号会混淆条件均值与条件方差。

残差诊断至少包括自相关、平方残差自相关、尾部分布和参数稳定。模型阶数可用训练期信息准则筛选，再用滚动样本外验证；不能在全样本 ACF 图上挑阶数后仍称测试独立。

\MFQTransition{协整向量的经济解释与不稳定性}

两个 $I(1)$ 序列若存在 $\beta$ 使 $z_t=y_t-\beta x_t$ 平稳，则共享随机趋势。$\beta$ 的尺度要固定，例如把 $y$ 系数规范为 1。纯统计搜索会从大量资产中找到偶然平稳组合，因此候选形成机制和多重检验必须记录。

滚动 $\hat\beta_t$ 变化意味着对冲比率变化，也可能说明关系失效。若每期重新估计并立刻用新 beta 重写过去价差，历史信号会被回看修订；研究账本要保存当时版本的 beta 和价差。

这组模型为后续路线 Capstone 提供状态、半衰期与失效边界；它们描述的是可检验的动态假设，不是收益承诺。

\chapter{估计与状态推断：滤波、切换与变点}
\label{ch:lower-stat-arb-estimation}

\begin{itemize}
  % Generated from curriculum/manifest.json; do not edit by hand.
\item 直接先修：下册《动态与长期关系：从单位根到 OU》。

  \item 学习产出：能区分预测、滤波与平滑信息集，并实现 Kalman、regime filtering 与变点评价。
  \item 研究边界：状态编号没有天然经济含义；状态数、分布与初值都属于模型选择。
\end{itemize}

\section{状态空间：三种条件分布}

线性高斯模型写成
\[
s_t=Fs_{t-1}+\eta_t,\qquad y_t=Hs_t+\epsilon_t,
\]
其中 $\eta_t\sim N(0,Q)$、$\epsilon_t\sim N(0,R)$。预测分布只条件于 $y_{1:t-1}$；滤波分布条件于 $y_{1:t}$；平滑分布条件于完整 $y_{1:T}$。Kalman 更新
\[
K_t=P_{t|t-1}H^{\mathsf T}
(HP_{t|t-1}H^{\mathsf T}+R)^{-1}
\]
该式表示用创新协方差缩放预测误差；实现中应使用线性求解而非显式求逆\cite{kalman1960new}。
定义在线信息集
\[
\mathcal F_t=\sigma(y_1,\ldots,y_t).
\]
交易决策最多使用条件于 $\mathcal F_t$ 的滤波量。

冻结四期观测 $0,0.2,2.8,3.1$ 后，在线滤波末值约 $2.3096$；第二期平滑值约 $1.0764$，明显使用了后两期信息。回测在第二期决策时读取该平滑值就是未来函数。

透明递推与 \texttt{statsmodels} 对照必须使用完全相同的 $F,H,Q,R$ 和初始分布；冻结实验的滤波与平滑路径最大差仅为浮点舍入量级。把 $Q$ 从 $0.1$ 提高到 $0.8$ 后，末期滤波状态变化约 $0.6614$，说明过程噪声不是无关紧要的实现参数。

\section{两状态过滤与可识别性}

隐状态 $S_t\in\{0,1\}$ 按转移矩阵 $P$ 演化，观测密度为 $p(y_t\mid S_t=j)$。Hamilton 过滤递推为
\[
\Pr(S_t=j\mid y_{1:t})\propto
p(y_t\mid S_t=j)\sum_i p_{ij}\Pr(S_{t-1}=i\mid y_{1:t-1}).
\]
取均值 $0,3$、方差 $0.2$、持续概率 $0.95$，最后两期大观测令状态 1 概率接近 1。透明实现验证给定参数的 Bayes 过滤；成熟库的 Markov switching 则估计参数并报告对数似然和归一化概率。二者假设不同，不能把数值差伪装成同一契约对拍。标签交换意味着“状态 1=牛市”不是识别约束\cite{hamilton1989new}。

\begin{diagnostic}
只报告最可能状态会隐藏参数不确定性。至少保存全概率路径、转移概率、初值、多起点结果与标签排序规则；若不同起点得到相近似然但不同解释，应报告不可识别。
\end{diagnostic}

\section{变点：误报、漏检与延迟}

变点算法必须冻结最短样本、阈值和损失。已知真值为索引 2 时，检测评价至少包含：提前警报数、首次有效警报、检测延迟和漏检。调低阈值通常增加误报，调高阈值通常增加漏检；只展示一条“漂亮”的检测线没有样本外含义。

regime switching 假设状态持续随机转换，change-point 假设参数在未知时点发生结构变化。二者可以描述相似图形，却回答不同问题。把事后变点分段当作事前可交易状态同样是泄漏。

\section{数值实现与失败恢复}

协方差矩阵应保持对称半正定；概率应归一化；对数似然应使用 log-sum-exp 防下溢。缺测规则必须写明“跳过更新”还是“插值后更新”。当优化器不收敛时，不应只增加迭代次数：先检查尺度、参数边界、初值、多峰似然与状态可分性。

\begin{implementationnote}
\begin{lstlisting}[language=bash]
uv run python notebooks/lower/stat_arb_estimation.py \
  evidence/stat-arb-estimation/oracle.json
\end{lstlisting}
实验逐项比较在线滤波与平滑，画出状态路径，并用真实时间验证器拒绝未来平滑状态。
\end{implementationnote}

本单元向 Route Capstone 交付仅使用当时信息的滤波状态、状态概率、变点指标和失效标志；成交与收益留到下一单元。

\section{分层练习}
\begin{enumerate}[leftmargin=*]
  \item \textbf{口述概念：}写出预测、滤波、平滑各自条件的信息集。
  \item \textbf{概念：}为什么 regime 标签会交换？
  \item \textbf{笔试推导：}推导标量 Kalman 增益与后验方差。
  \item \textbf{推导：}写出两状态过滤的一步预测与 Bayes 更新。
  \item \textbf{数值编程：}改变 $Q/R$，记录滤波路径的响应速度。
  \item \textbf{编程：}用多组初值拟合 Markov switching，比较似然与状态解释。
  \item \textbf{研究判断：}报告用全样本平滑状态回填历史交易。构造能执行的负例。
\end{enumerate}


\section{从预测到净收益的统计套利研究}
\label{ch:lower-stat-arb-research}

\begin{itemize}
  % Generated from curriculum/manifest.json; do not edit by hand.
\item 直接先修：下册《估计与状态推断：滤波、切换与变点》。

\end{itemize}

均值回归预测要先变成两腿订单，才能计算损益。以下区分预测变化、目标仓位和实际成交，并说明对冲比率更新为什么不等于持仓盈利。


\MFQTransition{先画时间轴，再拟合模型}

标签跨越 $h$ 期时，应删除结果尚未在拟合截止前完成、或与验证标签共享信息的训练样本。额外缓冲由标签完成、发布延迟与依赖结构决定，不是任何验证后都必须固定空置若干期。事先确定示例为训练索引 $0{:}5$、验证 $8{:}9$、交易 $12{:}13$，$h=2$、embargo=2。若训练扩到 6，其标签触及验证起点，验证器立即拒绝。随机 K 折、全样本 scaler、用交易期重新选阈值也由同一信息时间轴判断，而不是靠函数名称判断。

% Generated by tools/build_figures.py; edit figures/figure-specs.json instead.
\MFQFigure{tex/figures/generated/lower-sa-timeline.pdf}{滚动研究先用已有数据估计关系及选择阈值，再观察后续交易区间；重叠标签需要单独处理。}{lower-sa-timeline}


\MFQTransition{预测必须真实决定仓位}

本节先把预测明确定义为下一期价差变化 $\widehat{\Delta z}_{t+1|t}$，而非下一期价差水平。预测变化为正时做多价差，决策映射为
\[
s_t=\begin{cases}
+1,&\widehat{\Delta z}_{t+1|t}\ge c,\\
-1,&\widehat{\Delta z}_{t+1|t}\le-c,\\
0,&|\widehat{\Delta z}_{t+1|t}|<c,
\end{cases}
\qquad w_t=\min(L,|s_t|L)\operatorname{sign}(s_t).
\]
同时声明阈值 $c=0.5$、上限 $L=1$、持有期 1、每期调仓。预测 $[0.8,-0.7,0.1,0.9]$ 因而产生目标仓位 $[1,-1,0,1]$，不是从 解析参照 直接读取交易收益。

% Generated by tools/build_figures.py; edit figures/figure-specs.json instead.
\MFQFigure{tex/figures/generated/lower-sa-position-pnl.pdf}{预测先映射为有界仓位，仓位再作用于下一期收益；换手成本在仓位变化时扣除。}{lower-sa-position-pnl}


成交比例作用于订单增量，而不是目标仓位本身。令 $p_{t-1}$ 为旧仓位、$f_t$ 为成交比例，则
\[
q_t=w_t-p_{t-1},\qquad p_t=p_{t-1}+f_tq_t.
\]
成交比例 $[1,0.5,1,0]$ 因而产生实际仓位 $[1,0,0,0]$：第二期从 $+1$ 翻到 $-1$ 的订单量是 $-2$，成交一半后恰好平仓。最后一笔未成交时，仓位保持上一期的 0；更一般地，零成交必须保持旧仓位，而非强制归零。毛收益逐期为 $p_t r_{t+1}$，换手从空仓开始：
\[
\text{turnover}=|1-0|+|0-1|+|0-0|+|0-0|=2.
\]
实现收益 $[0.02,-0.01,0.03,0.04]$ 给出毛收益 $0.02$；单位换手成本 $0.002$ 给出成本 $0.004$，故 $R^{\mathrm{net}}=0.016$。把成本率加倍后，净收益应再下降 $0.004$，这是计算记录的敏感性检查。

透明实现逐期暴露订单、成交与旧仓位的语义；成熟库对照把同一递推写成有限维下三角线性系统并交给 SciPy 求解。后者适合独立核对小型事先确定计算记录，但稠密矩阵会随期数平方增长，也没有队列位置、撮合优先级、冲击或停牌状态，因此不能冒充真实执行引擎。长时间轴应利用双对角稀疏结构或直接递推，市场微观结构则需要另建状态模型。

\MFQTransition{失效检测与研究材料}

研究材料必须逐窗保存：估计截止、模型参数、残差诊断、状态概率、预测、目标仓位、成交比例、实际仓位、换手、成本与净收益。失效规则可包括协整残差单位根恶化、半衰期超出持有期、状态概率熵升高、变点警报和成交率下降。规则必须在交易期前事先确定；用全样本挑选“最早正确警报”没有意义。

错误研究材料审计要追问：预测是否真实进入仓位？收益日期是否晚于决策？未成交是否仍计收益？成本是否按实际仓位变化计算？scaler、协整向量与阈值是否只在训练窗口拟合？这些问题比多加一个模型名称更能决定研究可信度。

若数据不支持交易时点，项目应停在方法演示，不得包装成回测。

\MFQTransition{从价差 z-score 到目标仓位}

训练窗口估计价差均值和波动后，$z_t=(s_t-\hat\mu_t)/\hat\sigma_t$。最简单规则在 $z_t>z_{\rm entry}$ 时做空价差、$z_t<-z_{\rm entry}$ 时做多，在 $|z_t|<z_{\rm exit}$ 时平仓。入口和出口不同形成 hysteresis，减少阈值附近反复交易。

仓位可固定，也可按预期回归幅度、风险和半衰期缩放。无论哪种，必须设最大持有期、止损、关系失效和不可交易处理。参数只在内层验证选择；用全样本最优 z-score 是数据挖掘。

\MFQTransition{配对交易的自融资计算记录}

若价差为 $s_t=y_t-\beta_tx_t$，做空一单位价差意味着卖出 $y$、买入 $\beta_t$ 单位 $x$。beta 变化会触发再平衡成本。收益应从两腿实际持仓与下一期价格变化计算：
\[
\Pi_{t+1}-\Pi_t=w_{y,t}\Delta y_{t+1}
+w_{x,t}\Delta x_{t+1}-C_{t+1}.
\]
直接用 $-\Delta s$ 当收益只在 beta 固定、单位一致且忽略融资/成本时成立。

\MFQTransition{候选选择与多重检验}

从 1000 只股票枚举约 50 万对，再挑最平稳价差，会产生巨大选择偏差。先用行业、供应链或共享经济暴露缩小候选，再在训练期做协整检验；完整记录尝试数，并在独立未来期评价整个筛选程序。

篮子套利用 Johansen 或稀疏方法寻找多资产组合，实施复杂度随腿数上升：更多点差、借券和异步成交。统计上更平稳的篮子不一定更可交易。

\MFQTransition{收益分解与压力测试}

将净收益拆成均值回归收益、beta 更新、融资、点差、冲击、借券和未成交。压力测试包括相关关系突然断裂、波动翻倍、半衰期延长、单腿停牌和借券召回。失效规则应在压力前定义，而非根据损失路径挑一个恰好提前退出的阈值。

\begin{diagnostic}
如果预测序列与策略持仓没有确定映射，或收益来自 合成样本 中预先给定的 \texttt{trade\_return}，那只是并排展示模型和交易数字，不是模型驱动策略。
\end{diagnostic}

\section{本路线的综合项目}

选定一条协整或状态切换关系，沿时间轴完成候选选择、滚动估计、仓位映射和成本后评价。综合报告应同时展示关系失效、参数漂移、未成交和借券约束的一个反例，并把结论限定在声明的资产池与研究期内。

\subsection{变动 beta 为什么产生虚假的价差收益}

设 $x_t=50,y_t=101,\beta_t=2$，下一期 $x_{t+1}=51,y_{t+1}=102$。做空一单位价差的持仓为 $w_y=-1,w_x=2$，不计成本的损益为 $-1(102-101)+2(51-50)=1$。它等于固定 beta 下的 $-\Delta(y-2x)$。

若下一期重新估计得到 $\beta_{t+1}=2.1$，新定义价差为 $102-2.1(51)=-5.1$，相较旧价差 1 下降 6.1，却不是本期赚了 6.1。恒等式为
\[
 \Delta(y-\beta x)=\Delta y-\beta_t\Delta x-x_{t+1}\Delta\beta.
\]
最后一项 $-51(0.1)=-5.1$ 是定义更新，不是旧持仓损益。要把持仓调到新的 beta，还需买入 $0.1$ 单位 $x$，并在现金账户记录支出及成本。notebook 将直接显示两腿股数、价格变化与现金调整，避免用重新拟合的历史价差替代真实交易。

\section{综合练习}
\begin{enumerate}[leftmargin=*]
  \item \textbf{口述概念：}为什么 ADF 的 t 统计量不能查普通 t 分布？
  \item \textbf{概念：}高相关与协整分别约束什么？
  \item \textbf{笔试推导：}由离散 AR 系数 $\phi$ 推导 OU 的 $\kappa$ 与半衰期。
  \item \textbf{推导：}从 VAR 写出 VECM，并解释 $\Pi=\alpha\beta^{\mathsf T}$。
  \item \textbf{数值编程：}改变残差 $\phi$，画半衰期敏感性并标记无效区间。
  \item \textbf{编程：}比较透明 DF、\texttt{adfuller} 与 Engle--Granger 专用临界值。
  \item \textbf{研究判断：}一份报告先从 500 对资产中选最平稳价差，再对所选对报告 p 值。指出选择偏差并设计修复。
\end{enumerate}

\subsection*{基础衔接：独立计算}
$\phi=0.5$，采样间隔两天，求 OU 均值偏离的半衰期。

\subsection*{模拟频率与尺度}
\begin{enumerate}[leftmargin=*]
\item 共同趋势加平稳价差与独立随机游走，两组拒绝频率各估计什么？拒绝率 0.5 时，60 次和 300 次重复的模拟标准误各为多少？
\item 固定 $D,z$，将 $y=D\beta+cz$ 的 $c$ 从 0.5 改为 2，推导系数误差、残差与固定规格 DF 统计量的变化。
\end{enumerate}

\section{分层练习}
\begin{enumerate}[leftmargin=*]
  \item \textbf{口述概念：}写出预测、滤波、平滑各自条件的信息集。
  \item \textbf{概念：}为什么 regime 标签会交换？
  \item \textbf{笔试推导：}推导标量 Kalman 增益与后验方差。
  \item \textbf{推导：}写出两状态过滤的一步预测与 Bayes 更新。
  \item \textbf{数值编程：}改变 $Q/R$，记录滤波路径的响应速度。
  \item \textbf{编程：}用多组初值拟合 Markov switching，比较似然与状态解释。
  \item \textbf{研究判断：}报告用全样本平滑状态回填历史交易。构造能执行的负例。
\end{enumerate}

\MFQTransition{仓位与收益计算与应用分析题组}
\begin{enumerate}[leftmargin=*]
  \item \textbf{口述概念：}purge 与 embargo 分别阻断什么信息路径？
  \item \textbf{概念：}为什么目标仓位不能直接用于收益计算？
  \item \textbf{笔试推导：}逐期重建本章 $0.02$ 毛收益、2 单位换手和 $0.016$ 净收益。
  \item \textbf{数值编程：}把持有期改为 2，解释仓位和换手变化。
  \item \textbf{编程：}加入借券失败与停牌成交比例，扩展失败状态而非删除样本。
  \item \textbf{研究判断：}一项研究用完整样本拟合标准化参数，并按未成交的目标仓位计算收益。分别说明两种做法改变了什么信息或持仓假设，并重新写出正确的计算顺序。
  \item \textbf{综合项目：}选择一组可用数据，估计价差并制定仓位规则。比较简单基线，计算部分成交和成本后的结果，再分析关系变化时会发生什么。
\end{enumerate}

\subsection*{基础衔接：独立计算}
若上述交易的两腿总成本为 0.2，净损益是多少？下一期 beta 改变是否改变本期已实现损益？


\chapter{模型与表示：从基线到 PyTorch、Transformer 与文本适配}
\label{ch:lower-ml-model}

\begin{itemize}
  \item 直接先修：上册第 5、6、9、10、12、14、15、16、17 章。
  \item 学习产出：能从 NumPy 基线迁移到可复现的 PyTorch 训练，并证明序列模型确实保留顺序。
  \item 研究边界：本章的小模型用于暴露训练与信息协议，不声称达到生产模型容量。
\end{itemize}

\section{先冻结简单基线}

回归树桩在候选阈值 $c$ 上最小化两侧平方误差。候选集合必须来自排序后的唯一特征值中点；若直接使用输入相邻项，中点集合会依赖行顺序。透明实现先计算
\[
\mathcal C=\left\{\frac{x_{(j)}+x_{(j+1)}}2:x_{(j)}<x_{(j+1)}\right\},
\]
再逐个比较。成熟库对照使用 \texttt{DecisionTreeRegressor(max\_depth=1)}；两者在冻结数据上的预测差为 0。两轮残差 boosting 也与 \texttt{GradientBoostingRegressor} 对拍。这里的“双实现”只验证同一平方损失与树深，不代表不同库默认值天然等价。

线性、树和 boosting 不是被深度模型淘汰的装饰品。它们提供训练时间、解释难度和样本外损失的下界；复杂模型若不能在冻结验证协议下改善决策结果，应保留基线。资产定价中的系统比较也表明，模型灵活性必须与正则化和严格样本外评价一起解释\cite{gu2020mlassetpricing}。

\section{PyTorch：从张量到 checkpoint}

张量同时携带数值、形状、dtype 和 device。训练循环的最小闭环分为三步：
\begin{enumerate}
  \item \texttt{Dataset} 保存样本语义，\texttt{DataLoader} 决定批次与顺序；
  \item 前向传播产生损失，\texttt{autograd} 计算梯度，优化器更新参数；
  \item 在评估模式与无梯度上下文中冻结推理。
\end{enumerate}
代码还固定随机种子、CPU device、epoch 数与学习率，并把配置和参数字节 $D$ 哈希为 checkpoint 指纹 $h(D)$\cite{paszke2019pytorch}。

对四个样本训练两层 MLP，最终 MSE 约为 $6.34\times10^{-4}$。这个数不是模型质量结论；真正的证据是同一配置两次运行得到完全相同预测和指纹。若更换 GPU、算子或框架版本，必须重新声明可复现级别，而非默认逐位一致。

\begin{diagnostic}
只保存 \texttt{state\_dict} 不够：还要保存 tokenizer、特征顺序、预处理参数、随机种子、优化器状态、代码版本和数据截止。checkpoint 是推理契约，不只是权重文件。
\end{diagnostic}

\section{因果卷积、RNN、Attention 与 Transformer}

序列输入为 $X\in\mathbb R^{B\times T\times F}$。因果卷积的 $t$ 时点只读取 $\{X_s:s\le t\}$；RNN 递推 $h_t=\phi(W_xX_t+W_hh_{t-1})$；self-attention 通过
\[
\operatorname{softmax}\!\left(\frac{QK^{\mathsf T}}{\sqrt d}+M\right)V
\]
聚合位置，其中因果 mask $M$ 屏蔽未来。Transformer 再叠加前馈、残差与归一化\cite{vaswani2017attention}。没有位置编码的 attention 对置换近似等变，不能自动识别先后。

冻结负例把 $[1,2,4,8]$ 反转。时间均值差为 0，而因果卷积、RNN、带位置与因果 mask 的 Attention/Transformer 输出均发生变化。这个实验不证明模型预测有效，只证明实现没有把顺序悄悄压成均值。padding mask、缺失 mask 与因果 mask 必须分别记录。

\section{NLP/LLM：任务边界先于模型规模}

文本链从 tokenizer 和 embedding 开始，再进入预训练编码器、任务头与决策映射。零样本依赖任务提示或词义规则；少样本利用少量标注形成原型；全量微调更新编码器全部参数；PEFT 只训练 adapter、LoRA 或任务头。LoRA 的核心是冻结预训练权重并学习低秩更新\cite{hu2022lora}。冻结实验以小型编码器代理说明参数边界：可训练参数仅占总参数约 $15.8\%$，不是把它包装成基础模型。

真实预训练编码器必须绑定模型卡、许可、权重版本和 tokenizer。新闻或公告还要同时保存首次发布时间、抓取时间与修订时间；决策日后的修订不能回填历史。LLM 生成的标签需保留 prompt、采样参数、原文证据和人工抽查，零/少样本结果不能免除时间审计。

\begin{implementationnote}
\begin{lstlisting}[language=bash]
uv run python notebooks/lower/ml_alpha_model.py \
  evidence/ml-alpha-model/oracle.json
\end{lstlisting}
\end{implementationnote}

本单元向 Route Capstone 交付：基线表、PyTorch checkpoint、顺序负例、文本适配协议和算力/随机性限制。

\section{分层练习}
\begin{enumerate}[leftmargin=*]
  \item \textbf{口述概念：}tensor 的 shape、dtype、device 和梯度状态分别约束什么？
  \item \textbf{口述概念：}为什么树桩候选阈值必须先排序唯一特征值？
  \item \textbf{笔试推导：}写出两层 MLP 的前向传播与平方损失梯度链。
  \item \textbf{笔试推导：}解释没有位置编码的 self-attention 对时间置换为何不敏感。
  \item \textbf{数值编程：}把序列反转，比较均值、因果卷积、RNN、Attention 与 Transformer 输出。
  \item \textbf{数值编程：}保存并恢复 checkpoint，验证预测与哈希不漂移。
  \item \textbf{研究判断：}审计一个把随机特征均值池化称作“深度序列模型”的错误研究包。
  \item \textbf{研究判断：}比较全量微调、只训任务头和 LoRA/adapter 的证据边界。
\end{enumerate}


\chapter{现代人工智能方法：表示学习、关系建模与序贯决策}
\label{ch:lower-ml-deep}

怎样把一个非线性预测写成许多简单运算的组合，又怎样求出每个参数应当调整的方向？本章先完整计算一个两层网络，再推广到矩阵形式；卷积则展示同一组参数如何在多个位置重复使用。

树模型通过划分特征空间表示非线性关系，神经网络通过复合函数学习表示。两者组织近似的方式不同，适用性取决于数据结构、样本数量和所需的预测任务。

\section{计算图、张量与反向传播}

\MFQLead{先算一个两层网络}
取标量 $h=\operatorname{ReLU}(ax+b)=\max(0,ax+b)$、$\hat y=ch+d$、$L=(\hat y-y)^2/2$，参数为 $x=2,y=1,a=1,b=-1,c=2,d=0$。前向得到 $ax+b=1,h=1,\hat y=2,L=1/2$。设误差 $e=1$，因激活输入严格为正，链式法则给
\[
 \partial_cL=eh=1,\quad \partial_dL=e=1,\quad
 \partial_aL=ecx=4,\quad \partial_bL=ec=2.
\]
学习率为 $0.1$，同时用旧参数梯度更新，得到 $(a,b,c,d)=(0.6,-1.2,1.9,-0.1)$。新激活输入为零，输出 $-0.1$、损失 $0.605$，反而上升。梯度方向只保证足够小步长的局部下降；沿一步跨过 ReLU 边界尤其不能假定下降。取学习率 $0.01$ 可得到新输出 $1.781$、损失约 $0.3053$。notebook 逐项展示梯度并扫描步长，再用一个小网络学习带噪非线性函数。

下面把同样的链式法则写成矩阵形式，以同时处理多个输入与隐藏单元。

神经网络是参数化复合函数 $f_\theta$。以两层网络为例：
\[
z_1=W_1x+b_1,\quad h=\phi(z_1),\quad
\hat y=W_2h+b_2,\quad
L=\ell(\hat y,y).
\]
反向传播沿计算图反向使用链式法则，重复利用已经算出的局部导数：
\[
\frac{\partial L}{\partial W_2}
=\frac{\partial L}{\partial\hat y}h^{\mathsf T},
\qquad
\frac{\partial L}{\partial W_1}
=\left(W_2^{\mathsf T}\frac{\partial L}{\partial\hat y}
\odot\phi'(z_1)\right)x^{\mathsf T}.
\]
自动微分也按这些局部规则累积导数。因此，手算小网络的中间梯度可以帮助理解框架返回的梯度；实际训练的目标则由我们选择的损失函数决定。

张量的形状说明每一轴组织哪些观测。金融序列常写成 $X\in\mathbb R^{B\times T\times F}$，分别表示批次、时间和特征。若把资产维与时间维交换，代码仍可能运行，但模型学习的问题已经改变。例如最后一轴含收益和成交量两个特征时，交换时间轴与特征轴就改变了哪些数在一起参与运算。

\MFQLead{激活、初始化与梯度传播}

若每层都使用线性映射，多层网络仍等价于单个线性变换。ReLU $\phi(z)=\max(0,z)$ 提供分段线性表达，但负半轴梯度为零；GELU 以平滑门控常用于 Transformer。Sigmoid 和 $\tanh$ 在绝对值大时导数接近零，长链相乘容易导致梯度消失。

初始化要稳定前向方差与反向梯度。若输入分量独立、方差为 $v$，权重方差约取 $1/n_{\rm in}$ 可避免线性层方差随宽度爆炸；ReLU 常采用 He 初始化 $2/n_{\rm in}$。该近似基于对称的零均值预激活：ReLU 将二阶矩减半；激活后均值通常不为零，因此不能把二阶矩直接称为方差。

残差连接
\[
h_{l+1}=h_l+F_l(h_l)
\]
为梯度提供恒等路径。LayerNorm 对单个样本的特征维标准化，BatchNorm 使用批次统计；金融小批次、非平稳和跨资产混合会使 BatchNorm 的运行统计难以解释，不能照搬图像默认配置。

\section{优化器、正则化和训练循环}

随机梯度下降用抽取的小批次梯度近似训练集平均损失的梯度。Momentum 累积方向，Adam 对一阶和二阶矩做自适应缩放。AdamW 将权重衰减从梯度更新中解耦。优化器改变训练动力学，不改变数据的有效信息量；在只有数十个独立月份的样本上，百万行资产面板并不等于百万个独立观察。

% Generated by tools/build_figures.py; edit figures/figure-specs.json instead.
\MFQFigure{tex/figures/generated/lower-ml-training.pdf}{本例的训练损失继续下降，验证损失却开始上升，表现出过拟合；早停使用验证结果选择迭代轮数。}{lower-ml-training}


常见正则化包括权重衰减、dropout、数据增强、早停和容量限制。早停监控的是验证窗口，不是最终测试窗口。Dropout 在训练时随机屏蔽单元，推理时必须切到 \texttt{eval()}；忘记切换会使同一输入得到随机预测。

训练可以理解为重复三步：前向计算预测与损失，反向计算梯度，再按学习率更新参数。每轮之后在独立验证窗口计算损失，观察训练误差下降是否伴随未来误差下降。选择早停轮数属于模型选择，最终测试数据不参与这一步。

\section{卷积：局部连接与参数共享}

一维卷积核 $K\in\mathbb R^{k\times F\times C}$ 在时间轴滑动：
\[
H_{t,c}=\phi\!\left(b_c+
\sum_{u=0}^{k-1}\sum_{f=1}^F K_{u,f,c}X_{t-u,f}\right).
\]
因果卷积只读取 $t$ 及以前；普通“same”卷积若对称填充，会读取未来位置。扩张卷积以间隔 $d$ 取样，使较少层覆盖更长历史。$L$ 层、核宽 $k$、扩张率 $1,2,\ldots,2^{L-1}$ 的感受野为
\[
1+(k-1)(2^L-1).
\]

% Generated by tools/build_figures.py; edit figures/figure-specs.json instead.
\MFQFigure{tex/figures/generated/lower-ml-cnn.pdf}{一维因果卷积只汇总当前及过去窗口；感受野随层数扩大，而参数在时间位置间共享。}{lower-ml-cnn}


\MFQLead{卷积的手算例}

价格变化序列为 $(1,3,2,5)$，因果核为 $(1,-1)$，按“当前减前一期”定义输出，则有效位置得到 $(2,-1,3)$。反转序列后输出变为 $(-3,1,-2)$；时间均值虽然不变，卷积保留了局部顺序。比较两种次序的输出，可以看出卷积保留了哪些局部变化。

二维卷积适合图像、相关矩阵或订单簿“价格档位 $\times$ 时间”的网格表示，但网格的邻接关系必须有经济含义。把任意因子表 reshape 成图片并不会创造空间局部性。

\MFQLead{CNN 在量化研究中的三种典型输入}

第一类是多变量时间序列：收益、成交量、波动和盘口特征沿时间排列，用一维因果 CNN 提取局部形状。第二类是限价订单簿张量：价格档、买卖方向和事件时间形成局部结构。第三类是文本或事件 token 的局部 n-gram 表示；这类任务如今常由 Transformer 接管，但 CNN 仍是速度快、延迟低的基线。

所有输入都要明确采样时钟。日线卷积中“过去 20 根 bar”较清楚；事件时间卷积中 20 个事件可能跨越毫秒或数小时，需要额外时间间隔特征。对缺失交易日补零会把“未观察”混成“收益为零”，应使用 mask 或显式缺失类别。

\begin{diagnostic}
预测时点 $t$ 的结果时，以 $t$ 为中心、包含 $t$ 之后数据的滚动窗口会使用未来信息。若整个输入窗口都已在决策前观测到，窗口内部的双向运算本身并不造成泄漏。需要分别判断原始数据的可得时间与每个输出实际依赖的位置。
\end{diagnostic}

\section{分类、回归、排序与多任务学习}

方向分类使用交叉熵，收益回归常用 MSE/Huber，横截面选股更关注排序。Pairwise loss 比较同一时点资产对，Listwise loss 直接作用于整组排序。多任务学习可共享表示，同时预测收益、波动和成交概率：
\[
L=\lambda_r L_{\rm return}+\lambda_v L_{\rm vol}+
\lambda_f L_{\rm fill}.
\]
$\lambda$ 不只是调参；它定义了研究目标的相对单位。可以用验证期决策损失选权重，但不能看到测试期收益后再改。

深度集成、bootstrap、分位数损失和 conformal 方法从不同角度估计预测不确定性，各自需要相应假设。MC dropout 产生多组预测，但这些变化是否对应可靠的概率，需要通过校准或覆盖率另行考察。后续决策还要说明预测不确定性如何影响仓位。

\MFQLead{何时深度模型没有带来价值}

复杂度本身无法说明预测是否改善。可以先比较树与网络在相同时间窗口上的预测损失，再考察不同随机初始化下的变化。若讨论交易，还应使用相同的仓位与成本约定，把预测改善与更高换手带来的影响分开。

两层网络展示链式法则如何传递误差，卷积展示同一组参数如何作用于多个位置。下一章从这些运算继续研究序列依赖与 Attention。

\section{序列模型与大语言模型：从 RNN 到 LLM}
\label{ch:lower-ml-sequence}


金融信息按时间到达，但“使用序列模型”并不自动保证时间合法。本章从递归状态开始，推导 Attention 与 Transformer，再讨论文本预训练、LLM 适配和量化研究中的证据边界。

\MFQLead{RNN：用状态压缩历史}

最简单的循环网络为
\[
h_t=\phi(W_xx_t+W_hh_{t-1}+b),\qquad
\hat y_t=W_oh_t+c.
\]
$h_t$ 是对 $x_{1:t}$ 的有限维摘要。通过时间反向传播会反复乘以 $W_h$ 和激活导数；谱半径与激活饱和共同造成梯度消失或爆炸。梯度裁剪只限制数值爆炸，不解决长期信息已被压缩的问题。

LSTM 用遗忘门、输入门和输出门控制记忆：
\[
f_t=\sigma(W_f[x_t,h_{t-1}]+b_f),\quad
i_t=\sigma(W_i[x_t,h_{t-1}]+b_i),
\]
\[
c_t=f_t\odot c_{t-1}+i_t\odot\tilde c_t,
\qquad h_t=o_t\odot\tanh(c_t).
\]
当 $f_t$ 接近 1 时，记忆梯度可沿 $c_t$ 较稳定传播。GRU 合并部分门，参数更少。两者仍需处理 padding、变长序列和隐状态跨资产串线：不同证券的 hidden state 不得因批处理顺序意外相连。

\MFQLead{Attention：内容寻址而非免费记忆}

给定 $Q=XW_Q$、$K=XW_K$、$V=XW_V$，缩放点积注意力为
\[
A=\operatorname{softmax}\!\left(\frac{QK^{\mathsf T}}{\sqrt{d_k}}+M\right),
\qquad H=AV.
\]
$1/\sqrt{d_k}$ 防止维度增大时点积方差过大、softmax 过度饱和。因果 mask 令 $M_{ts}=-\infty$（当 $s>t$），padding mask 则屏蔽不存在的 token；两者目的不同，不能共用一个布尔数组后不再检查。

% Generated by tools/build_figures.py; edit figures/figure-specs.json instead.
\MFQFigure{tex/figures/generated/lower-ml-attention.pdf}{因果 mask 将未来键位置置为不可见；每一行注意力权重只能分配给当前及过去 token。}{lower-ml-attention}




多头注意力让不同子空间学习不同关系，但头数增加不保证可解释性。注意力权重描述模型内部加权，不等同于特征的因果贡献或忠实解释。

\MFQLead{三 token 手算}

若单头、$d_k=1$，查询 $q_2=1$，键为 $(0,1,2)$，在因果位置 $t=2$ 只能看到前两个 token，则权重为
\[
\operatorname{softmax}(0,1)=
\left(\frac1{1+e},\frac e{1+e}\right).
\]
第三个 token 即使键更大也必须被 mask 成零权重。这个简单例子应成为实现因果 mask 的单元 oracle。

\MFQLead{位置、时间间隔与金融日历}

没有位置编码的 self-attention 对输入置换近似等变。绝对位置 embedding 区分第几个 token，相对位置编码表达间隔，旋转位置编码把相对位置信息嵌入点积。金融数据还需要真实时间差：两个相邻事件可能相隔 10 微秒或一夜。可以加入 $\Delta t$、交易阶段、隔夜标志和日历 embedding，但这些特征必须在事件发生时已知。

长序列的标准 Attention 计算和内存复杂度为 $O(T^2)$。稀疏/局部注意力、状态空间模型或分块缓存可以降低成本，但改变了可见历史和归纳偏置。模型延迟是交易约束，不只是工程优化指标。

\MFQLead{Transformer 块与训练稳定性}

Transformer 块由多头注意力、前馈网络、残差和归一化组成\cite{vaswani2017attention}：
\[
U=X+\operatorname{MHA}(\operatorname{LN}(X)),\qquad
Y=U+\operatorname{FFN}(\operatorname{LN}(U)).
\]
前馈层逐位置应用两层 MLP，通常先扩维再压回。Pre-LN 在深层训练中常比 Post-LN 稳定。训练仍需学习率 warmup、梯度裁剪、混合精度和 checkpoint；这些是数值协议，不应与研究有效性混为一谈。

% Generated by tools/build_figures.py; edit figures/figure-specs.json instead.
\MFQFigure{tex/figures/generated/lower-ml-transformer.pdf}{Transformer 保持批次、时间、注意力头和通道维度；残差路径连接注意力与逐位置 MLP。}{lower-ml-transformer}




编码器适合表示一段已经完整可见的历史。自回归解码器通过因果掩码预测下一个词元；编码器--解码器结构适合条件生成。用未来完整窗口做“编码器输入”再预测窗口中间标签是典型泄漏。

\MFQLead{从语言模型目标到 LLM}

自回归语言模型最大化
\[
\sum_{t=1}^T\log p_\theta(w_t\mid w_{<t}),
\]
masked language model 则重建被遮蔽 token。Tokenizer 把文本映射为 subword ID；更换 tokenizer 会改变序列长度、未知词和数值表示，因此它是 checkpoint 契约的一部分。

预训练让模型获得通用语言表示，但不会自动获得“当时可用的信息”。用于历史回测时必须记录：模型训练数据截止、权重发布日期、新闻首次发布时间、所用文本版本和决策时点。现代模型可能在预训练语料中见过历史事件的后续总结，因此即使输入文本来自过去，也可能发生参数记忆污染。

\MFQLead{微调、指令学习与参数高效适配}

全量微调更新全部参数，成本高且易在小样本上遗忘。冻结编码器只训练任务头是最稳基线。LoRA 将权重更新限制为低秩矩阵
\[
W'=W+BA,qquad A\in\mathbb R^{r\times d_{\rm in}},
\quad B\in\mathbb R^{d_{\rm out}\times r},\quad r\ll d,
\]
显著减少可训练参数\cite{hu2022lora}。Adapter、prefix tuning 和 prompt tuning 是其他参数高效路径。参数少不意味着选择偏差少：秩、目标层、学习率和 prompt 模板仍需在训练/验证协议中冻结。

指令微调和偏好对齐改变模型回答方式，不保证事实正确。量化文本任务应优先使用结构化输出 schema、引用原文片段和拒答条件。把模型生成的情绪标签直接当真值，会把提示词偏差和幻觉写进因子。

\MFQLead{检索增强、工具使用与结构化抽取}

RAG 把检索文档 $D$ 作为条件：$p(y\mid x,D)$。历史研究必须使用按决策时点构建的版本化索引；今天的知识库检索 2018 年问题会泄漏后续文档。向量库还要保存 embedding 模型、切块规则、索引时间和过滤条件。

% Generated by tools/build_figures.py; edit figures/figure-specs.json instead.
\MFQFigure{tex/figures/generated/lower-ml-llm-pipeline.pdf}{文本检索、表示与信号构造依次将语言信息转为预测输入；历史预测只能使用相应时点可得的文本与模型信息。}{lower-ml-llm-pipeline}




工具调用适合把计算交给确定性程序，例如财务比率、日期解析和单位换算。LLM 负责选择工具与组织解释，数值 oracle 由代码给出。研究包应保存 prompt、工具输入输出、模型版本、温度、重试和人工覆盖记录。

\MFQLead{LLM 的评价不能只看一个准确率}

分类/抽取任务报告精确率、召回率、校准和按时间切分的错误；生成任务评价事实一致性、引用覆盖、结构合法率和拒答质量。对最终交易决策，还要经过
\[
\text{文本}\to\text{结构化信号}\to\text{目标仓位}
\to\text{成交}\to\text{净收益}.
\]
任何只展示几条“看起来不错”的回答都不是样本外证据。

\begin{diagnostic}
LLM 最大的量化研究风险往往不是一次明显幻觉，而是不可见的时间污染：模型权重、检索索引或文本修订包含决策日之后的信息。没有版本化语料就只能做当前时点的方法演示，不能做历史收益声明。
\end{diagnostic}

本章建立了从 RNN 到 LLM 的完整对象链。下一章讨论不规则关系结构、序贯决策和仍在快速发展的前沿方向。

\section{选读：图神经网络、强化学习与多模态研究}
\label{ch:lower-ml-frontiers}

本节的可操作实验限于三节点传播和两期动态规划，分别体现关系共享与行动反馈。它们不训练大型 GNN 或深度强化学习模型。


这一节提供后续学习入口。基础目标仅为识别图上的邻域聚合、序贯决策的状态与动作，以及多模态数据的不同输入；不要求初学者凭本节掌握 GCN、GraphSAGE、GAT、DQN 或 PPO 的完整训练理论。可在学完树模型与验证章节后跳过本节，继续 Alpha 决策。需要深入时，先选一个问题，再学习对应算法、完整训练实验与专门教材。

\MFQLead{图表示与消息传递}

图 $G=(V,E)$ 由节点、边和可选属性组成。量化研究中的节点可以是公司、基金、行业或账户；边可以来自供应链、共同持股、行业隶属、新闻共现或历史相关。前四类更接近外生关系，相关边则由数据估计，必须按时间滚动并防止使用未来窗口。

% Generated by tools/build_figures.py; edit figures/figure-specs.json instead.
\MFQFigure{tex/figures/generated/lower-ml-gnn.pdf}{GNN 通过邻居消息的聚合与更新传播关系信息；每条边的含义和可得时间决定它能提供什么信息。}{lower-ml-gnn}


消息传递 GNN 的一般形式为
\[
m_v^{(l)}=\operatorname{AGG}^{(l)}
\{\psi^{(l)}(h_v^{(l)},h_u^{(l)},e_{uv}):u\in\mathcal N(v)\},
\]
\[
h_v^{(l+1)}=\phi^{(l)}(h_v^{(l)},m_v^{(l)}).
\]
聚合必须对邻居顺序不敏感，常用 sum、mean 或 attention。$L$ 层后节点能接收 $L$ 跳邻域信息；层数过深会造成 over-smoothing，使节点表示趋同。

\MFQTransition{GCN、GraphSAGE 与 GAT}

GCN 常写成
\[
H^{(l+1)}=\sigma\!\left(
\tilde D^{-1/2}\tilde A\tilde D^{-1/2}H^{(l)}W^{(l)}
\right),
\]
其中 $\tilde A=A+I$。归一化防止高阶节点简单主导聚合。GraphSAGE 对邻居采样并学习聚合器，适合大图和新节点；GAT 用注意力学习邻居权重，但权重仍不是因果影响。

\MFQTransition{三节点手算}

考虑链 $1-2-3$，每个节点加自环，标量特征为 $(1,2,4)$。若暂用简单均值聚合，第一层得到
\[
h_1'=1.5,\qquad h_2'=\frac{1+2+4}{3}=\frac73,
\qquad h_3'=3.
\]
节点 2 同时吸收两端信息。若边 $2-3$ 在预测日后才出现，把当前完整图用于历史训练就会泄漏。图快照必须像价格数据一样带生效时间。

图任务包括节点回归（预测个股收益/风险）、边预测（关系形成）、图分类（组合或市场状态）和异构图学习。供应商、客户、行业和高管是不同边类型，强行压成无类型邻接会丢失语义。关系方向、权重、发布日期和删除时间都要进入数据协议。

异构图可为每种关系 $r$ 使用独立变换：
\[
h_v^{(l+1)}=\phi\!\left(
W_0h_v^{(l)}+\sum_r\sum_{u\in\mathcal N_r(v)}
\frac{1}{c_{v,r}}W_rh_u^{(l)}
\right).
\]
它能区分“客户--供应商”和“同属行业”，但参数量随关系类型增加。样本稀疏时，关系专属矩阵可能只记住少数大公司；可以共享低秩基、合并经济含义相近的关系，或把关系属性作为消息输入。无论采用哪种方法，都要报告每类边的覆盖率、更新时间和孤立节点比例。

动态图需要区分事件发生时间、数据库入库时间和研究者首次可见时间。若边在 $t$ 时刻出现，预测 $t$ 之前节点时不能让该边参与消息传递。常见实现是为每个决策日构造快照 $G_t$，或让消息携带时间编码并仅聚合 $t_{uv}\le t$ 的边。仅按最终图切分节点无法阻止关系泄漏，因为训练节点会通过未来边把测试期信息传播回来。

链路预测常用负采样，但“没有记录的边”不一定是真负例，可能只是数据源未覆盖。随机负采样还会生成经济上不可能的公司组合，使任务虚假简单。更可靠的评估按行业、规模和时间匹配候选负边，并分别报告新边发现、旧边持续和删除边识别。AUC 很高但只能识别数据覆盖差异的模型，没有形成可交易关系信号。

GNN 还面对 over-squashing：远处大量信息被压入固定维度向量，梯度和信号都可能衰减。增加层数同时加剧 over-smoothing 与过拟合。残差、跳连、图重连和子图采样可以缓解数值问题，却不能替代“究竟需要几跳经济关系”的研究假设。每增加一跳，都应做删除该跳关系的消融并检查样本外稳定性。

\MFQTransition{强化学习：预测问题变成控制问题}

马尔可夫决策过程由 $(\mathcal S,\mathcal A,P,r,\gamma)$ 构成。策略 $\pi(a\mid s)$ 的价值为
\[
V^\pi(s)=\mathbb E_\pi\left[\sum_{t=0}^{\infty}
\gamma^t r_{t+1}\mid S_0=s\right],
\]
动作价值满足 Bellman 方程
\[
Q^\pi(s,a)=\mathbb E\left[r_{t+1}+\gamma
\mathbb E_{a'\sim\pi}Q^\pi(S_{t+1},a')\mid s,a\right].
\]
在交易中，状态可能包含价格、库存和剩余时间，动作是订单或目标仓位，奖励必须扣除成本、风险和约束罚项。若奖励只用账面收益，代理会通过极端杠杆或不成交订单钻空子。

% Generated by tools/build_figures.py; edit figures/figure-specs.json instead.
\MFQFigure{tex/figures/generated/lower-ml-drl.pdf}{策略选择动作，动作影响状态与奖励，奖励再用于改进策略；离线历史数据没有直接给出未执行动作的后果。}{lower-ml-drl}


\MFQTransition{两步执行算例}

需在两期买入 2 单位，第一期最多买 1 单位，第二期最多买 2 单位。第一期立即买 1 单位成本为 1，等待到第二期买剩余单位的期望成本为 3；若第一期全部等待，第二期买 2 单位期望成本为 6。无折扣时，第一期买 1 的总期望成本为 $1+3=4$，优于等待的 6。这个小型动态规划是复杂 RL 环境的 解析参照：训练出的策略若连它都不能恢复，不能进入市场仿真。

\MFQTransition{DQN、策略梯度与 Actor--Critic}

DQN 用神经网络近似 $Q(s,a)$，目标为
\[
y=r+\gamma\max_{a'}Q_{\theta^-}(s',a'),
\qquad L=(Q_\theta(s,a)-y)^2.
\]
经验回放打破相邻样本相关，目标网络降低自举目标的快速漂移。连续动作更适合策略梯度或 actor--critic。策略梯度定理给出
\[
\nabla_\theta J(\theta)=
\mathbb E\big[\nabla_\theta\log\pi_\theta(a_t\mid s_t)
A^\pi(s_t,a_t)\big].
\]
PPO 用裁剪概率比限制单次更新，但不能消除环境错设和样本外制度变化。

金融 RL 的难点不是调用算法，而是环境反事实不可观察：历史只记录实际采取的一条轨迹，换一个订单会改变成交和市场冲击。纯历史回放无法准确评价未执行动作。离线 RL 还面对行为策略覆盖不足；策略选择了数据中从未出现的动作时，Q 值可能任意外推。

若历史由行为策略 $\mu$ 产生，目标策略 $\pi$ 的轨迹重要性权重为
\[
w_{0:T}=\prod_{t=0}^{T}
\frac{\pi(a_t\mid s_t)}{\mu(a_t\mid s_t)}.
\]
普通重要性采样用 $w_{0:T}$ 加权回报，在策略差异或期限较长时方差会爆炸；逐步重要性采样、加权归一化和 doubly robust 估计可降低方差，但都要求行为概率可估且具有支持覆盖。若 $\mu(a\mid s)=0$ 而 $\pi(a\mid s)>0$，仅凭该离线数据无法识别目标动作价值。

保守离线 RL 在拟合 Bellman 目标之外，惩罚数据外动作的高 Q 值；行为克隆约束则限制策略偏离历史行为。这能减少外推风险，也可能保留历史策略的缺陷。改变保守程度，并与简单规则策略比较，可以观察估计价值、行为偏离和数据支持范围之间的取舍。

交易奖励通常是多目标的。一个执行奖励可写成
\[
r_t=-x_t(P_t-S_0)-\lambda_q q_t^2-\lambda_u u_t^2
-\lambda_f\mathbf 1\{t=T,\ q_t\ne0\},
\]
其中 $q_t$ 是剩余库存，$u_t$ 是本期动作。各项单位必须统一为现金或基点。改变 $\lambda_q$ 不只是“调参”，而是在改变风险偏好；应绘制成本、完成率、尾部损失和库存暴露的 Pareto 前沿，而不是只选总奖励最大的种子。

部分可观测市场更接近 POMDP：代理只能从订单簿切片、成交和延迟推断隐含流动性状态。堆叠最近若干帧、使用循环状态或贝叶斯 belief 都是近似。若模拟器把真实隐状态直接提供给代理，得到的策略使用了额外信息。它可以作为理想信息情形的比较对象，但与仅使用可观测量的策略回答不同的问题。

\begin{diagnostic}
固定历史价格路径的回放，可用于研究价格接受者的决策。执行和做市问题还涉及动作对成交、库存及市场冲击的影响；忽略这些反馈，会改变待求解的控制问题。因此需要比较不同成交与冲击假设下的策略表现。
\end{diagnostic}

\MFQTransition{自监督学习与表示预训练}

未标注金融数据远多于可靠收益标签。自监督任务可以遮蔽重建、下一事件预测、对比学习或跨模态对齐。对比损失将同一对象的增强视图拉近、不同对象推远：
\[
L_i=-\log\frac{\exp(\operatorname{sim}(z_i,z_i^+)/\tau)}
{\sum_j\exp(\operatorname{sim}(z_i,z_j)/\tau)}.
\]
关键在于增强是否保持经济语义。随意打乱时间、缩放价格或替换行业可能改变标签含义。预训练窗口也必须早于下游评价期。

Masked autoencoder 可在收益/盘口序列上重建遮蔽片段；预测任务过容易时，模型只学局部插值。应使用线性 probe、事先确定编码器和完整微调三档比较，判断价值来自表示还是下游容量。

不同预训练目标强调不同关系：下一事件预测利用局部动态，遮蔽重建利用输入窗口内的上下文，对比学习利用增强前后应保持的特征。双向编码器能够联合读取窗口两端的信息；用于在线预测时，整个输入窗口都应已在决策前可得。比较从头训练、冻结编码器和端到端微调时，还需使用一致的语料截止与下游划分，才能区分预训练信息和参数更新的作用。

对比学习还可能发生表示塌缩：所有样本映射到近似同一点。负样本、停止梯度、方差正则或教师--学生更新可避免平凡解，但会引入新的批大小和动量选择。研究报告至少给出表示维度的方差、协方差谱、最近邻语义和下游稳定性；只报告较低预训练损失无法证明表示有用。

\MFQTransition{多模态：数值、文本、图和图像如何对齐}

公司研究同时包含财务表、价格序列、公告文本、供应链图和电话会音频。早期融合直接拼接表示，晚期融合分别预测后再组合，cross-attention 让一个模态查询另一个模态。最难的是实体和时间对齐：公告属于哪家公司、何时首次可见、财务数据何时发布、文本修订何时发生。

缺失模态不能简单补零。训练时可做 modality dropout，推理时报告可用模态集合。若某模态只在大型公司可得，模型可能把“数据可得性”当成规模信号。

跨模态对齐常以配对损失学习共同空间。例如文本表示 $z_i^T$ 与数值表示 $z_i^N$ 的对称对比损失同时让正确配对相互接近：
\[
L=\frac12\bigl[L_{T\to N}+L_{N\to T}\bigr].
\]
配对键必须是“实体--事件--首次可见时点”，而不是只按公司和自然日连接。一天内多份公告、盘前与盘后新闻、跨时区财报都可能造成一对多关系。错误对齐会让大模型学到时间邻近或公司规模，而不是事件语义。

融合模型的消融应包含单模态基线、随机打乱某模态、删除模态和时间错位负例。若加入文本后只在文本覆盖最密集的资产上改善，应进一步区分信息增量和股票池变化。语音转写和文本推理还需要时间。如果输出晚于决策时点，这些特征就无法用于原先定义的预测时刻。

\MFQTransition{生成模型、扩散模型和合成数据}

扩散模型通过逐步加噪和反向去噪学习数据分布，可生成情景或高维序列。GAN 和 VAE 也可用于市场模拟。合成数据的首要评价不是视觉逼真，而是是否保留边际分布、依赖、尾部、条件响应和策略排序，同时不复制训练样本。

生成模型不能凭空提供未观察制度下的因果反事实。用同一生成器训练和评价策略会产生“模拟器过拟合”。至少要跨生成机制评价、保留真实留出期，并明确合成结果只证明在假设环境中的行为。

扩散模型的前向过程逐步加入高斯噪声，反向模型学习噪声或 score：
\[
q(x_t\mid x_{t-1})=
N\!\left(\sqrt{1-\beta_t}x_{t-1},\beta_tI\right),
\qquad
L_{\rm simple}=\mathbb E\|\varepsilon-
\varepsilon_\theta(x_t,t,c)\|_2^2.
\]
条件 $c$ 可以是波动状态、宏观情景或资产属性。生成器若未显式条件于这些变量，抽到的极端路径可能只是在边际分布上罕见，却不对应给定压力情景。反之，条件标签来自未来全样本分类也会泄漏制度信息。

合成市场的验证至少分三层：第一层比较边际、尾部、自相关和横截面相关；第二层比较状态持续、冲击响应与成交机制；第三层比较多种事先确定策略的相对排序及失败模式。没有任何单层指标能证明“足够真实”。策略只在一个生成器上占优时，应把结论写成模拟器条件结论，而不是市场结论。

\MFQTransition{怎样继续学习}
图模型需要先掌握矩阵乘法与监督学习；强化学习需要条件期望、Markov 链与动态规划；生成模型需要概率变换与梯度优化。读者可只选择其中一个问题，先做完下面的小型计算，再进入专门教材与完整训练实验。


\subsection{Bellman 递推与图上的重复平均}
前述两期买入例中，终端成本为 $V_2(R)=3R$，第一期可选 $a\in\{0,1\}$，每单位成本为 1。因此 $V_1(2)=\min\{V_2(2),1+V_2(1)\}=\min\{6,4\}=4$。若允许第一期买 2，答案将变成 2；动作集合是问题的一部分。

三节点均值传播写成 $h'=Ah$，其中
\[
 A=\begin{pmatrix}1/2&1/2&0\\1/3&1/3&1/3\\0&1/2&1/2\end{pmatrix}.
\]
行和为 1 保证每次是邻域凸组合。反复应用会减小节点间差异；在这个连通且有自环的图上，极限由稳态权重 $(2,3,2)/7$ 决定，对初值 $(1,2,4)$ 得共同极限 $16/7$，不是普通均值 $7/3$。归一化方式因此改变传播的含义。notebook 画出每层三个节点的值，并枚举两期的全部合法动作。
\section{现代人工智能方法综合练习}
\begin{enumerate}[leftmargin=*]
  \item \textbf{口述概念：}tensor 的 shape、dtype、device 和梯度状态分别约束什么？
  \item \textbf{口述概念：}因果卷积、RNN、self-attention 和 Transformer 如何表示历史？
  \item \textbf{口述概念：}GNN 和 DRL 分别把普通监督学习问题增加了什么数学对象？
  \item \textbf{笔试推导：}写出两层 MLP 的前向传播与平方损失梯度链，并解释残差连接为何改善梯度路径。
  \item \textbf{笔试推导：}计算核 $(1,-1)$ 对序列 $(1,3,2,5)$ 的因果卷积；再说明普通对称卷积为何可能泄漏。
  \item \textbf{笔试推导：}解释没有位置编码的 self-attention 对时间置换为何不敏感，并手算两 token softmax 权重。
  \item \textbf{笔试推导：}写出 GCN 一层传播和 Bellman 方程，分别指出图快照泄漏与环境错设发生在哪里。
  \item \textbf{数值编程：}把序列反转，比较均值、因果卷积、RNN、Attention 与 Transformer 输出，并验证 padding mask 与 causal mask 的差别。

  \item \textbf{数值编程：}用两层标量网络比较学习率 $0.1$ 与 $0.01$ 的一次更新，解释损失是否下降。
  \item \textbf{数值编程：}构造三节点供应链图，比较静态全图与逐期图快照的预测差异；不得使用预测日后新增的边。
  \item \textbf{研究判断：}审计一个把均值池化随机特征称作“深度序列模型”的研究材料。
  \item \textbf{研究判断：}比较全量微调、只训任务头和 LoRA/adapter 的证据边界，并列出 LLM 历史回测的参数记忆污染。
  \item \textbf{研究判断：}设计一个执行 RL 的最小解析或独立计算参照，并说明为什么固定历史价格回放不足以评价未执行动作。
  \item \textbf{开放项目：}从序列或文本预测、动态图预测、执行决策中选一个问题，比较简单基线与所学模型，并分析成本及一个模型假设变化的影响。
\end{enumerate}

\subsection*{两层网络的全部梯度：基础练习}
上述旧参数下若 $x=0$，哪些梯度为零？

\subsection*{记忆衰减与注意力的数值含义：基础练习}
$h_t=0.5h_{t-1}+x_t$，输入 $(1,0,0)$，求三个隐状态。

\subsection*{Bellman 递推与图上的重复平均：基础练习}
若三个节点初始特征都为 5，一次均值传播后是多少？


\chapter{机器学习研究设计：验证、监控与 Alpha 决策}
\label{ch:lower-ml-validation}

\begin{itemize}
  % Generated from curriculum/manifest.json; do not edit by hand.
\item 直接先修：下册《表格统计学习：从决策树到梯度提升》。

\end{itemize}

排序正确与概率正确是不同目标。一个模型可以把所有上涨样本排在前面，却把 60\% 的上涨概率报成 99\%；下面把时间验证与概率评价分别计算。


\section{嵌套时间轴与选择权}

按信息完成时间建立训练与验证，而不只按行号。若第 $t$ 行标签在 $t+h$ 才完成，拟合时只能使用已完成标签。以下为一种保守的分块选择：
\[
 \max I_{\rm train}+h<\min I_{\rm val},\qquad
 \max I_{\rm val}+h<\min I_{\rm test}.
\]
若另外需要缓冲 $e$，应在标签完成时间之外再加 $e$。例如训练最后索引 7，$h=2$，验证从 10 开始；验证最后索引 11，其标签到 13 完成，最终测试从 14 开始。额外缓冲并非通用必需条件，不能保证数据独立。

% Generated by tools/build_figures.py; edit figures/figure-specs.json instead.
\MFQFigure{tex/figures/generated/lower-ml-nested-timeline.pdf}{内层窗口用于模型选择，外层窗口评价按这一选择得到的预测；汇总外层结果即可比较整套学习方法。}{lower-ml-nested-timeline}


内层只能在外层训练区内选择超参数；外层测试衡量整个选择程序。合法特征只可由决策时点的信息集 $\mathcal I_t$ 构造。把测试结果反馈给模型后，它就变成新的验证集。尝试次数、搜索空间和提前停止规则都属于研究选择权计算记录。

\section{Cross-fitting 与概率校准}

cross-fitting 对每个验证块只使用更早训练块拟合，再把预测写回原时间索引。透明 NumPy 路径显式构造截距、只惩罚斜率并拒绝倒序或重叠折；成熟库路径使用相同折和相同 ridge 强度的 \texttt{Ridge}。二者事先确定预测差低于数值容差。前者适合手算与暴露约束，后者适合扩展特征维度；若库默认截距、标准化或惩罚约定不同，就不能把差异归咎于浮点误差。cross-fitting 减少同一样本同时拟合 nuisance model 和评价主效应的污染，但不能自动解决非平稳或标签泄漏。事先确定两折产生索引 4--7 的四个真正样本外预测。

% Generated by tools/build_figures.py; edit figures/figure-specs.json instead.
\MFQFigure{tex/figures/generated/lower-ml-calibration.pdf}{校准图将预测概率与样本中的事件频率比较；排序能力相近的模型，概率数值仍可能有明显差异。}{lower-ml-calibration}


概率校准拟合 $\Pr(Y=1\mid s)$，例如 Platt scaling 的 logistic 映射。Brier 分数
\[
\operatorname{BS}=n^{-1}\sum_i(p_i-y_i)^2
\]
在较早校准窗口拟合 logistic 映射，再在后期评价窗口计算分数；绝不在同一批标签上自评改善。事先确定后期 Brier 的实际数值由 notebook 重建。校准改善仍不代表未来可靠；可靠性图、样本数与置信区间应同时报告\cite{niculescu2005probabilities}。

\MFQTransition{偏差、方差与金融样本的有效规模}

监督学习把泛化误差分解为不可约噪声、模型偏差和估计方差只是起点。金融面板中的行数 $N\times T$ 不能直接当作独立样本数：同一时点资产共享市场冲击，同一资产跨期相关，重叠标签又共享未来收益。模型容量应相对于独立时间块和横截面簇来判断。

学习曲线要同时沿两个方向画：增加历史长度和增加横截面资产。若只增加同一短时期的资产就改善，模型可能主要学习截面结构；若跨年份扩展后性能下降，说明制度漂移比方差下降更重要。训练误差很低、时序验证误差高是过拟合；两者都高则可能是特征、目标或容量不足。

树深、boosting 轮数、网络宽度、dropout、学习率和 lookback 都是选择权。包含 20 个候选配置、5 个随机种子和 4 种标签的研究实际上做了 400 次尝试。报告只展示获胜模型会隐藏 data snooping；至少保留候选空间、失败结果和停止规则。

\MFQTransition{嵌套 walk-forward 的完整算法}

设外层评价窗口为 $O_k$，其之前的开发数据为 $D_k$。在 $D_k$ 内再建立若干按时间前推的内层训练/验证折，选择超参数 $\lambda_k$；随后用整个 $D_k$ 重新拟合并只预测 $O_k$。不同外层折允许因市场变化选择不同超参数，但最终报告必须聚合整个选择程序，而不是从每折候选中事后再选一次。

标签 $y_t=r_{t+1:t+h}$ 跨越 $h$ 期时，任何训练样本的标签区间与验证区间相交都应 purge。Embargo 是验证/测试边界后的额外空窗，用于阻断特征或标签的缓慢传播。它们按信息区间定义，不是机械删除固定行数。

预处理器也在折内拟合：缺失填充、winsorize、标准化、PCA、目标编码、特征筛选都只能读取训练子集。把全样本 scaler 放在 pipeline 外面，是最常见且最隐蔽的泄漏之一。

\section{指标、解释与漂移响应}

回归报告 MSE、MAE、$R^2_{\mathrm{OOS}}$；横截面排序报告 Pearson IC、Rank IC、Top--Bottom spread 和分组单调性；分类报告 ROC-AUC、PR-AUC、log loss 与 Brier。类别极不平衡时 ROC-AUC 可能看似良好而精确率很差。

% Generated by tools/build_figures.py; edit figures/figure-specs.json instead.
\MFQFigure{tex/figures/generated/lower-ml-drift.pdf}{输入分布、预测分布和条件收益关系可能分别变化；区分这些变化有助于判断重新估计哪些部分。}{lower-ml-drift}


经济指标必须建立在事先确定的仓位映射后，包括净收益、波动、回撤、换手、成本、容量和暴露。模型 A 的 MSE 较低不保证组合收益更高：改善可能集中在接近零的样本，而排序头尾没有变化。评价单位应匹配决策；逐行置信区间会忽略同日相关，通常要按日期聚合或使用双向聚类、块 bootstrap。

\MFQTransition{校准、阈值与拒绝决策}

概率模型的 $0.7$ 应意味着在相似条件下约 70\% 发生，而不只是比 $0.6$ 排名更高。Platt scaling 假设 logistic 映射，isotonic 更灵活但小样本易过拟合。校准器只能在模型未用于拟合的预测上训练，常用 cross-fitted scores。

若只有预测概率超过 $q$ 才交易，$q$ 必须在验证期结合成本和容量选择。更稳健的系统允许拒绝：当不确定性过大、输入漂移或关键特征缺失时输出“无交易”，而不是强制每个样本产生方向。

\MFQTransition{解释：全局结构、局部归因与反事实}

树的 split gain、permutation importance、PDP/ALE、SHAP 和神经网络梯度回答不同问题。相关特征下，单列置换会制造不现实样本；PDP 也可能穿越数据支持之外。ALE 在局部条件区间累积效应，通常更适合相关特征，但仍是模型解释而非因果效应。

局部归因的加和恒等式只能证明解释器分配了预测差，不证明分配唯一或经济正确。应改变训练窗口、随机种子和背景样本检查重要性稳定性，并用删除特征组后的样本外损失做消融。财务特征存在会计恒等式，把市值单独减半却保持股价和股本不变不是合法反事实。

\MFQTransition{解释稳定性与漂移响应}

Top-$k$ 重要特征集合的 Jaccard 衡量两个窗口使用的特征是否一致。本章 Top-2 只有 $1/3$，必须报告不稳定；即使为 1，也不证明因果。正式研究可加入 permutation importance、SHAP 或消融，但每种方法都要声明背景分布和相关特征处理。

漂移至少区分输入漂移 $P(X)$、标签漂移 $P(Y)$、概念漂移 $P(Y\mid X)$ 与执行漂移。事先确定均值距离为 $0.75$，超过 0.5 后停止自动上线并进入重校准/重训练审批。响应动作需事先确定，不能看到亏损后再挑指标解释。

输入漂移可用 PSI、KS、Wasserstein 距离或分类器区分新旧样本；标签漂移要等标签成熟后观察；概念漂移比较条件误差或残差结构；执行漂移来自成交、成本和覆盖率变化。多指标持续监控本身又形成多重检验，阈值和联合报警规则需事前设定。

报警后的动作可分级：记录但不行动、降低仓位、回退基线、重新校准、重新训练或停止模型。重新训练不能自动吸收刚发生的异常期；需要最小样本、审批与新的外层留出。监控面板应同时显示数据质量、模型输出、组合暴露和执行状态，避免把所有问题都归咎于模型。

\begin{heuristic}
漂移统计量回答“分布是否变化”，不是“变化是否有经济意义”。极显著但幅度很小的漂移可能无关紧要；样本少但影响成交规则的制度变化可能必须立即停用。
\end{heuristic}

本单元向路线 Capstone 交付：嵌套切分、cross-fitted 预测、校准报告、解释稳定性和漂移响应记录。
\subsection{Brier 分数为何鼓励诚实概率}

设真实成功概率为 $q$，预测为 $p$。条件期望损失为
\[
 \mathbb E(p-Y)^2=q(p-1)^2+(1-q)p^2
 =(p-q)^2+q(1-q).
\]
第二项与预测无关，第一项在且仅在 $p=q$ 时最小，所以 Brier 是严格适当评分。若 $q=0.6$，报告 $p=0.6$ 的期望损失为 $0.24$，报告 $0.9$ 为 $0.33$；更极端不等于更好。

有限评价样本 $p=(0.2,0.8,0.6,0.4)$、$y=(0,1,1,0)$ 的 Brier 为 $(0.04+0.04+0.16+0.16)/4=0.1$。把概率分箱后比较频率可以展示校准，但每箱样本少时波动很大。notebook 用较早数据拟合预测，后续数据拟合校准器，再在更晚数据评价；不能在同一组标签上调概率又宣称校准改善。

\section{从预测损失到组合净收益}
\label{ch:lower-ml-research}

\begin{itemize}
  % Generated from curriculum/manifest.json; do not edit by hand.
\item 直接先修：下册《验证与监控：嵌套时序、校准、解释与漂移》。

\end{itemize}

同一个分数可以对应不同仓位，因而对应不同收益。先指定损失对应的预测对象，再用成本和风险决定是否交易，比直接最大化一条回测曲线更容易解释。


\MFQTransition{预测目标与决策目标}

% Generated by tools/build_figures.py; edit figures/figure-specs.json instead.
\MFQFigure{tex/figures/generated/lower-ml-decision-pnl.pdf}{模型分数经过阈值、仓位上限和成本映射后才形成净收益；预测误差较小并不保证净收益为正。}{lower-ml-decision-pnl}


MSE 衡量点预测误差；pairwise ranking loss 惩罚收益次序与分数次序冲突；收益加权目标直接奖励 $s_ir_i$。三个目标回答不同问题：
\[
L_{\rm rank}=\frac1{|P|}\sum_{(i,j)\in P}
\log\!\left(1+e^{-\operatorname{sign}(r_i-r_j)(s_i-s_j)}\right),
\qquad L_R=-\frac1n\sum_i s_ir_i.
\]
横截面 Alpha 常关心排序，却最终承担成本后的组合收益，预测误差、排序表现与决策收益因此分别回答不同的问题。收益加权损失若不限制杠杆会鼓励无界分数，不能替代组合约束。

\MFQTransition{例：概率与交易阈值}

假设上涨时每单位盈利 $a>0$，下跌时亏损 $b>0$，往返成本为 $c\geq0$，预测上涨概率为 $p$。买入的条件期望净收益为 $pa-(1-p)b-c$，所以仅在
\[
 p>\frac{b+c}{a+b}
\]
时为正。$a=b=1\%$、$c=0.2\%$ 时阈值是 0.6，准确率超过一半仍不够。此推导假定预测概率校准、盈亏幅度固定且成本正确；改变幅度分布后要重新求条件期望。

若连续仓位 $w$ 的收益目标为 $\mu w-\gamma\sigma^2w^2/2-c|w|$，从空仓出发，最优解为 $\operatorname{soft}(\mu,c)/(\gamma\sigma^2)$，再按仓位上限截取。成本使 $|\mu|\leq c$ 的信号不交易。notebook 会画出概率阈值和无交易区，并逐期更新部分成交后的实际仓位。

\MFQTransition{从预测函数到组合策略的六个映射}

这些变量之间的关系可以写为
\[
x_t\xrightarrow{f_\theta}s_t
\xrightarrow{g}\tilde w_t
\xrightarrow{\Pi_\mathcal C}w_t^*
\xrightarrow{\mathrm{orders}}q_t
\xrightarrow{\mathrm{fills}}w_t
\xrightarrow{r_{t+1}}R_{t+1}^{\rm net}.
\]
$s_t$ 是分数，$g$ 可做排序、z-score 或概率到预期收益映射；$\Pi_\mathcal C$ 把目标投影到杠杆、行业、风险和个股约束；订单与成交决定实际仓位；最后才得到净收益。每一步都可能吞掉模型优势。

最简单的横截面方案是按分位数组合。更连续的方案令原始权重与去均值分数成比例：
\[
\tilde w_{i,t}=\frac{s_{i,t}-\bar s_t}
{\sum_j|s_{j,t}-\bar s_t|},
\]
再做行业中性和风险缩放。分母很小时必须拒绝或沿用旧仓位，否则微小分数噪声会被放大成满杠杆。

\MFQTransition{信号滞后、持有期与组合重叠}

预测 $t+1$ 收益不代表能以 $t$ 收盘价成交。必须明确特征截止、模型推理、订单提交和成交价格。若公告在收盘后发布，最早可交易时点通常是下一交易日。日历、停牌和涨跌停会改变持有期。

$h$ 期标签对应的策略若每天建仓，会同时持有 $h$ 个 vintage。组合收益是各 vintage 的叠加，换手也由实际总仓位变化计算；不能把每个标签样本当作独立完整交易再求平均。重叠持有还使统计标准误需要处理序列相关。

\MFQTransition{风险约束与模型不确定性}

目标权重常通过优化得到：
\[
\max_w\quad \hat\mu_t^{\mathsf T}w
-\frac\gamma2w^{\mathsf T}\Sigma_tw
-\kappa\|w-w_{t-1}\|_1,
\]
约束包括净/毛敞口、行业和风格中性、单资产上限及可交易性。机器学习只提供 $\hat\mu_t$ 或排序；协方差、成本和约束来自其他模型。

预测不确定性可以进入 shrinkage 或仓位缩放。例如令
\[
\mu_{i,t}^{\rm adj}=
\frac{\hat\mu_{i,t}}{1+\lambda\hat\sigma_{i,t}},
\]
但 $\hat\sigma$ 必须经过样本外校准。简单把神经网络 softmax 当置信度会产生过度集中。

\MFQTransition{分数、目标仓位与部分成交}

每期做多最高一只、做空最低一只，毛敞口为 2。第一期分数 $[0.9,0.2,-0.8]$ 产生目标 $[1,0,-1]$。第二期目标为 $[-1,1,0]$；从旧仓位到目标的订单增量为 $[-2,1,1]$。成交比例 $[0.5,1,0]$ 后，实际仓位为 $[0,1,-1]$：未成交的第三只仍保留旧空仓，而不是强制归零。

两期毛收益分别为 $0.03$ 与 $0.02$。实际成交换手为 $2+2=4$，单位成本 $0.001$，所以
\[
R^{\mathrm{net}}=0.05-4(0.001)=0.046.
\]
这个计算的关键是逐期更新“旧仓位加成交增量”。未成交部分不会自动消失，目标仓位也不能直接替代实际仓位。配套实验显示每期分数、目标与成交后的持仓，便于改变成交比例后重算收益。这里的比例仍是外部给定的简化假设，没有进一步模拟队列、借券或市场冲击。

\MFQTransition{成本、冲击与容量}

线性成本近似
\[
C_t^{\rm lin}=\sum_i c_{i,t}|\Delta w_{i,t}|
\]
适合佣金和半点差。市场冲击常随交易量占比非线性增长，例如平方根近似
\[
C_{i,t}^{\rm impact}\approx
\sigma_{i,t}|\Delta w_{i,t}|
\eta\sqrt{\frac{Q_{i,t}}{V_{i,t}}}.
\]
它是经验模型，参数需按市场和订单类型校准。容量分析应逐步放大资本，重算订单量、冲击和借券约束，观察净收益在何处归零，而不是把历史百分比收益线性乘资本。

成本模型也必须只用当时可得的流动性。用当天收盘后才完整知道的成交量估算当日开盘订单冲击，是未来信息。

\MFQTransition{模型集成与简单基线}

树、神经网络、文本和图模型可做简单平均、基于验证表现加权或 stacking。元模型只能读取 base model 的 out-of-fold 预测；用同一训练样本的拟合值训练 stacking 会严重乐观。高度相关模型的平均并不会带来多少分散，应报告预测相关和边际贡献。

比较集成模型时，可以保留一条预先选定的树模型或线性模型作为简单基线。复杂模型在某些窗口表现较差，是理解其适用范围的材料；根据最终评价结果反复更换基线，会使比较本身参与模型选择。

\MFQTransition{历史评价与向前观察}

历史回测考察给定数据与成交假设下的表现，向前观察则考察新数据到来时信号怎样变化。在不实际交易的观察中，成交与市场影响仍需假设。因此，不同评价方式各自提供不同信息，不能简单按一种强弱次序替代。

最终测试结果用于修改模型后，这段数据便参与了开发。之后的表现可以解释开发过程，却需要新的未来区间来评价调整后的预测能力。

\begin{diagnostic}
若收益只来自极少数日期、单一行业或无法成交的小盘股，模型指标再漂亮也不是稳健 Alpha。逐笔归因、暴露分解和流动性过滤应在宣称改进之前完成。
\end{diagnostic}

\MFQTransition{文本信息的可得时间}

文本的首次发布与所使用的修订版本都应早于决策时点。例如，1 月 3 日作出预测时，1 月 5 日才发布的修订内容尚不可用，即使该文本最初发布得更早也一样。

预训练数据也可能包含评价期之后的信息。分析历史文本时，除了文本版本，还需考虑模型训练截止与可能的记忆影响。若只能获得当前修订版本，实验可以说明一种分析方法，但难以据此重建历史可用的信息。

\section{本路线的综合项目}

选择一个明确的预测问题，比较一种较复杂的模型与简单基线，并在未参与选择的时间区间评价。将预测转为仓位后，进一步考察部分成交、成本和分布变化怎样影响结果。综合项目应解释预测改善与净收益改善之间的联系。

\section{综合练习}
\begin{enumerate}[leftmargin=*]
  \item \textbf{口述概念：}purge 与 embargo 分别阻断哪条信息路径？
  \item \textbf{口述概念：}内层验证与外层测试为何不能互换？
  \item \textbf{笔试推导：}验证事先确定索引满足两条严格不等式。
  \item \textbf{笔试推导：}写出 Brier 分数并解释校准改善为何不等于收益改善。
  \item \textbf{数值编程：}实现两折时间 cross-fitting，并把预测写回原索引。
  \item \textbf{数值编程：}改变 top-$k$，画出重要性 Jaccard 敏感性。
  \item \textbf{研究判断：}审计在最终测试集上 early stopping 后仍报告“样本外”的研究材料。
  \item \textbf{研究判断：}分别举例说明输入、标签、条件关系和成交变化怎样影响预测或收益，并讨论重新校准与重新估计各能处理什么问题。
\end{enumerate}

\subsection*{基础衔接：独立计算}
真实成功率 $q=0.3$ 时，恒定预测 0.5 比诚实预测多多少期望 Brier 损失？

\MFQLead{Alpha 决策、交易摩擦与研究包题组}
\begin{enumerate}[leftmargin=*]
  \item \textbf{口述概念：}MSE、排序损失与收益加权损失分别服务什么决策？
  \item \textbf{口述概念：}为什么未成交不能把旧仓位直接归零？
  \item \textbf{笔试推导：}重建两期目标、订单增量、实际仓位与 0.046 净收益。
  \item \textbf{笔试推导：}解释无约束收益加权目标为何可能鼓励无界分数。
  \item \textbf{数值编程：}加入仓位上限和借券失败，重算 turnover 与成本。
  \item \textbf{数值编程：}对文本首次发布、抓取与修订时间注入三个泄漏负例。
  \item \textbf{研究判断：}审计一个预测优秀但成本后为负、仍宣称可交易的错误研究包。
  \item \textbf{开放 Capstone：}交付模型、验证、文本协议、组合、摩擦、漂移与限制的一键研究包。
\end{enumerate}


\chapter{随机分析与无套利：从二次变差到风险中性测度}
\label{ch:lower-derivatives-stochastic}

若价格不断受到随机冲击，普通微积分还能计算持仓收益吗？若能通过交易复制期权，价格又为何不需要直接代入股票的历史平均收益率？这两个问题把本章连在一起：先建立随机积分，再由链式法则的修正推导复制价格。

先修是上册条件期望、平方可积收敛、Brownian 运动与多元微分。目标是推导 $W_t^2$ 和几何 Brownian 运动的公式，理解自融资复制如何导出定价 PDE。一般测度变换与市场完备性放在选读部分，引用的定理不在本章完整证明。

% Generated by tools/build_figures.py; edit figures/figure-specs.json instead.
\MFQFigure{tex/figures/generated/lower-dv-route.pdf}{随机分析描述价格变化与复制关系，数值方法近似求价，离散对冲则展示有限调仓频率下的风险。}{lower-dv-route}


\section{为什么平方增量不能忽略}
标准 Brownian 运动 $W$ 从零出发，具有连续路径与独立正态增量。对确定分割 $\pi=\{0=t_0<\cdots<t_m=T\}$，令 $h_i=t_{i+1}-t_i$，$\Delta W_i=W_{t_{i+1}}-W_{t_i}$，$Q_\pi=\sum_i(\Delta W_i)^2$。正态二阶、四阶矩给出
\[
\mathbb E(\Delta W_i)^2=h_i,\qquad \mathbb E(\Delta W_i)^4=3h_i^2.
\]
增量独立，因而
\[
\mathbb E Q_\pi=T,\qquad \operatorname{Var}(Q_\pi)=2\sum_i h_i^2\le2T|\pi|.
\]
所以网格宽度 $|\pi|\to0$ 时，$Q_\pi\to T$ 在 $L^2$ 中成立，也依概率成立。这解释了平方增量为何不能忽略。光滑路径的平方增量和则趋于零。

% Generated by tools/build_figures.py; edit figures/figure-specs.json instead.
\MFQFigure{tex/figures/generated/lower-dv-qv.pdf}{固定起点与终点时，所有一阶增量之和恒为端点之差；Brownian 平方增量和则随分割加密，在概率意义下趋于时间长度。}{lower-dv-qv}

图中将同一细路径的增量聚合成粗增量。有限路径上的误差不必单调减小；在多条路径上重复才看得见误差分布的收缩。均匀 $m$ 步的理论方差是 $2T^2/m$。

每个有限分割都满足代数恒等式
\[
W_T^2=2\sum_iW_{t_i}\Delta W_i+\sum_i(\Delta W_i)^2.
\]
因此中间的和若有积分极限，就应满足
\[
\int_0^TW_t\,dW_t=\tfrac12(W_T^2-T).
\]
它的期望为零，与当前持仓无法预知下一次独立冲击相符。下一节给出这个极限的构造。

\section{从可用信息构造随机积分}
滤过 $(\mathcal F_t)$ 表示截至 $t$ 的信息，并假定 $W$ 相对该滤过为 Brownian 运动。简单过程
\[
H_t=\sum_iH_i\mathbf1_{(t_i,t_{i+1}]}(t)
\]
中 $H_i$ 对 $\mathcal F_{t_i}$ 可测，$\mathbb EH_i^2<\infty$。它表示每段开始时决定持仓。定义 $I(H)=\sum_iH_i\Delta W_i$。条件于段首信息，每个增量均值为零，所以 $\mathbb EI(H)=0$。

平方展开中的交叉项也为零：对 $i<j$，$H_i\Delta W_iH_j$ 在 $t_j$ 已知，再对最后的 $\Delta W_j$ 条件化即可。对角项为 $\mathbb E[H_i^2(\Delta W_i)^2]=h_i\mathbb EH_i^2$。因此
\[
\mathbb E[I(H)^2]=\mathbb E\int_0^TH_t^2dt.
\]
这就是 Itô 等距。若简单可预测过程 $H^{(n)}$ 在右侧范数下逼近 $H$，等距保证 $I(H^{(n)})$ 是 $L^2$ Cauchy 列。利用完备性定义积分为其极限；再次用等距可知极限与逼近方式无关。由此得到平方可积可预测过程的 Itô 积分。渐进可测且满足同样可积条件的过程也可使用这一构造；单独写“适应”还不足以交代时间可测性。

对 $H_t=W_t$，$\mathbb E\int_0^TW_t^2dt=T^2/2<\infty$。左端阶梯近似的均方积分误差不超过 $T|\pi|$，所以前一节的和确实收敛到该积分，$W_T^2$ 的公式得到证明。

\begin{example}[看见未来会改变什么]
若先观察 $\Delta W_i$ 再选 $H_i=\operatorname{sign}(\Delta W_i)$，乘积变为 $|\Delta W_i|$，期望为 $\sqrt{2h_i/\pi}>0$。它不满足段首信息约束，因而不是上述可交易简单过程。
\end{example}

\section{Itô 公式与几何 Brownian 运动}
设 $dX_t=\mu_tdt+\sigma_tdW_t$，系数满足积分的可测性与局部可积条件，$f\in C^{1,2}$。Itô 公式为
\[
df(t,X_t)=\left(f_t+\mu_tf_x+\tfrac12\sigma_t^2f_{xx}\right)dt+\sigma_tf_xdW_t.
\]
从 Taylor 展开理解它：一阶项给出 $f_t\Delta t+f_x\Delta X$，二阶项含 $f_{xx}(\Delta X)^2/2$。Brownian 增量的平方累积保留为 $dt$，时间增量的平方及交叉项消失。一般公式的严格证明还需控制余项并作局部化；这里的解释不代替这些步骤。

对 $dS_t=\mu S_tdt+\sigma S_tdW_t$、$S_0>0$，取 $f(s)=\log s$，由 $f'=1/s,f''=-1/s^2$ 得
\[
d\log S_t=(\mu-\sigma^2/2)dt+\sigma dW_t,\qquad
S_T=S_0\exp\{(\mu-\sigma^2/2)T+\sigma W_T\}.
\]
反过来对这个显式正过程用 Itô 公式可验证原 SDE，也保证取对数合法。正态指数矩给出 $\mathbb ES_T=S_0e^{\mu T}$；对数的期望却是 $\log S_0+(\mu-\sigma^2/2)T$。价格均值和对数均值是不同的量。

\section{复制如何导出定价方程}
假定无股息股票服从上述模型，$\sigma>0$，银行账户 $B_t=e^{rt}$，允许无摩擦连续交易。用 $\Delta_t$ 股股票与银行账户组成价值 $V_t$ 的自融资组合，则
\[
dV_t=\Delta_tdS_t+r(V_t-\Delta_tS_t)dt.
\]
自融资表示调仓由组合内部支付，并不是对 $\Delta_tS_t$ 求导后丢掉调仓项。若期权价值为光滑函数 $V(t,S)$，则
\[
dV=(V_t+\mu SV_S+\tfrac12\sigma^2S^2V_{SS})dt+\sigma SV_SdW.
\]
比较随机项得 $\Delta=V_S$；比较漂移项，两边的 $\mu SV_S$ 抵消，于是
\[
V_t+\tfrac12\sigma^2S^2V_{SS}+rSV_S-rV=0,\qquad V(T,S)=g(S).
\]
历史漂移退出公式，是因为复制消去了股票风险。在适当增长条件下，Feynman--Kac 定理给出
\[
V(t,S_t)=e^{-r(T-t)}\mathbb E^{\mathbb Q}[g(S_T)\mid\mathcal F_t],
\qquad dS_t=rS_tdt+\sigma S_tdW_t^{\mathbb Q}.
\]
到期收益不光滑时，在到期前解释 PDE，用终端极限指定边界。取 $g(S)=S$，条件期望为 $S_te^{r(T-t)}$，贴现价格恰为 $S_t$，与持有一股的复制相符。

% Generated by tools/build_figures.py; edit figures/figure-specs.json instead.
\MFQFigure{tex/figures/generated/lower-dv-replication.pdf}{股票与现金的动态组合匹配衍生品价值变化；自融资复制关系给出无套利价格。}{lower-dv-replication}

图把收益、动态持仓和现金账户联系起来。实际离散调仓留下对冲误差；下一章将讨论它与定价离散误差的区别\cite{black1973pricing,merton1973rational,shreve2004stochastic}。

\section{选读：风险中性测度从哪里来}
常数参数下令 $\theta=(\mu-r)/\sigma$，定义
\[
Z_T=\exp(-\theta W_T-\tfrac12\theta^2T),\qquad d\mathbb Q=Z_Td\mathbb P.
\]
正态指数矩给出 $\mathbb EZ_T=1$，且 $Z_T>0$，所以定义了与原测度等价的概率。独立增量还给出 $\mathbb E[Z_T\mid\mathcal F_t]=Z_t$。Girsanov 定理说明 $W_t^{\mathbb Q}=W_t+\theta t$ 在新测度下为 Brownian 运动。代入 $dW=dW^{\mathbb Q}-\theta dt$，
\[
dS/S=(\mu-\sigma\theta)dt+\sigma dW^{\mathbb Q}=r\,dt+\sigma dW^{\mathbb Q}.
\]
例如 $\mu=0.08,r=0.02,\sigma=0.2$，则 $\theta=0.3$。新测度改变路径的权重，并非把已发生的路径移动。

% Generated by tools/build_figures.py; edit figures/figure-specs.json instead.
\MFQFigure{tex/figures/generated/lower-dv-measure.pdf}{Girsanov 变换通过密度过程改变路径权重，使可交易资产的贴现价格在风险中性测度下成为鞅。}{lower-dv-measure}

对随机可预测 $\theta_t$，密度改为随机指数。常用充分条件是 Novikov 条件
\[
\mathbb E\exp\left(\tfrac12\int_0^T\theta_t^2dt\right)<\infty.
\]
它保证密度过程是真鞅。任意有限常数 $\theta$ 在有限期限上均满足此条件，没有固定能量上限。若 $\theta_t=(T-t)^{-1/2}$，则 $\int_0^{T-\varepsilon}\theta_t^2dt=\log(T/\varepsilon)$，完整端点处不再平方可积，不能直接用此式定义终端积分。一般而言，Novikov 不成立表示这个充分条件不可用，不等价于所有测度变换均不存在。

\MFQLead{多维与完备性的阅读边界}
若 $dX=b\,dt+\sigma\,dW$，二阶项为 $\tfrac12\operatorname{tr}(\sigma\sigma^T\nabla^2f)dt$，即协变差与 Hessian 的配对。一般连续时间无套利定理需要 NFLVR 等技术条件，本章不作证明。在 Brownian 自然滤过下的标准 Black--Scholes 模型，可复制性与鞅表示相联系；增加不可交易风险源通常会使价格不再由复制唯一确定。更换正的计价资产还会更换相应测度。这些是进一步学习的入口。

\section{章末练习}
\begin{enumerate}
  \item \textbf{口述概念：}区分依概率收敛、单路径数值接近和离散恒等式。
  \item \textbf{口述概念：}为什么 Itô 被积过程的适应性是信息约束？
  \item \textbf{笔试推导：}由离散平方恒等式推导 $W_T^2$ 的 Itô 表达。
  \item \textbf{笔试推导：}对 GBM 应用 Itô 公式推导 $\log S_T$。
  \item \textbf{数值编程：}从同一细网格构造四级嵌套分割并画出二次变差。
  \item \textbf{数值编程：}验证 $\mu-\sigma\theta=r$，再注入错误符号。
  \item \textbf{研究判断：}若风险价格为有限常数，说明为何 Novikov 条件成立；再讨论“观察到一条路径的能量为 100”是否足以判断指数矩。
\end{enumerate}
% MFQ source replacements removed=1 unit=DerivativesStochastic

% Source-mapped interview replacements.  The questions are independent Chinese
% adaptations; only the short source titles are retained for auditability.
\providecommand{\MFQInterviewFollowup}[1]{\par\noindent\textbf{面试官追问：}#1}

\providecommand{\MFQMappedUpperChFiveQuestions}{
\subsection*{笔试推导（来源映射替换）}
\begin{enumerate}[nosep,leftmargin=*]
\item \MFQInterviewMeta{笔试推导}{12 分钟}{中}{一阶条件与边界比较}{GB-3-02 Maximum and minimum}给定 $f(x)=x^4-4x^2+cx$，写出随参数 $c$ 变化的驻点判定流程；说明只解 $f'(x)=0$ 为什么不足以确定全局最小值。
\end{enumerate}
\subsection*{研究判断（来源映射替换）}
\begin{enumerate}[nosep,leftmargin=*]
\item \MFQInterviewMeta{研究判断}{12 分钟}{中}{Taylor 余项与误差控制}{GB-3-08 Taylor series}在 $x=0$ 展开 $\log(1+x)$ 到三阶，用余项给出 $|x|\le0.2$ 的统一误差界；据此审计“在所有收益率上都可安全替代对数收益”的报告。
\end{enumerate}}
\providecommand{\MFQMappedUpperChFiveHints}{\par\textbf{来源映射题提示：}先比较全部驻点、边界与无穷远行为；Taylor 题用四阶导数在区间上的上界控制 Lagrange 余项。}
\providecommand{\MFQMappedUpperChFiveSolutions}{
\section*{来源映射替换题}
\begin{enumerate}[leftmargin=*]
\item \textbf{GB-3-02。}$f'(x)=4x^3-8x+c$ 只给候选点；还需用 $f''(x)=12x^2-8$ 分类，并比较所有局部极小值。因首项系数为正，$f(x)\to\infty$，故全局最小一定在驻点中；若定义域有边界还要加入端点。参数追问应讨论三次方程重根处的分岔，而不能只给数值根。
\item \textbf{GB-3-08。}$\log(1+x)=x-x^2/2+x^3/3+R_4$，且 $f^{(4)}(u)=-6/(1+u)^4$。在 $|x|\le0.2$ 上 $|R_4|\le |x|^4/[4(0.8)^4]$。下一项为负，因此正 $x$ 附近三阶截断高估对数收益；负 $x$ 时应直接依余项而不是机械套交错级数。
\end{enumerate}\MFQInterviewFollowup{GB-3-02：定义域变成紧集时还要加哪些候选点？GB-3-08：如何用同一余项界选取收益率阈值？}}

\providecommand{\MFQMappedUpperChSixQuestions}{
\subsection*{笔试推导（来源映射替换）}
\begin{enumerate}[nosep,leftmargin=*]
\item \MFQInterviewMeta{笔试推导}{12 分钟}{中}{QR 与最小二乘}{GB-3-16 QR decomposition}对满列秩 $X=QR$，从残差正交性推导 $R\hat\beta=Q^Ty$，并解释为何不必形成 $X^TX$。
\end{enumerate}
\subsection*{口述概念（来源映射替换）}
\begin{enumerate}[nosep,leftmargin=*]
\item \MFQInterviewMeta{口述概念}{10 分钟}{中}{谱与经济方向}{GB-3-17 Determinant, eigenvalue and eigenvector}解释协方差矩阵的特征向量、特征值和行列式分别表达什么；零行列式为什么不等于“所有风险都为零”？
\end{enumerate}
\subsection*{研究判断（来源映射替换）}
\begin{enumerate}[nosep,leftmargin=*]
\item \MFQInterviewMeta{研究判断}{10 分钟}{中}{半正定与模型边界}{GB-3-18 Positive semidefinite and positive definite matrices}证明 $B^TB$ 半正定，再审计“将任意估计矩阵换成 $B^TB$ 就保留了原风险模型”的报告。
\end{enumerate}
\subsection*{数值编程（来源映射替换）}
\begin{enumerate}[nosep,leftmargin=*]
\item \MFQInterviewMeta{数值编程}{12 分钟}{中}{LU/Cholesky 分解与失败诊断}{GB-3-19 LU decomposition and Cholesky decomposition}对比 LU 与 Cholesky 的前提和计算结构；实现带对角门禁的 Cholesky，并用一个对称但不正定矩阵证明失败是模型证据而非“库坏了”。
\end{enumerate}}
\providecommand{\MFQMappedUpperChSixHints}{\par\textbf{来源映射题提示：}利用 $Q^TQ=I$；谱题区分零方向与全部方向；$x^TB^TBx=\|Bx\|^2$；Cholesky 的根号参数必须为正。}
\providecommand{\MFQMappedUpperChSixSolutions}{
\section*{来源映射替换题}
\begin{enumerate}[leftmargin=*]
\item \textbf{GB-3-16。}$\|y-QR\beta\|^2=\|Q^Ty-R\beta\|^2+\|(I-QQ^T)y\|^2$，故最小化第一项得到 $R\hat\beta=Q^Ty$。形成正规方程会把奇异值平方，使条件数平方。
\item \textbf{GB-3-17。}特征向量是风险二次型的正交主方向，特征值是这些方向的方差；行列式是特征值乘积，反映总体积尺度。行列式为零只说明至少一个线性组合方差为零，其他方向仍可有正风险。
\item \textbf{GB-3-18。}$x^TB^TBx=\|Bx\|_2^2\ge0$，所以半正定。其正定当且仅当 $Bx=0$ 只有零解，即 $B$ 满列秩。
\item \textbf{GB-3-19。}LU 面向一般方阵且通常需选主元；Cholesky 利用对称正定结构，将 $A$ 写成 $LL^T$。逐列更新时若 $a_{jj}-\sum_{k<j}l_{jk}^2\le0$，则正定假设失败。用 $\left(\begin{smallmatrix}1&2\\2&1\end{smallmatrix}\right)$ 时第二个枢轴为 $1-4<0$；应报告负特征方向和输入来源，不能加任意 jitter 后假装原矩阵有效。
\end{enumerate}\MFQInterviewFollowup{GB-3-16：列近似相关时 QR 为什么比正规方程稳定？GB-3-17 与 GB-3-18：零特征值对风险模型意味什么？GB-3-19：何时允许收缩而不是盲目加 jitter？}}

\providecommand{\MFQMappedUpperChSevenQuestions}{
\subsection*{笔试推导（来源映射替换）}
\begin{enumerate}[nosep,leftmargin=*]
\item \MFQInterviewMeta{笔试推导}{10 分钟}{中}{置换对称}{GB-4-03 Drunk passenger}第一位旅客随机占一个座位，后续旅客若自己的座位被占就随机选空座。证明最后一位坐到自己座位的概率，并指出递推中的最小状态。
\item \MFQInterviewMeta{笔试推导}{10 分钟}{中}{停止时与对称随机游走}{GB-4-01 Coin toss game I}公平硬币重复抛掷，正面累计领先反面两次时停止。写出吸收随机游走模型，并求从零开始先到 $+2$ 而非 $-2$ 的概率。
\item \MFQInterviewMeta{笔试推导}{12 分钟}{中}{组合计数}{GB-4-05 Poker hands}从标准扑克牌抽五张，推导恰为两对的概率；分母、点数组合、两对花色和第五张牌必须分别计数。
\item \MFQInterviewMeta{笔试推导}{10 分钟}{中}{卷积与独立性}{GB-4-37 Sum of random variables}独立 $X,Y\sim U(0,1)$，推导 $S=X+Y$ 的分段密度，并指出若只知道边际分布而不知道联合分布，为何答案不唯一。
\end{enumerate}
\subsection*{数值编程（来源映射替换）}
\begin{enumerate}[nosep,leftmargin=*]
\item \MFQInterviewMeta{数值编程}{12 分钟}{中}{指示变量与期望}{GB-4-38 Coupon collection}从 $m$ 种等概率标签中独立抽取，推导集齐全部标签的期望抽数，并写程序对比精确值与 Monte Carlo 结果。
\item \MFQInterviewMeta{数值编程}{12 分钟}{中}{碰撞概率与数值稳定}{GB-4-10 Birthday problem II}推导 $n$ 人生日至少一次重复的精确概率，并实现对数域计算；比较常用指数近似在 $n=23,50,100$ 时的误差。
\item \MFQInterviewMeta{数值编程}{12 分钟}{中}{Poisson 独立增量}{GB-4-31 Poisson process property}模拟强度 $\lambda$ 的 Poisson 过程，验证不相交区间计数近似独立且均值等于区间长度乘 $\lambda$；说明有限模拟不能证明独立性。
\end{enumerate}
\subsection*{研究判断（来源映射替换）}
\begin{enumerate}[nosep,leftmargin=*]
\item \MFQInterviewMeta{研究判断}{10 分钟}{中}{几何概率与分布假设}{GB-4-29 Meeting probability}两人在一小时内独立均匀到达、各等待十二分钟时计算相遇概率；再审计直接将该结果用于日内非均匀订单到达的报告。
\item \MFQInterviewMeta{研究判断}{12 分钟}{难}{相关违约与尾部依赖}{GB-4-39 Joint default}两个主体一年违约概率均为 $p$，相关系数给定。判断这些信息是否足以唯一确定联合违约概率；写出 Bernoulli 情形的可行界，并说明相关系数为何不能替代 copula 或条件违约模型。
\end{enumerate}}
\providecommand{\MFQMappedUpperChSevenHints}{\par\textbf{来源映射题提示：}旅客过程只追踪两个关键座位；随机游走利用反射对称；扑克牌按点数与花色分层计数；卷积看单位正方形截线长度；生日题使用对数乘积；Poisson 题区分理论性质与模拟证据；联合违约先写 Fr\'echet 界。}
\providecommand{\MFQMappedUpperChSevenSolutions}{
\section*{来源映射替换题}
\begin{enumerate}[leftmargin=*]
\item \textbf{GB-4-03。}中间被迫选择时，尚未解决的关键座位只有第一人和最后一人；先选到前者则链终止且最后一人成功，先选到后者则失败。两者在对称候选中等可能，故概率 $1/2$。逐位枚举不是必要状态。
\item \textbf{GB-4-01。}状态是领先数 $D_n\in\{-2,-1,0,1,2\}$，边界 $\pm2$ 吸收。令 $h(d)$ 为先到 $+2$ 的概率，则 $h(d)=[h(d-1)+h(d+1)]/2$、$h(-2)=0,h(2)=1$，线性解给 $h(0)=1/2$。若硬币有偏，差分方程的系数必须改变。
\item \textbf{GB-4-05。}总手数 $\binom{52}{5}$。选两种对子点数 $\binom{13}{2}$，各选两张花色 $\binom42^2$，第五张从剩余十一种点数与四种花色选 $11\cdot4$，故概率为 $\binom{13}{2}\binom42^2(44)/\binom{52}{5}$。第五张不能与对子同点，否则牌型改变。
\item \textbf{GB-4-37。}$f_S(s)=\int f_X(x)f_Y(s-x)\,dx$，积分区间长度在 $0\le s\le1$ 为 $s$，在 $1<s\le2$ 为 $2-s$，其余为零。卷积使用独立性；相同均匀边际可以通过不同 copula 产生不同和分布。
\item \textbf{GB-4-29。}令到达时刻缩放到 $[0,1]^2$，等待十二分钟即 $0.2$。不相遇区域是对角线两侧两个边长 $0.8$ 的直角三角形，总面积 $0.8^2$，故相遇概率 $1-0.64=0.36$。
\item \textbf{GB-4-38。}已有 $k$ 种时，新标签概率为 $(m-k)/m$，等待期望 $m/(m-k)$。总期望为 $m\sum_{j=1}^m1/j=mH_m\sim m\log m+\gamma m$。线性期望不要求各阶段等待独立。
\item \textbf{GB-4-10。}无重复概率为 $\prod_{k=0}^{n-1}(365-k)/365$，故重复概率为一减该乘积；实现时计算 $\sum_k\log(1-k/365)$ 再取指数。近似 $1-\exp[-n(n-1)/(730)]$ 忽略高阶碰撞项，$n$ 增大时必须报告误差。
\item \textbf{GB-4-31。}定义要求 $N(t)-N(s)\sim\mathrm{Poisson}(\lambda(t-s))$，且不相交区间增量独立。模拟可比较样本均值、方差和交叉协方差是否接近理论值，但“协方差接近零”不证明独立，且罕见尾部需要更大样本。
\item \textbf{GB-4-39。}联合概率 $q=P(D_1=D_2=1)$ 满足 $\max(0,2p-1)\le q\le p$；Bernoulli 相关系数若精确给定且两边际相同，可由 $q=p^2+\rho p(1-p)$ 恢复，但还必须落在可行界内。真实信用组合还需期限结构、共同因子和尾部依赖，单一相关系数不能跨状态外推。
\end{enumerate}\MFQInterviewFollowup{GB-4-01：有偏硬币时命中概率如何变化？GB-4-03：首位旅客按非均匀分布选座时结论如何改？GB-4-05：允许 joker 后计数如何重做？GB-4-10：闰日和非均匀出生率会破坏哪一步？GB-4-29：到达分布非均匀时如何积分？GB-4-31：非齐次强度怎样改变增量分布？GB-4-37：怎样用特征函数复核卷积？GB-4-38：非等概标签的瓶颈是什么？GB-4-39：怎样用单因子模型产生尾部相关？}}

\providecommand{\MFQMappedUpperChEightQuestions}{
\subsection*{口述概念（来源映射替换）}
\begin{enumerate}[nosep,leftmargin=*]
\item \MFQInterviewMeta{口述概念}{8 分钟}{中}{随机化提纯}{GB-4-16 Fair coin from an unfair coin}一枚硬币正面概率未知但每次独立且不变。构造无偏比特并计算每个输出比特的期望抛掷次数。
\item \MFQInterviewMeta{口述概念}{8 分钟}{中}{可辨识性与似然}{GB-4-15 Unfair coin}只观察一枚硬币的有限次结果，解释为什么不能“证明”它公平；写出估计正面概率的似然、置信区间与顺序检验所需的停止规则。
\end{enumerate}
\subsection*{笔试推导（来源映射替换）}
\begin{enumerate}[nosep,leftmargin=*]
\item \MFQInterviewMeta{笔试推导}{10 分钟}{中}{条件概率与无放回抽样}{GB-4-25 Aces}从洗匀牌堆依次翻牌，已知前两张至少有一张 A，求两张都是 A 的概率；分别比较“至少一张”与“指定第一张是 A”的条件事件。
\end{enumerate}
\subsection*{数值编程（来源映射替换）}
\begin{enumerate}[nosep,leftmargin=*]
\item \MFQInterviewMeta{数值编程}{12 分钟}{中}{Poisson 流量与暴露时间}{GB-4-28 Cars on road}假设车辆按强度 $\lambda$ 的 Poisson 过程到达长度为 $L$ 的路段、速度固定为 $v$。推导任意时刻路段车辆数分布并模拟核验；再指出拥堵使速度依赖密度时为何模型失效。
\end{enumerate}
\subsection*{研究判断（来源映射替换）}
\begin{enumerate}[nosep,leftmargin=*]
\item \MFQInterviewMeta{研究判断}{8 分钟}{中}{条件概率与信息更新}{GB-4-20 Monty Hall}主持人知道奖品位置且总会打开一扇空门时证明换门胜率；再审计“主持人随机开门也可直接套用 $2/3$”的结论。
\item \MFQInterviewMeta{研究判断}{10 分钟}{中}{条件生存与协议差异}{GB-4-24 Russian roulette}左轮有六个弹膛和相邻两发子弹。已知第一次扣动为空，比较“直接再扣”与“重新随机旋转后再扣”的风险；明确枪膛是否顺序前进这一机制条件。
\end{enumerate}}
\providecommand{\MFQMappedUpperChEightHints}{\par\textbf{来源映射题提示：}无偏提纯只接受 HT/TH；有限样本不能证明点假设；A 题先写清条件事件；路段车辆数利用到达率乘暴露时间；Monty Hall 与轮盘题都要把主持或机械协议写进似然。}
\providecommand{\MFQMappedUpperChEightSolutions}{
\section*{来源映射替换题}
\begin{enumerate}[leftmargin=*]
\item \textbf{GB-4-16。}成对抛掷，HT 输出 1、TH 输出 0，HH/TT 丢弃。两种接受序列概率同为 $p(1-p)$，故无偏；每对接受概率 $2p(1-p)$，期望对数为其倒数，期望抛数 $1/[p(1-p)]$。独立且固定 $p$ 是必要条件。
\item \textbf{GB-4-15。}若观察 $h$ 次正面、$t$ 次反面，似然为 $L(p)=p^h(1-p)^t$，极大似然估计 $\hat p=h/(h+t)$。任何有限序列在所有 $0<p<1$ 下都有正概率，故不能逻辑证明 $p=1/2$；只能给区间或在预先规定的显著性与停止规则下检验。边看边停会改变错误率。
\item \textbf{GB-4-25。}两张无序抽样中，条件事件“至少一张 A”有 $\binom{52}{2}-\binom{48}{2}=198$ 种，二者均 A 有 $\binom42=6$ 种，故概率 $6/198=1/33$。若指定第一张是 A，则第二张为 A 的概率为 $3/51=1/17$；两个条件事件不能混用。
\item \textbf{GB-4-28。}\par 一辆车停留时间为 $L/v$。稳态下，路段车辆数等于过去 $L/v$ 时间内的到达数，故服从 $\mathrm{Poisson}(\lambda L/v)$。模拟应从稳态开始或丢弃预热期。若速度随密度变化，停留时间与车辆数相互依赖，简单的无限服务器结论失效。
\item \textbf{GB-4-20。}初选中奖概率 $1/3$，未选两门合计 $2/3$；知情主持人必排除其中空门，所以换到剩余门继承 $2/3$。若主持人随机开门且可能揭奖，观察到空门的似然随奖品位置不同，必须用 Bayes 公式重新条件化。
\item \textbf{GB-4-24。}两发相邻时共有四个空膛位置。已知当前为空后，这四种位置等可能，其中只有一个位置的下一膛是子弹，故直接再扣风险为 $1/4$；重新均匀旋转后风险回到 $2/6=1/3$，所以不旋转更安全。若弹膛不是确定前进一步或初始旋转不均匀，需重写条件样本空间。
\end{enumerate}\MFQInterviewFollowup{GB-4-15：序贯检验如何控制可选停止？GB-4-16：偏置 $p$ 随时间漂移会破坏哪一步？GB-4-20：主持人策略未知时需要估计什么似然？GB-4-24：三发相邻时结论如何变化？GB-4-25：若抽样信息来自“随机告诉你一张 A”会怎样？GB-4-28：随机速度如何改变占用分布？}}

\providecommand{\MFQMappedUpperChNineQuestions}{
\subsection*{研究判断（来源映射替换）}
\begin{enumerate}[nosep,leftmargin=*]
\item \MFQInterviewMeta{研究判断}{12 分钟}{中}{次序统计量与假设边界}{GB-4-40 Expected maximum and minimum}若 $X_1,\ldots,X_n$ 独立均匀于 $[0,1]$，推导最大值与最小值的分布和期望；再说明样本非独立时为何不能直接复用。
\item \MFQInterviewMeta{笔试推导}{10 分钟}{中}{正态矩与递推}{GB-4-32 Normal moments}对 $Z\sim N(0,1)$，用分部积分证明 $E[Z^{2k}]=(2k-1)E[Z^{2k-2}]$，并解释奇数阶矩为何为零。
\item \MFQInterviewMeta{研究判断}{12 分钟}{难}{相关界与可达性}{GB-4-41 Correlation maximum and minimum}给定随机变量 $X$ 的边际分布，讨论 $\mathrm{Corr}(X,g(X))$ 的上下界；说明 Cauchy--Schwarz 的 $[-1,1]$ 只是普适界，何时才能取到等号。
\end{enumerate}}
\providecommand{\MFQMappedUpperChNineHints}{\par\textbf{来源映射题提示：}$P(M\le x)=x^n$；正态矩把 $z\phi(z)$ 写成 $-\phi'(z)$；相关取等要求中心化变量几乎处处线性相关。}
\providecommand{\MFQMappedUpperChNineSolutions}{
\section*{来源映射替换题}
\begin{enumerate}[leftmargin=*]
\item \textbf{GB-4-40。}$F_M(x)=x^n$，故密度 $nx^{n-1}$、$E[M]=n/(n+1)$。$P(m>x)=(1-x)^n$，所以 $E[m]=\int_0^1(1-x)^n dx=1/(n+1)$。仿射变换给 $a+(b-a)$ 乘相应结果；独立同分布是推导前提。
\item \textbf{GB-4-32。}标准正态密度满足 $\phi'(z)=-z\phi(z)$。分部积分得 $E[Z^{2k}]=(2k-1)E[Z^{2k-2}]$，从 $E[1]=1$ 递推为 $(2k-1)!!$；奇次幂乘偶密度是奇函数，积分为零。非中心正态应先标准化或用矩母函数。
\item \textbf{GB-4-41。}Cauchy--Schwarz 给 $|\mathrm{Corr}(X,g(X))|\le1$；等号当且仅当 $g(X)-Eg(X)=a(X-EX)$ 几乎处处。若 $g$ 被限制为非线性、单调或离散决策类，可达界通常更窄，必须在允许函数类与边际分布下重新优化，不能把普适界当成可实现值。
\end{enumerate}\MFQInterviewFollowup{GB-4-32：如何由矩母函数复核偶数矩？GB-4-40：样本有厚尾时，极值渐近要改用哪类条件？GB-4-41：限制 $g$ 为指标函数后如何求最优阈值？}}

\providecommand{\MFQMappedUpperChElevenQuestions}{
\subsection*{研究判断（来源映射替换）}
\begin{enumerate}[nosep,leftmargin=*]
\item \MFQInterviewMeta{研究判断}{15 分钟}{难}{模式等待与 Markov 状态}{GB-5-03 Coin triplets}反复抛公平硬币，比较模式 HHT 与 HTH 的首次出现等待时间；构造前缀状态，并审计“长度相同所以等待时间相同”的结论。
\item \MFQInterviewMeta{笔试推导}{12 分钟}{中}{吸收概率与差分方程}{GB-5-01 Gambler's ruin II}资本在 $0$ 与 $N$ 之间作步长 $\pm1$ 的随机游走，向上概率为 $p$。推导从 $i$ 出发先到 $N$ 的概率，并单独处理 $p=1/2$。
\item \MFQInterviewMeta{数值编程}{12 分钟}{中}{首达时间与状态截断}{GB-5-05 Drunk man}随机游走者从整数点 $i$ 出发，在 $0,N$ 吸收。建立期望吸收时间的线性方程并用三对角求解器计算；与 Monte Carlo 比较时报告置信误差。
\item \MFQInterviewMeta{研究判断}{12 分钟}{中}{排队状态与独立假设}{GB-5-07 Ticket line}售票队列中到达与服务均被建模为 Poisson/指数。写出出生--死亡状态，并说明只匹配均值为何不足以支持等待时间结论。
\item \MFQInterviewMeta{笔试推导}{12 分钟}{中}{动态规划与系列赛}{GB-5-11 World series}两队进行先赢四场的系列赛，每场 A 胜率固定为 $p$ 且独立。写出从比分 $(a,b)$ 出发 A 最终夺冠的递推和边界，并给出 $p=1/2$ 时开赛前概率。
\end{enumerate}}
\providecommand{\MFQMappedUpperChElevenHints}{\par\textbf{来源映射题提示：}模式状态取仍可延伸目标的最长后缀；破产概率解二阶差分；吸收时间右端多一个常数项；排队先核验 Markov/无记忆假设；系列赛以当前胜场为状态。}
\providecommand{\MFQMappedUpperChElevenSolutions}{
\section*{来源映射替换题}
\begin{enumerate}[leftmargin=*]
\item \textbf{GB-5-03。}对每个模式建立前缀自动机，状态为已匹配长度 $0,1,2,3$，吸收态为 3；对 $i<3$ 写 $E_i=1+(E_{T(i,H)}+E_{T(i,T)})/2$。求解得到 HHT 的期望等待 $8$，HTH 的期望等待 $10$。差异来自 HTH 失败后保留的前缀结构，不能用“长度都是三”断言相同。
\item \textbf{GB-5-01。}令 $h_i=P_i(\tau_N<\tau_0)$，则 $h_i=ph_{i+1}+(1-p)h_{i-1}$，$h_0=0,h_N=1$。若 $p\ne1/2$，$h_i=[1-((1-p)/p)^i]/[1-((1-p)/p)^N]$；公平时取极限得 $i/N$。边界和步长改变时不能套用闭式。
\item \textbf{GB-5-05。}期望吸收时间满足 $e_i=1+pe_{i+1}+(1-p)e_{i-1}$、$e_0=e_N=0$，形成三对角线性系统。公平时闭式为 $e_i=i(N-i)$。模拟均值应附标准误或区间；厚尾的停止时间会使有限样本收敛看起来很慢。
\item \textbf{GB-5-07。}状态可取系统人数 $n$，转移率为 $n\to n+1$ 的 $\lambda$ 与 $n\to n-1$ 的 $\mu$（$n>0$）。稳态要求 $\rho=\lambda/\mu<1$。若到达有聚集、服务时长非指数或服务能力随时段变化，均值相同也会产生不同等待尾部。
\item \textbf{GB-5-11。}令 $V(a,b)$ 为 A 夺冠概率，边界 $V(4,b)=1,V(a,4)=0$，递推 $V(a,b)=pV(a+1,b)+(1-p)V(a,b+1)$。$p=1/2$ 且开赛前双方对称，故 $V(0,0)=1/2$。若胜率随主客场或比分变化，状态必须扩充。
\end{enumerate}\MFQInterviewFollowup{GB-5-01：加入每步不同赌注后状态如何改变？GB-5-03：如何由前缀自动机直接生成转移矩阵并比较模式竞赛胜率？GB-5-05：如何避免无限循环模拟拖垮测试？GB-5-07：怎样从数据检验指数服务假设？GB-5-11：七局赛改成动态胜率时如何后向计算？}}

\providecommand{\MFQMappedUpperThirteenQuestions}{
\subsection*{研究判断（来源映射替换）}
\begin{enumerate}[nosep,leftmargin=*]
\item \MFQInterviewMeta{研究判断}{12 分钟}{中}{约束与乘子解释}{GB-3-10 Lagrange multipliers}最小化 $w^T\Sigma w$，约束 $\mathbf1^Tw=1$ 与 $\mu^Tw=r_0$。写出 KKT 线性系统，并审计“求解器成功即表示乘子和经济约束都可识别”的结论。
\end{enumerate}}
\providecommand{\MFQMappedUpperThirteenHints}{\par\textbf{来源映射题提示：}把两个等式约束矩阵记为 $A^Tw=b$，联立 $2\Sigma w+A\lambda=0$。}
\providecommand{\MFQMappedUpperThirteenSolutions}{
\section*{来源映射替换题}
\begin{enumerate}[leftmargin=*]
\item \textbf{GB-3-10。}KKT 系统为 $\left(\begin{smallmatrix}2\Sigma&A\\A^T&0\end{smallmatrix}\right)(w,\lambda)^T=(0,b)^T$，其中 $A=[\mathbf1,\mu]$、$b=(1,r_0)^T$。若 $A$ 列相关或与 $\Sigma$ 的零空间形成退化，约束资格或严格凸性失效；程序返回一个解不证明经济约束可识别。
\end{enumerate}\MFQInterviewFollowup{GB-3-10：加入不等式限制后哪些 KKT 条件需要新增？}}

\providecommand{\MFQMappedUpperFourteenQuestions}{
\subsection*{数值编程与研究判断（来源映射替换）}
\begin{enumerate}[nosep,leftmargin=*]
\item \MFQInterviewMeta{数值编程}{12 分钟}{中}{Newton 收敛域}{GB-3-09 Newton's method}用 Newton 法求 $x^3-2x-2=0$，记录不同初值的迭代；给出停止准则，并构造一个进入循环或远离根的初值。
\item \MFQInterviewMeta{笔试推导}{12 分钟}{中}{有限差分误差}{GB-7-15 Finite difference}从 Taylor 展开推导中心差分的一阶导数公式与 $O(h^2)$ 截断误差，再解释 $h$ 过小时舍入误差为何上升。
\item \MFQInterviewMeta{数值编程}{15 分钟}{中}{Monte Carlo 误差预算}{GB-7-14 Monte Carlo}用 Monte Carlo 估计 $E[e^Z]$（$Z\sim N(0,1)$），给出独立理论值、标准误和置信区间；再用 antithetic variates 比较方差而不是只比较单次误差。
\end{enumerate}}
\providecommand{\MFQMappedUpperFourteenHints}{\par\textbf{来源映射题提示：}Newton 需同时检查残差与步长；中心差分把两侧 Taylor 展开相减；Monte Carlo 先用矩母函数给独立 oracle，再比较配对样本的方差。}
\providecommand{\MFQMappedUpperFourteenSolutions}{
\section*{来源映射替换题}
\begin{enumerate}[leftmargin=*]
\item \textbf{GB-3-09。}$x_{k+1}=x_k-(x_k^3-2x_k-2)/(3x_k^2-2)$。停止应同时要求 $|f(x_k)|$ 和相对步长足够小，并设最大迭代；导数接近零或初值落在非吸引区域时可能跳远/循环。面试追问不应以某个库的 success 标志替代理论收敛域。
\item \textbf{GB-7-15。}两式相减得 $[f(x+h)-f(x-h)]/(2h)=f'(x)+h^2f^{(3)}(\xi)/6$。截断误差随 $h^2$ 降低，但两近数相减的浮点绝对误差约为 $\epsilon|f|$，除以 $h$ 后放大为 $O(\epsilon/h)$，因此存在最优中间尺度。
\item \textbf{GB-7-14。}正态矩母函数给 $E[e^Z]=e^{1/2}$，且 $\mathrm{Var}(e^Z)=e^2-e$。独立样本均值的标准误为样本标准差除以 $\sqrt n$。反变量配对使用 $[e^Z+e^{-Z}]/2$；应在相同随机预算下比较配对估计量的样本方差和区间宽度，而不是挑一次更接近真值的结果。
\end{enumerate}\MFQInterviewFollowup{GB-3-09：如何用阻尼或信赖域避免跳远？GB-7-14：控制变量与重要性抽样分别需要什么额外结构？GB-7-15：如何从截断与舍入两项估计最佳 $h$？}}

\providecommand{\MFQMappedDerivativesStochasticQuestions}{
\section*{研究判断（来源映射替换）}
\begin{enumerate}[leftmargin=*]
\item \MFQInterviewMeta{研究判断}{15 分钟}{难}{停止时与反射原理}{GB-5-15 Stopping and first passage}对标准 Brownian 运动推导首次越过 $a>0$ 在 $T$ 前发生的概率，并审计直接将无界停止时代入鞅等式的证明。
\end{enumerate}}
\providecommand{\MFQMappedDerivativesStochasticHints}{\par\textbf{来源映射题提示：}用反射原理把越界后终点低于 $a$ 的路径映到终点高于 $a$。}
\providecommand{\MFQMappedDerivativesStochasticSolutions}{\section*{来源映射替换题}\begin{enumerate}[leftmargin=*]\item \textbf{GB-5-15。}反射原理给 $P(\tau_a\le T)=2P(W_T\ge a)=2[1-\Phi(a/\sqrt T)]$。若改用可选停止证明，必须先对停止时截断并验证一致可积或相应有界条件；直接把无界停止时代入鞅等式是常见错误。\end{enumerate}\MFQInterviewFollowup{GB-5-15：若边界改为随时间变化的 $a(t)$，反射证明的哪一步不再直接成立？}}

\providecommand{\MFQMappedDerivativesNumericsQuestions}{
\section*{研究判断（来源映射替换）}
\begin{enumerate}[leftmargin=*]
\item \MFQInterviewMeta{研究判断}{10 分钟}{中}{无套利复制}{GB-6-02 Put-call parity}在含连续股息率 $q$ 的市场推导欧式看涨--看跌平价，并审计一个忽略 bid/ask、融资、借券与结算时点便声称可套利的报告。
\end{enumerate}}
\providecommand{\MFQMappedDerivativesNumericsHints}{\par\textbf{来源映射题提示：}比较 $C+Ke^{-rT}$ 与 $P+S_0e^{-qT}$ 的终值。}
\providecommand{\MFQMappedDerivativesNumericsSolutions}{\section*{来源映射替换题}\begin{enumerate}[leftmargin=*]\item \textbf{GB-6-02。}两组合到期均为 $\max(S_T,K)$，故 $C-P=S_0e^{-qT}-Ke^{-rT}$。若左侧过高则卖出昂贵组合、买入便宜组合；真实可执行性还受 bid/ask、融资、股息预测、借券、行权与结算时点约束。\end{enumerate}\MFQInterviewFollowup{GB-6-02：美式期权或随机股息下，复制组合的哪个等式需要重写？}}

\providecommand{\MFQMappedPortfolioOptimizationQuestions}{
\section*{来源映射替换题}
\begin{enumerate}[leftmargin=*]
\item \MFQInterviewMeta{笔试推导}{15 分钟}{难}{均值方差与估计误差}{GB-6-14 Portfolio optimization}推导预算约束下的最小方差权重；随后说明均值方差解为何会放大预期收益估计误差，并提出至少两种稳健化方案。
\end{enumerate}}
\providecommand{\MFQMappedPortfolioOptimizationHints}{\par\textbf{来源映射题提示：}先解 $\min w^T\Sigma w$、$\mathbf1^Tw=1$；含均值时观察 $\Sigma^{-1}\hat\mu$。}
\providecommand{\MFQMappedPortfolioOptimizationSolutions}{\section*{来源映射替换题}\begin{enumerate}[leftmargin=*]\item \textbf{GB-6-14。}$w^*=\Sigma^{-1}\mathbf1/(\mathbf1^T\Sigma^{-1}\mathbf1)$。含预期收益的解依赖 $\Sigma^{-1}\hat\mu$，小特征方向和噪声均值会被放大。可用协方差收缩、均值不确定集、权重/换手约束或 Bayesian shrinkage，并必须重新比较样本外风险和成本。\end{enumerate}\MFQInterviewFollowup{GB-6-14：禁止卖空后闭式解为何失效，应怎样识别活跃约束？}}

\providecommand{\MFQMappedPortfolioTailQuestions}{
\section*{来源映射替换题}
\begin{enumerate}[leftmargin=*]
\item \MFQInterviewMeta{口述概念}{10 分钟}{中}{VaR 的性质与反例}{GB-6-15 Value at Risk}解释历史 VaR、参数 VaR 与 ES 的估计对象；构造一个 VaR 不次可加的简单离散损失例，并说明这对风险限额意味着什么。
\end{enumerate}}
\providecommand{\MFQMappedPortfolioTailHints}{\par\textbf{来源映射题提示：}让两个头寸各自以小概率损失、且损失事件互斥；单项 VaR 可为零而组合 VaR 为正。}
\providecommand{\MFQMappedPortfolioTailSolutions}{\section*{来源映射替换题}\begin{enumerate}[leftmargin=*]\item \textbf{GB-6-15。}历史法取经验损失分位数，参数法在分布假设下求分位数，ES 是越过分位点的尾部平均。令两个头寸各以 $4\%$ 概率损失 1，事件互斥；在 $95\%$ 水平各自 VaR 为 0，而组合有 $8\%$ 概率损失 1，VaR 为 1，故不次可加。限额系统不能把分散收益当作由 VaR 自动保证。\end{enumerate}\MFQInterviewFollowup{GB-6-15：把置信度提到 $99\%$ 时，需要多少尾部观测才能支撑可辩护的 ES？}}

\MFQMappedDerivativesStochasticQuestions


\chapter{定价、校准与对冲：从数值方法到误差分布}
\label{ch:lower-derivatives-numerics}

\begin{itemize}
  % Generated from curriculum/manifest.json; do not edit by hand.
\item 直接先修：下册《随机分析与无套利：从二次变差到风险中性测度》。

\end{itemize}

复制得到价格方程之后，还要把它算出来。以下用一层树说明风险中性概率的来源，再把空间导数变成差分系数，最后讨论报价反演的敏感度。


\MFQTransition{所求解的定价问题}
上一章已由自融资复制推导 Black--Scholes PDE。本节沿用同一方程与终端支付；数值方法改变的是求解方式，而非另行假定价格期望。含连续股息率 $q$ 时，股票方向的一阶系数为 $(r-q)S$。树法、差分与模拟必须针对同一合约比较。

\section{同一价格，不同误差来源}

欧式看涨闭式价格为
\[
 C=S_0\Phi(d_1)-Ke^{-rT}\Phi(d_2),\qquad
 d_{1,2}=\frac{\log(S_0/K)+(r\pm\sigma^2/2)T}{\sigma\sqrt T}.
\]
在 $S_0=K=100,r=2\%,\sigma=20\%,T=1$ 下，解析值是 $8.916037$。它是当前模型中的参考值，不是真实市场的“正确价格”。

% Generated by tools/build_figures.py; edit figures/figure-specs.json instead.
\MFQFigure{tex/figures/generated/lower-dv-methods.pdf}{不同定价方法使用相同的合约与模型输入，却引入不同的近似误差；加密网格、增加样本及比较方法可以帮助理解这些误差。}{lower-dv-methods}


这些方法近似同一个模型价格，但误差来源不同：
\begin{center}
\begin{tabular}{lll}
\toprule
方法 & 首要误差 & 误差的考察方式\\
\midrule
闭式 & 参数/模型错设 & 边界、单位与极限情形\\
CRR 树 & 时间离散 & $p\in[0,1]$ 与步数敏感性\\
隐式 PDE & 空间、时间、截断 & 边界条件与三维网格敏感性\\
Monte Carlo & 采样与方差 & 独立种子、标准误与区间\\
\bottomrule
\end{tabular}
\end{center}

\MFQTransition{例：一层树中的复制}

取 $S_0=100,u=1.2,d=0.8,K=100$，一期无风险总回报为 $R=1.02$。终端看涨支付为 $20,0$。复制要求 $120\Delta+RB=20$、$80\Delta+RB=0$，相减得 $\Delta=1/2$、$B=-40/R$，初始价格为 $50-40/R\approx10.7843$。风险中性概率 $p=(R-d)/(u-d)=0.55$ 给同一个价格 $(0.55\cdot20)/R$。因为 $d<R<u$，两状态概率均严格为正。

CRR 树取 $u=e^{\sigma\sqrt{\Delta t}},d=u^{-1}$ 和
\[
 p=\frac{e^{r\Delta t}-d}{u-d}.
\]
需要先验证 $p\in[0,1]$，再从终端收益反推。256 步价格为 $8.908302$，其误差主要来自时间离散\cite{cox1979option}。

\MFQTransition{从空间导数到三对角系数}

令剩余期限为 $\tau$、$S_j=j\Delta S$，无股息空间算子为 $LV=\tfrac12\sigma^2S^2V_{SS}+rSV_S-rV$。中心差分给
\[
 (L_hV)_j=A_jV_{j-1}+B_jV_j+C_jV_{j+1},
\]
\[
 A_j=\tfrac12\sigma^2j^2-\tfrac12rj,\quad
 B_j=-\sigma^2j^2-r,\quad
 C_j=\tfrac12\sigma^2j^2+\tfrac12rj.
\]
从终端支付 $V_j^0=(S_j-K)_+$ 向增加的 $\tau$ 推进，隐式一步解 $(I-\Delta\tau L_h)V^{n+1}=V^n$；已知边界项移到右端。notebook 会显示一次树的倒推，再比较多层树与闭式价格，并用二分反演隐波。

隐式有限差分在 $S\in[0,S_{\max}]$ 上离散 PDE。每个时间层解三对角线性系统
\[
 (I-\Delta t L_h)V^{n+1}=V^n,
\]
并使用 $V(\tau,0)=0$、$V(\tau,S_{\max})\approx S_{\max}-Ke^{-r\tau}$。事先确定网格给出 $8.908382$。误差同时包含空间离散、时间离散和右边界截断；只加密时间而固定 $S_{\max}$ 不能消除全部误差。

教学实验保持其他维度不变，分别把空间节点减半、时间节点减半，并在保持 $\Delta S$ 近似不变时把右边界从 $400$ 缩到 $300$。相对基准网格的差分别为 $0.020505$、$0.000857$ 与约 $4.4\times10^{-14}$。这些是敏感性诊断，不是严格误差上界；边界差很小只说明本合约与本参数下 $S_{\max}=300$ 已足够远。

\MFQTransition{差分格式与误差来源}

把 PDE 离散为显式、隐式或 Crank--Nicolson 格式。显式法便于理解但受稳定步长约束；隐式法每步解三对角系统，稳定性更好；Crank--Nicolson 二阶时间精度但不光滑 payoff 附近可能振荡，可先做 Rannacher 平滑。

% Generated by tools/build_figures.py; edit figures/figure-specs.json instead.
\MFQFigure{tex/figures/generated/lower-dv-pde.pdf}{有限差分从到期支付出发，沿剩余期限增加的方向求解价格；空间边界和时间步长共同影响数值近似。}{lower-dv-pde}


误差预算至少分为空间截断、时间截断、$S_{\max}$ 边界和线性求解误差。只把网格加密一次不能识别哪一项主导。对无套利价格还应检查非负、随 $S$ 单调、对 $K$ 凸等离散性质。

风险中性 Monte Carlo 使用
\[
 S_T=S_0\exp\left((r-\tfrac12\sigma^2)T+\sigma\sqrt T Z\right)
\]
和贴现收益均值，其中 Brownian 过程记为 $W_t$，且 $Z=W_T/\sqrt T\sim N(0,1)$。20,000 条路径给出 $8.779633$，标准误 $0.096831$；解析值落在约两个标准误内。随机误差按 $N^{-1/2}$ 收敛，方差缩减、准 Monte Carlo 或多层方法改变的是计算效率，不修复模型错设\cite{glasserman2004montecarlo}。

树的步数、PDE 的时间和空间网格、Monte Carlo 的路径数，分别影响不同类型的近似误差。比较价格时，需要同时比较相应的误差大小。解析值偶尔落在模拟置信区间之外，也可能只是区间本身未覆盖真值。

同一计算还可以用以下方法比较：
\begin{itemize}
  \item 闭式用 SciPy \texttt{ndtr}，树用 SciPy 二项分布概率质量；
  \item PDE 用 \texttt{solve\_banded}；Monte Carlo 的独立参照用 SciPy 自适应积分直接积分正态密度下的贴现收益。
\end{itemize}
这些方法使用相同的合约和参数，却采用不同的数值计算方式。确定性方法之间的差异主要来自舍入与离散化；Monte Carlo 与数值积分的差异还包含随机抽样误差，应结合标准误解释。

\section{先反演单点，再校准曲面}

给定单个报价 $C^{\mathrm{mkt}}$，隐含波动率解
\[
 C^{BS}(S,K,r,T,\sigma_{\mathrm{impl}})=C^{\mathrm{mkt}}.
\]
这是一维根求解。每个节点各自找到一个 $\sigma_j$，目标损失接近零是必然结果，不能称为“曲面模型校准”。根必须被静态套利上下界 bracket；Vega 太小时还应报告病态性。

% Generated by tools/build_figures.py; edit figures/figure-specs.json instead.
\MFQFigure{tex/figures/generated/lower-dv-vol-surface.pdf}{隐含波动率同时随执行价和期限变化；切片形状必须与报价单位及无套利条件一起解释。}{lower-dv-vol-surface}


二分法的优点是只依赖单调性和 bracket；Newton 法利用 Vega 收敛更快，但在深度实值、深度虚值或短期限节点可能跳出合法区间。可靠实现可先用套利界拒绝不可能报价，再用带 bracket 的方法求根，同时保存价格残差和 Vega。bid/ask 区间内存在许多可接受隐波时，返回过多小数位只会制造虚假精度。

第二阶段才拟合参数化总方差。事先确定教学模型令 $k=\log(K/F_T)$，
\[
 w(k,T)=\sigma^2(k,T)T=aT+bk^2+cT^2.
\]
对 15 个非等距执行价与期限节点先逐点反演，再以加权最小二乘恢复 $(a,b,c)=(0.035,0.080,0.006)$。这个例子说明报价怎样决定一组共享参数，以及拟合残差怎样反映参数化限制。该形式本身不保证一般的曲面无套利；更深入的研究可参见局部波动率等模型\cite{dupire1994pricing}。

校准试图由多个报价确定一组共享参数。设计矩阵接近奇异时，不同参数可能产生近似相同的样本内价格，却对缺失期限给出不同外推。可以比较参数约束、报价权重和留出节点的重定价误差。若权重来自买卖价差所反映的报价精度，应先定义这一关系；根据拟合残差事后降低某些权重，会改变原先的估计问题。

\MFQTransition{非等距凸性与期限条件}

看涨价格随执行价递减。非等距网格的凸性应比较斜率
\[
 \frac{C_i-C_{i-1}}{K_i-K_{i-1}}
 \le
 \frac{C_{i+1}-C_i}{K_{i+1}-K_i},
\]
不能直接用 $C_{i-1}-2C_i+C_{i+1}\ge0$。后者只在等距网格对应离散二阶差分。

固定执行价的看涨价格随期限不降，需要无分红与非负利率等条件。存在分红、负利率或不同远期时，应先明确贴现因子与远期归一化，再选择适用的日历无套利条件；期限之间的比较因此依赖所采用的贴现、远期与执行价约定。

归一化模式令 $F_T=S_0e^{(r-q)T}$、$D_T=e^{-rT}$，在固定
$k=K/F_T$ 上比较 $C/(D_TF_T)$。教学实验取非零股息率 $q=3\%$，由每个节点携带
\texttt{forward} 与 \texttt{discount\_factor}；网格不完整或原始固定执行价模式与假设冲突时均拒绝。这样“通过日历检查”同时携带远期、贴现和 moneyness 口径。

\MFQTransition{从公开利率到贴现输入}

合成报价节点用于理解参数反演，公开利率例子则演示如何把决策日前可得的利率转换为贴现因子；国债 par yield 也不等于期权定价所需的完整零息曲线。该快照不能验证隐波、校准或交易收益。

从 par yield 到零息贴现因子通常还需 bootstrap，并明确日计数、复利方式和结算日。本例直接使用连续复利，便于理解贴现计算；具体合约仍需要相应的曲线、计息和结算约定。把 1 年 par yield 直接塞进所有期限的 Black--Scholes 公式，会同时制造期限错配与插值偏差。

定价方法近似当前模型下的价格，校准则从报价反推参数。接下来考察这些近似怎样随网格和输入变化，并进一步联系到离散对冲误差。

\MFQTransition{抽样与反演的联系}
欧式看涨的 Monte Carlo 可用 $S_T$ 作控制变量，其风险中性期望为 $S_0e^{(r-q)T}$；路径依赖支付则还要分析路径离散误差。隐波反演把价格误差放大约为 $\Delta\sigma\approx\Delta C/\mathrm{Vega}$，因此 Vega 很小时，价格数值误差与 bid/ask 都可能造成较大的隐波不确定性。

\begin{diagnostic}
校准误差接近零可能只是每个报价都有一个自由参数。模型风险要看参数稳定、插值外推、无套利和对未参与校准报价的表现，而不是训练目标是否为零。
\end{diagnostic}

定价与校准结果会在衍生品路线综合项目中同对冲误差、成本和压力情景合并评价；单点拟合并不等于可实施模型。

\section{离散对冲与模型失效}
\label{ch:lower-derivatives-hedging}

\begin{itemize}
  % Generated from curriculum/manifest.json; do not edit by hand.
\item 直接先修：下册《定价与校准：从闭式、树和 PDE 到参数化曲面》。

\end{itemize}

连续复制允许不断调整 Delta，实际交易只能在有限时点成交。下面把股票与现金逐笔相加，说明对冲误差如何产生，也说明手续费为什么会积累。


\MFQLead{Greeks 是局部敏感度，不是损益保证}

在 Black--Scholes 模型中
\[
 \Delta=\frac{\partial C}{\partial S}=\Phi(d_1),\qquad
 \Gamma=\frac{\partial^2C}{\partial S^2},\qquad
 \mathrm{Vega}=\frac{\partial C}{\partial\sigma}.
\]
解析 Greek 应与中心差分在一组逐步缩小但不过度受舍入支配的步长上比较。差分一致只检查导数实现；若价格模型漏掉股息、跳跃或曲面动态，Greek 仍可能在经济上错误。

% Generated by tools/build_figures.py; edit figures/figure-specs.json instead.
\MFQFigure{tex/figures/generated/lower-dv-greeks.pdf}{Delta、Gamma 与 Vega 描述不同局部扰动方向；同一头寸的风险会随标的价格变化而重新分配。}{lower-dv-greeks}


教学实验以现货步长 $1,0.5,0.25,0.125$ 检查 Delta 与 Gamma，并按比例缩小波动率步长检查 Vega。由最粗到最细步长，三者绝对差分别从
$6.51\times10^{-5}$、$3.10\times10^{-6}$、$6.53\times10^{-4}$ 收敛到
$1.02\times10^{-6}$、$4.84\times10^{-8}$、$1.02\times10^{-5}$；notebook 同时保存并绘制全部四个步长。步长继续缩小时舍入误差可能反弹，因此不以单个巧合步长充当证明。

小变动的 P\&L 展开为
\[
 \Delta C\approx \Delta\,\Delta S+\tfrac12\Gamma(\Delta S)^2
 +\mathrm{Vega}\,\Delta\sigma+\Theta\,\Delta t.
\]
Delta 中性只消除一阶现货项。Gamma、波动率、时间衰减、离散调仓和交易成本仍进入损益。

若实际实现波动率为 $\sigma_{\mathrm{real}}$，而期权按 $\sigma_{\mathrm{impl}}$ 定价，对于持有期权并做空 Delta 股票的方向，连续小步近似下 P\&L 的核心项与
\[
 \frac12\Gamma S^2\left(\sigma_{\mathrm{real}}^2-\sigma_{\mathrm{impl}}^2\right)\Delta t
\]
相关。它解释了波动率观点如何进入对冲损益，但仍忽略跳跃、曲面重标定、交易成本和 Greek 离散误差，不能当作逐路径恒等式。

\MFQLead{自融资离散台账}

路径仍由风险中性 Brownian 运动 $W_t$ 驱动；实际对冲评估则必须说明使用哪个测度和真实世界路径模型。

期初卖出一份期权，买入 $\Delta_0$ 股，其余进入现金账户。每个调仓时点先让现金按无风险利率增长，再用当前价格买卖 $\Delta_{t_i}-\Delta_{t_{i-1}}$ 股，并扣除
\[
 \mathrm{cost}_{t_i}=c\,S_{t_i}\left|\Delta_{t_i}-\Delta_{t_{i-1}}\right|.
\]
到期组合价值减支付收益得到复制误差。必须保存旧 Delta、交易量、现货价格、现金、成本和到期 payoff；只保存最终误差无法定位符号、融资或时点错误。

以一步路径为手算基准时，期初成本为 $c|\Delta_0|S_0$，到期前不再调仓，故成本拖累应等于该成本按现金账户增长后的值。仓库保留这条独立手算测试，用来捕获“成本没有融资”“到期重复调仓”或“成本符号加反”等错误。多路径实验是在这个小计算记录通过以后才有意义。

透明实现逐路径执行该计算记录，第二个独立向量实现以 SciPy \texttt{ndtr} 计算批量 Delta，并用 NumPy 同时更新 512 条路径。两者最大逐路径误差差约为 $5.7\times10^{-14}$；它们共享模型假设，不能互相证明 GBM 合理。

\MFQLead{从一条路径升级为分布}

单路径误差可能恰好很小或很大。固定 512 条路径、24 个时间步后，至少报告
\[
 \mathrm{bias}=N^{-1}\sum_j e_j,\qquad
 \mathrm{RMSE}=\sqrt{N^{-1}\sum_j e_j^2},
\]
以及 5\%、50\%、95\% 分位数。教学实验无成本 bias/RMSE 为 $(-0.086348,1.416360)$；加入比例成本后为 $(-0.192103,1.434142)$，误差分位数为 $(-2.485457,-0.052646,1.960850)$。平均名义成本为 $0.104415$，95\% 分位为 $0.149720$。

% Generated by tools/build_figures.py; edit figures/figure-specs.json instead.
\MFQFigure{tex/figures/generated/lower-dv-hedge-error.pdf}{离散对冲误差在不同路径上形成一个分布；在本例模型中，提高调仓频率减少离散误差，却可能增加交易成本。}{lower-dv-hedge-error}


这些数值只描述固定随机种子和参数。正式实验还应改变调仓频率、成本、跳跃强度、实现波动率、隐含波动率与曲面动态，报告误差的抽样不确定性。提高调仓频率通常降低离散误差，却可能增加成本；不存在脱离摩擦的“越频繁越好”。

同一规则在 12、24、52 次调仓下的成本后 RMSE 为 $2.040789$、$1.434142$、$0.919696$，平均名义成本为 $0.081246$、$0.104415$、$0.139579$。该组实验同时展示离散误差下降与成本上升，不能从中挑一个维度宣称“更优”。

bias 与 RMSE 本身也有抽样误差。可以对路径级误差 bootstrap，或用多组独立 seed 重复整个实验并报告区间。若路径共享市场因子或波动率状态，简单把每条路径当独立观察会低估不确定性；分组重采样单位应与共同冲击结构一致。

\MFQTransition{跳跃与离散调仓}
跳跃、随机波动率和延迟使实际路径偏离连续复制条件。对固定策略比较不同调仓频率与成本率，应同时观察误差分布和总交易成本；路径多只能减小给定模型内的抽样误差。

\MFQTransition{Greek 的单位与组合聚合}

Delta 是价格对标的的一阶导，Gamma 是二阶曲率，Vega 是对波动率“绝对变化 1”还是“一个波动率点”的敏感度必须注明。Theta 的日历日/交易日单位、Rho 的利率单位也常被混淆。组合 Greek 先按合约乘数、数量和币种换算后再相加。

有限差分检查采用中心差分
\[
\Delta\approx\frac{V(S+h)-V(S-h)}{2h},\qquad
\Gamma\approx\frac{V(S+h)-2V(S)+V(S-h)}{h^2}.
\]
$h$ 太大截断误差高，太小则浮点消减严重；应画 $h$ 的误差曲线，而不是只挑一个能对上的步长。

\MFQTransition{离散 Delta 对冲误差的来源}

持有期权、做空 Delta 股票的损益可局部近似分解为
\[
\Delta\Pi\approx
\frac12\Gamma S^2(\sigma_{\rm realized}^2-\sigma_{\rm hedge}^2)\Delta t
-\text{cost}+\text{jump/model residual}.
\]
后文卖出期权并持有复制组合的误差方向相反。这是局部近似，揭示 Gamma 暴露如何把实现方差与对冲波动率差转成损益。更频繁调仓降低离散误差，却增加点差和冲击，因此存在成本依赖的最优频率，而不是“越频繁越好”。

跳跃时 Delta 无法提前复制不连续变动；随机波动率引入 Vega 风险；曲面移动还涉及 sticky-strike、sticky-delta 等约定。对冲报告应按这些情景分组，而非只给总体 RMSE。

\MFQTransition{从模型价格到市场报价}

做市价格还要覆盖库存、资本、模型不确定性和流动性。市场 bid/ask 不是同一个“真波动率”的噪声观察；用中价校准后按中价计算对冲收益会忽略入场和退出价差。

曲面随时间更新时，参数跳动可能来自数据噪声、报价稀疏或优化局部极值。应记录报价过滤、权重、初值、约束、收敛状态和参数时间序列。无法校准时应保留失败状态，不用昨日参数静默填充。

\MFQTransition{对冲实验的统计设计}

多路径实验至少报告 bias、RMSE、中位数、尾部分位、成本分布和最大资金需求。比较两个对冲策略时使用共同随机数，让路径噪声配对抵消。模型和策略调参使用开发情景，最终压力情景只打开一次。

真实历史回放中只有一条市场路径，可按滚动起点形成多个重叠对冲 episode，但标准误要处理重叠。历史回放不能评价订单对市场的反事实影响，成交模型仍需单独压力测试。

\begin{diagnostic}
对冲误差小不等于定价模型正确：高频调仓和低成本假设可以掩盖模型错设。必须同时报告交易次数、成本、未成交、资本占用和压力期表现。
\end{diagnostic}

\section{本路线的综合项目}

以同一合约和参数贯通复制、定价、校准与离散对冲。路线 Capstone 的综合结果应给出价格与 Greeks 的独立核对、隐波或曲面的适用范围，以及多路径对冲误差、成本和压力情景分布；单条漂亮路径不能替代这些信息。

\subsection{两次持仓与现金更新}

设卖出期权收到价格 10，$S_0=100$，初始 Delta 为 $0.6$，利率暂为零、成本率为 $c=0.001$。买股花费 60，初始费用 0.06，现金为 $B_0=10-60-0.06=-50.06$。下一时点价格 110，新的 Delta 为 $0.8$；加买 0.2 股花费 22，费用 0.022，所以 $B_1=-72.082$。

到期 $S_T=105,K=100$，若将股票按中价估值而不另收平仓费，复制组合价值为 $0.8(105)-72.082=11.918$，支付为 5，误差为 $6.918$。这些 Delta 是教学输入，并非声称与某一波动率模型匹配。该例完整说明现金流符号，实际模型 Delta 会在 notebook 中由价格公式计算。

一般递推为 $B_i=B_{i-1}e^{r\Delta t_i}-S_i(\Delta_i-\Delta_{i-1})-cS_i|\Delta_i-\Delta_{i-1}|$。到期不再重新估计一个用于未来持有的 Delta，而是计算 $B_{m-1}e^{r\Delta t_m}+\Delta_{m-1}S_T-g(S_T)$；若要求实际平仓，再单独扣平仓费用。notebook 比较多条路径和不同调仓频率，保留成本前后误差的配对关系。

\section{分层练习}
\begin{enumerate}
  \item \textbf{口述概念：}树、PDE 和 Monte Carlo 各自主要有哪些误差？
  \item \textbf{口述概念：}为什么逐点隐波反演不是参数化曲面校准？
  \item \textbf{笔试推导：}写出隐式 PDE 三对角系统的三个系数。
  \item \textbf{笔试推导：}推导非等距执行价网格的斜率凸性条件。
  \item \textbf{数值编程：}分别加密空间、时间和截断边界，制作 PDE 误差表。
  \item \textbf{数值编程：}把参数化总方差系数恢复与价格残差分开报告。
  \item \textbf{研究判断：}何时固定执行价的日历单调性不宜直接使用？
  \item \textbf{研究判断：}财政部 par yield 快照还缺哪些输入才能成为期权贴现曲线？
\end{enumerate}

\MFQLead{Greeks、离散对冲与误差分布题组}
\begin{enumerate}
  \item \textbf{口述概念：}Delta 中性以后还剩哪些风险？
  \item \textbf{口述概念：}为什么单路径复制误差不能评价对冲质量？
  \item \textbf{笔试推导：}写出包含融资与比例成本的离散自融资递推。
  \item \textbf{笔试推导：}说明 bias、RMSE 和分位数回答的不同问题。
  \item \textbf{数值编程：}把调仓步数改为 12、24、52，并比较误差与成本。
  \item \textbf{数值编程：}加入一处价格跳跃，比较成本前后误差分布。
  \item \textbf{研究判断：}审计只展示一条“完美复制”路径的报告。
  \item \textbf{开放 Capstone：}结合本章模型，解释预测、成交和成本如何共同决定结果，并说明所用假设。
\end{enumerate}

\subsection*{基础衔接：独立计算}
在上述例子中增加到期平仓费用，误差减少多少？


\chapter{投资组合与风险管理：从协方差估计到压力测试}
\label{ch:lower-portfolio-estimation}

\begin{itemize}
  % Generated from curriculum/manifest.json; do not edit by hand.
\item 直接先修：上册第 6 章《线性代数、分解、投影与条件数》、上册第 10 章《估计、检验、渐近、回归与多重检验》、上册第 14 章《浮点数、数值线代与病态问题》、上册第 16 章《研究有效性、样本外与交易摩擦》、上册第 17 章《可复现研究审计》。

  \item \textbf{学习产出：}构造并诊断风险模型，完成风险预算、Black--Litterman、CVaR 与稳健组合优化，再用尾部统计、非线性重估和反向压力寻找失败阈值。
  \item \textbf{研究边界：}风险矩阵、收益观点和尾部分布都是估计对象；优化器与风险数字不会修复错误数据、制度变化或不可成交权重。
\end{itemize}

收益共同波动决定组合风险。先把收益矩阵中心化，再看为什么协方差必半正定、何时必奇异，才能理解收缩为何改善某些方向的稳定性。


\section{先确定收益矩阵的含义}

令 $R\in\mathbb R^{T\times n}$ 的第 $t$ 行是同一可用时点的资产收益，列顺序在研究期内固定。样本协方差为
\[
 \widehat\Sigma=\frac{1}{T-1}\sum_{t=1}^{T}(r_t-\bar r)(r_t-\bar r)^\top.
\]
该公式看似简单，却依赖收益口径、时区、公司行动、缺失值和资产池协议。成对删除会让不同矩阵元素来自不同日期集合，拼出的矩阵甚至可能不再半正定。

% Generated by tools/build_figures.py; edit figures/figure-specs.json instead.
\MFQFigure{tex/figures/generated/lower-pf-risk-map.pdf}{协方差矩阵把权重方向映射为风险椭圆；主轴方向和长度分别对应特征向量与特征值。}{lower-pf-risk-map}


任意权重 $w$ 都满足
\[
 w^\top\Sigma w=\operatorname{Var}(w^\top r)\ge0,
\]
因此半正定不是数值偏好，而是协方差定义的必要条件。可逆也不等于可靠：若最小特征值很小，微小估计误差会被逆矩阵放大。报告至少应保存样本窗口、维度、最小特征值、条件数和缺失值策略。

\section{收缩与因子风险模型}

线性收缩写成
\[
 \widehat\Sigma_\lambda=(1-\lambda)\widehat\Sigma+\lambda T,
 \qquad 0\le\lambda\le1.
\]
目标 $T$ 可以是对角矩阵、等相关矩阵或结构模型。收缩用偏差交换方差，目的在于改善样本外稳定性，而不是让历史拟合更漂亮。教学实验以平均方差乘单位阵为目标，$\lambda=0.25$；透明公式与 sklearn \texttt{ShrunkCovariance} 完全一致，最小特征值为 $3.58\times10^{-5}$\cite{ledoit2004well}。

% Generated by tools/build_figures.py; edit figures/figure-specs.json instead.
\MFQFigure{tex/figures/generated/lower-pf-shrinkage.pdf}{收缩在高方差样本协方差与结构化目标之间权衡，通常以少量偏差换取更低估计波动。}{lower-pf-shrinkage}


因子模型把收益写成
\[
 r=Bf+\varepsilon,\qquad
 \Sigma=B\Sigma_fB^\top+D,
\]
其中需要 $\operatorname{Cov}(f,\varepsilon)=0$，且 $D$ 的对角元素非负。它减少需要估计的自由参数，也引入因子遗漏、暴露漂移与特异风险相关性等错设。事先确定三资产、两因子矩阵的迹为 $0.1013$；该值可从 $B,\Sigma_f,D$ 独立重建。

\begin{diagnostic}
协方差收缩和因子模型回答“如何约束估计误差”，并不回答“未来制度是否与估计窗口相同”。应把结构选择、强度选择和窗口选择写进样本外协议。
\end{diagnostic}

\section{边际风险、成分风险与风险平价}

组合方差 $v(w)=w^\top\Sigma w$ 的梯度为 $2\Sigma w$。采用方差口径时，第 $i$ 个资产的成分贡献定义为
\[
 \mathrm{RC}_i=w_i(\Sigma w)_i,
 \qquad \sum_i\mathrm{RC}_i=w^\top\Sigma w.
\]
若使用波动率口径，则还需除以 $\sqrt{w^\top\Sigma w}$；两种口径不能混写。等风险预算要求
\[
 \mathrm{RC}_1=\cdots=\mathrm{RC}_n.
\]

对角协方差的波动率为 $(0.2,0.3,0.4)$ 时，权重与逆波动率成正比，因此
\[
 w=\frac{(5,10/3,5/2)}{5+10/3+5/2}
   =\left(\frac6{13},\frac4{13},\frac3{13}\right).
\]
三个方差贡献恰好相等。仓库的透明实现求解严格凸问题
\[
 \min_{y_i>0}\left\{\frac12y^\top\Sigma y-\sum_i\log y_i\right\},
\]
其一阶条件 $y_i(\Sigma y)_i=1$ 在存在负相关时仍成立；Newton 法求出 $y$ 后再归一化。独立 SLSQP 直接最小化贡献差，最大权重差约 $5.8\times10^{-9}$。双实现一致只验证计算，不证明风险预算目标符合投资者负债或收益目标。

\MFQLead{估计误差必须进入报告}

若把 $\widehat\Sigma$ 当作真值，最终权重会显得过度精确。均值与协方差估计误差会被优化器放大，这是经典的研究有效性问题\cite{michaud1989markowitz}。教学实验对收益行做 500 次 IID bootstrap，每次重算组合波动率；点估计 $0.007950$，90\% 区间为 $[0.005984,0.009304]$。区间宽度提醒读者：优化器小数点后六位不等于输入有六位信息。

普通 IID bootstrap 需要独立同分布及统计量的相应正则性；只有可交换性不够。波动聚集、重叠收益、共同宏观冲击或横截面分组会破坏该假设，此时应改用区块、簇或状态分层 bootstrap。重采样单位必须与依赖结构一致。

协方差迹为 $7.415213\times10^{-5}$，权重为 $(0.434184,0.565816)$。这演示 provenance、时点协议和真实输入适配；正文明确禁止把两列宏观增长率冒充可交易资产收益或风险模型有效性证据。

本单元向本路线的综合项目交付风险矩阵、风险预算、估计区间、失败诊断和数据时点协议。

\MFQLead{样本协方差为何在高维下失效}

样本协方差
\[
S=\frac1{T-1}\sum_{t=1}^T(r_t-\bar r)(r_t-\bar r)^{\mathsf T}
\]
在 $N\ge T$ 时必然奇异；即使 $N<T$，最小特征值也可能主要是估计噪声。优化器会在这些看似低风险方向上建立巨大多空仓位。报告维数比 $N/T$、条件数和特征值谱比只报“矩阵可逆”更有信息。

收缩估计 $\hat\Sigma=(1-\delta)S+\delta F$ 把样本矩阵拉向结构化目标 $F$。目标可取常相关、对角或因子模型。$\delta$ 越大方差越低、结构偏差越高；应通过解析收缩或训练期验证选择，而不是为得到漂亮权重手调。

\MFQLead{因子与特异风险的归因}
组合因子暴露为 $B^Tw$，在因子与特异项不相关的假设下，总方差分解为 $(B^Tw)^T\Sigma_f(B^Tw)+w^TDw$。若特异项相关，$D$ 应包含这些非对角协方差；把它强行设成对角会改变风险模型。

\MFQLead{动态协方差与状态依赖}

固定窗口协方差假设窗口内稳定。指数加权、DCC-GARCH 或状态条件协方差允许时变，但引入更多参数和模型风险。危机期相关通常上升，使用平静期矩阵会低估压力风险。

在决策时只能使用此前收益更新矩阵。若当天收盘组合使用包含当天收盘收益的协方差并按同一收盘价交易，需明确是否有足够计算和成交时间，否则应滞后一日。

这些估计对象将在组合风险路线 Capstone 中与优化、尾部压力和限制报告放在同一条证据链上。
\subsection{秩、收缩与风险贡献的同一例子}

记中心化矩阵为 $X$，则 $S=X^TX/(T-1)$。任意 $w$ 有 $w^TSw=\|Xw\|^2/(T-1)\geq0$；且中心化使行和为零，所以 $\operatorname{rank}S\leq\min(N,T-1)$。这是 $N\geq T$ 时必奇异的直接证明。

取 $S=\left(\begin{smallmatrix}1&1\\1&1\end{smallmatrix}\right)$，特征值为 2 和 0。向单位阵收缩得到 $S_\lambda=(1-\lambda)S+\lambda I$，特征值为 $2-\lambda,\lambda$；$0<\lambda\leq1$ 时可逆，条件数为 $(2-\lambda)/\lambda$。但低方差方向被抬高来自结构假设，不是新增观测。

对 $\Sigma=\operatorname{diag}(0.04,0.16)$，逆波动权重为 $(2/3,1/3)$，方差贡献分别为 $0.04(2/3)^2=0.16(1/3)^2=4/225$。总波动为 $\sqrt{8/225}$，波动贡献各为其一半。notebook 同时画出收缩后的谱与权重，并通过重新抽样观察估计变化。

\section{观点后验、尾部目标与稳健实施}
\label{ch:lower-portfolio-optimization}

\begin{itemize}
  % Generated from curriculum/manifest.json; do not edit by hand.
\item 直接先修：下册《风险估计：从协方差到风险预算》、上册第 13 章《凸优化、KKT 与对偶》。

\end{itemize}

观点可信多少与持仓承担多少风险是不同输入。先用一个标量后验计算理解信息合并，再用不确定集推导保守收益目标。


\MFQTransition{Black--Litterman：把观点和置信度显式化}

% Generated by tools/build_figures.py; edit figures/figure-specs.json instead.
\MFQFigure{tex/figures/generated/lower-pf-frontier.pdf}{增加约束或计入成本会缩小可达到的收益与风险组合；最优选择也随风险偏好与换手代价变化。}{lower-pf-frontier}


设先验均值 $\pi$ 的不确定性为 $\tau\Sigma$，观点矩阵 $P$、观点收益 $q$、观点误差协方差 $\Omega$ 满足
\[
 P\mu=q+\eta,\qquad \eta\sim N(0,\Omega).
\]
高斯更新给出
\[
 \mu_{\mathrm{BL}}=
 \left[(\tau\Sigma)^{-1}+P^\top\Omega^{-1}P\right]^{-1}
 \left[(\tau\Sigma)^{-1}\pi+P^\top\Omega^{-1}q\right].
\]
后验均值的不确定性为括号中精度矩阵的逆；若预测收益本身也随机，后验预测协方差还要加回 $\Sigma$\cite{black1992global}。

算例的先验均值为 $(0.05,0.04)$，观点是第一资产相对第二资产超额 $0.03$，$\tau=0.05$，观点方差为 $0.0025$。后验均值为 $(0.05375,0.03)$。透明精度公式与独立 Woodbury 更新最大差 $2.8\times10^{-17}$。当 $\Omega$ 变小，观点权重上升；这表示模型更相信观点，不表示观点变得更真实。

\MFQTransition{CVaR 的线性规划形式}

给定 $N$ 个情景收益 $r_s$，损失 $L_s(w)=-r_s^\top w$。Rockafellar--Uryasev 表示为
\[
 \operatorname{CVaR}_\alpha(w)=
 \min_{z}\left[z+\frac{1}{(1-\alpha)N}
 \sum_{s=1}^{N}(L_s(w)-z)_+\right].
\]
引入辅助变量 $u_s\ge0$ 得到线性规划
\[
 \min_{w,z,u}\quad z+\frac{1}{(1-\alpha)N}\sum_su_s,
 \qquad u_s\ge-r_s^\top w-z,\quad
 \mathbf1^\top w=1,\quad 0\le w_i\le\bar w.
\]
这不是把样本 VaR 当作平滑目标，而是同时选择权重和尾部阈值\cite{rockafellar2000optimization}。

事先确定八情景、$\alpha=0.75$、单资产上限 $0.8$ 的 LP 得到 $(0.2,0.8)$，CVaR 为 $0.013$。独立 0.1 网格枚举重算每个候选的最坏两情景均值，得到相同权重和目标。若枚举忘记权重上限，会错误选择第二资产 0.9；这正说明可行域也是 解析参照 的一部分。

\MFQTransition{把估计误差、成本和可交易性放在一起}

稳健再平衡采用
\[
 \max_w\quad
 \widehat\mu^\top w
 -\rho\lVert U w\rVert_2
 -\frac\gamma2w^\top\Sigma w
 -c\lVert w-w^-\rVert_1,
\]
并约束满仓、长仓、最大权重和不可交易资产权重不变。$U$ 编码期望收益不确定性尺度，$\rho$ 是最坏情形惩罚。它与协方差 ridge 都能平滑权重，但经济对象不同：一个约束均值误差，一个稳定风险矩阵。

事先确定结果从 $(0.4,0.6)$ 调整为 $(0.8,0.2)$。双边换手为 $0.8$；十万元资本、每单位换手成本 $0.001$ 的现金成本为
\[
 100000\times0.001\times0.8=80.
\]
比例成本与收益采用同一口径；乘以资本规模后得到相应的现金金额。若机构采用单边换手，必须显式除以二；基点、百分数和小数的混用会造成百倍或万倍错误。

\MFQTransition{均值--方差优化与输入放大}

无约束均值--方差解为
\[
w^*=\frac1\gamma\Sigma^{-1}\mu.
\]
它清楚展示了两个放大器：均值估计误差直接进入权重，协方差小特征值又通过逆矩阵放大。实际系统必须使用权重界、暴露约束、均值收缩和交易成本，而不是把解析解当推荐组合。

若目标是相对基准 $b$ 的主动组合，使用主动权重 $a=w-b$ 和跟踪误差 $a^{\mathsf T}\Sigma a$。行业中性、beta 中性和换手上限都可写成线性/凸约束。不可交易资产的旧仓位要固定，而不是在求解后裁成零。

\MFQTransition{观点的单位与尾部数据}
相对观点 $P=(1,-1)$ 约束的是两资产均值差，绝对观点 $P=(1,0)$ 约束的是第一资产均值。观点误差方差与收益均值采用同一期间。CVaR 的辅助阈值属于损失坐标，情景数少时线性规划仍有解，却可能由极少数观测决定。

\MFQTransition{稳健优化与分布不确定性}

若均值属于椭球不确定集，最坏情形均值会产生额外范数惩罚；协方差不确定可用情景集合或分布鲁棒方法。稳健半径越大，组合越保守。半径应由估计误差或验证协议决定，不能为获得目标仓位反推。

优化报告应包含 shadow price 或活跃约束：它们说明性能被哪个限制约束。若小幅改变上限就导致完全不同组合，说明解位于尖锐边界，应做参数和输入扰动。

\begin{diagnostic}
求解器返回 optimal 只表示数值模型被解出。若预期收益来自未来修订、成本单位错误或场景不代表尾部，最优解仍是错误问题的精确答案。
\end{diagnostic}

优化结果将在组合风险路线综合项目中接受尾部压力、输入扰动和实施限制的共同复核。
\subsection{精度加权与最坏均值}

标量先验 $\mu\sim N(\pi,v)$，观测观点 $q=\mu+\eta$，$\eta\sim N(0,\omega)$ 独立。配方得到后验均值 $m=(\omega\pi+vq)/(v+\omega)$，方差 $v\omega/(v+\omega)$。$\pi=0.04,q=0.08,v=\omega$ 时，$m=0.06$；若观点方差降为 $v/3$，则 $m=0.07$。高置信度提高观点权重，但不检验观点是否正确。

若均值不确定集为 $\mu=\hat\mu+U^Tz$、$\|z\|_2\leq\rho$，则
\[
 \min_\mu\mu^Tw=\hat\mu^Tw+\min_{\|z\|\leq\rho}z^TUw
 =\hat\mu^Tw-\rho\|Uw\|_2.
\]
Cauchy--Schwarz 给下界；当 $Uw\ne0$ 时取 $z=-\rho Uw/\|Uw\|$ 达到下界。这解释了稳健目标中的范数惩罚来自哪里。notebook 显示标量后验随置信度变化，并枚举两资产权重、计算每个权重的尾部损失和成本。

\section{尾部风险与压力测试}
\label{ch:lower-portfolio-tail}

\begin{itemize}
  % Generated from curriculum/manifest.json; do not edit by hand.
\item 直接先修：下册《组合优化：从观点后验到稳健实施》。

\end{itemize}

分位数只给一个阈值，期望短缺还要把超过阈值的概率质量计入平均。离散样本的边界质量最容易算错，因此先做一个非整数尾部个数的例子。


\MFQTransition{从分位数到尾部平均}

对损失 $L=-R_p$，左分位数定义为
\[
 \operatorname{VaR}_\alpha(L)=\inf\{\ell:\mathbb P(L\le\ell)\ge\alpha\},
\]
积分分位数 ES 为
\[
 \operatorname{ES}_\alpha(L)=\frac1{1-\alpha}
 \int_\alpha^1\operatorname{VaR}_u(L)\,\mathrm du.
\]
离散样本在分位点可能有概率质量，因此软件默认的线性插值、高位次序统计量和左分位数会得到不同答案。本节沿用上述经验左分位数定义，下面直接按概率质量计算。

\MFQTransition{离散 ES 的边界质量}

四个等概率损失为 $(0,1,2,10)$，$\alpha=0.6$。左分位数为 2，因为损失不超过 1 的概率是 0.5，不超过 2 的概率是 0.75。最坏 40\% 包含损失 10 的全部 25\% 质量，以及损失 2 的 15\% 质量，故
\[
 \operatorname{ES}_{0.6}=\frac{0.25(10)+0.15(2)}{0.4}=7.
\]
直接取所有 $L\geq2$ 的条件平均得到 6，是最坏 50\% 的平均，不是这里的 ES。代入辅助阈值公式 $z+\mathbb E(L-z)_+/(1-\alpha)$，在 $z=2$ 得 $2+(0.25\cdot8)/0.4=7$，两条计算一致。

% Generated by tools/build_figures.py; edit figures/figure-specs.json instead.
\MFQFigure{tex/figures/generated/lower-pf-tail.pdf}{VaR 给出分位点，ES 按尾部概率平均损失；分位点上的概率质量有时只计入一部分。尾部观察稀少时，估计的不确定性较大。}{lower-pf-tail}


VaR 一般不满足次可加性；ES 在常见可积条件下是相干风险度量，但更依赖尾部样本。关于市场、信用与操作风险中的标准定义、估计与回测框架，可参见 McNeil、Frey 与 Embrechts\cite{mcneil2015quantitative}。有效尾部观察数为
\[
 n_{\mathrm{tail}}=n(1-\alpha).
\]
例如 250 期观测在 99\% 水平只对应约 2.5 个尾部观测。没有一个普遍适用的“达到 20 或 30 就可信”定理；可信度还取决于尾部形状、依赖和持仓稳定性。读者应比较置信水平和样本窗口变化，并解释少量极端观测对结果的影响。

\MFQTransition{历史尾部、参数尾部和压力测试回答不同问题}

历史法问“过去样本中尾部如何”；参数法问“若分布模型成立，尾部如何”；压力测试问“指定联合冲击下损失多少”。三者不能互相替代。轻尾正态可能漏掉跳跃、波动聚集和相关性上升；历史样本可能没有经历当前持仓结构；压力情景则没有自动的发生概率。

联合情景描述哪些价格和流动性变化同时发生，并将它们作用于同一时点的持仓。历史重演保留曾发生的联合冲击，假设情景改变利率、波动率、价差或相关性，反向压力则从资本或流动性阈值反推触发条件。

\MFQTransition{非线性头寸不能只乘 Delta}

设每个风险因子的线性暴露为 $d_i$、二阶暴露为 $g_i$，冲击为 $x_i$，二阶损益近似
\[
 \Delta V(x)=d^\top x+\frac12\sum_i g_ix_i^2,
 \qquad L(x)=-\Delta V(x).
\]
取冲击方向 $x=(-0.10,0.04)$、$d=(60,40)$、$g=(80,-20)$。线性损失为 $4.4$，加入二阶项后为 $4.016$。这说明非线性重估不一定单调放大损失；Gamma 符号、交叉项、波动率曲面和路径依赖都会改变结果。

完整重估在新的风险因子取值下重新计算价值，适合检验二阶近似在大冲击下的误差。期权需要现货、波动率、时间和利率共同变化；债券需要曲线重构；信用与流动性头寸还需要违约、价差和成交假设。含期权时还会出现 $\mathrm{Vega}\,\Delta\sigma$ 与交叉 Gamma 项。大冲击、障碍附近或临近到期时，Taylor 近似可能失效，应直接重算估值；各风险因子的冲击还需能共同构成一个有意义的情景。

\MFQTransition{反向压力寻找的是第一处阈值穿越}

沿方向 $x$ 放大尺度 $s\ge0$，反向压力问题为
\[
 s^*=\inf\{s:L(sx)\ge C\},
\]
其中 $C$ 是资本或损失阈值。阈值取 $C=10$，需要求解二次不等式。

代入本章给定暴露与冲击，二阶损失为 $L(s)=4.4s-0.384s^2$，阈值 10 的两根为 $3.125$ 与 $25/3$。从零增加压力时先碰到 $3.125$；只检查很大的 $s$，反而可能看到损失重新低于阈值。二阶近似在大冲击下本就可能不可靠，因此求得根后还要判断近似适用范围。若只在某个方向放大，结果仅是该方向的阈值；更广泛的反向压力需要另给允许方向与幅度。Notebook 画出这条曲线，并允许改变暴露与阈值。


\MFQTransition{从样本误差到模型误差}

样本较短时，尾部平均主要由少数极端观测决定。可用块 bootstrap 保留一部分时间依赖，比较重抽样得到的风险估计；区间本身仍受块长度和未观察尾部的限制。极值理论用高阈值以上的超额拟合尾部分布，但阈值太高使数据稀少，太低又使尾部近似不合适，形成偏差与方差的权衡。

VaR 的样本外评价可记录实际损失超过预测分位点的次数：若条件分位点模型正确且分布连续，每期超损的条件概率应为 $1-\alpha$。只比较总次数还不够，连续聚集的超损可能提示条件波动未被解释。ES 还涉及这些极端损失的大小，因此不能仅凭超损次数评价尾部平均。

\section{本路线的综合项目}

把协方差估计、优化约束和尾部压力放在同一个组合中，比较历史、假设和反向压力情景。报告应同时给出风险贡献、成本、约束的影子价格和失败阈值，并说明哪些结果只依赖当前样本，哪些可以迁移到下一段数据。

\section{分层练习}
\begin{enumerate}[leftmargin=*]
  \item \textbf{口述概念：}为什么协方差半正定仍不足以保证估计可靠？
  \item \textbf{笔试推导：}证明方差成分贡献之和等于组合方差。
  \item \textbf{笔试推导：}对角协方差下推导风险平价权重与逆波动率成正比。
  \item \textbf{数值编程：}改变收缩强度，报告最小特征值、条件数和风险平价权重。
  \item \textbf{数值编程：}把 IID bootstrap 改为长度为 4 的移动区块 bootstrap。
  \item \textbf{研究审计：}找出成对删除导致非半正定矩阵的最小例子。
  \item \textbf{研究判断：}因子模型何时比样本协方差更稳定，却更可能遗漏风险？
  \item \textbf{开放任务：}为真实收益数据写出资产池、时间索引、公司行动与缺失值协议。
\end{enumerate}

\section{分层练习}
\begin{enumerate}[leftmargin=*]
  \item \textbf{口述概念：}Black--Litterman 中 $P,q,\Omega,\tau$ 分别编码什么？
  \item \textbf{笔试推导：}从高斯精度相加推导后验均值。
  \item \textbf{笔试推导：}从 $(L-z)_+$ 推导 CVaR 的辅助变量 LP。
  \item \textbf{数值编程：}用情景枚举核对 LP 的权重、阈值和目标值。
  \item \textbf{数值编程：}把最大权重从 0.8 改为 0.6，重算 CVaR 和活跃约束。
  \item \textbf{错误研究包审计：}找出“先优化再裁剪权重”破坏的约束与目标。
  \item \textbf{研究判断：}比较均值不确定性惩罚、协方差 ridge 和交易成本的经济含义。
  \item \textbf{开放任务：}加入部分成交与整数头寸，写出恢复协议和独立现金台账。
\end{enumerate}

\MFQTransition{尾部、压力测试与模型风险题组}
\begin{enumerate}[leftmargin=*]
  \item \textbf{口述概念：}VaR、ES、历史压力和反向压力分别回答什么问题？
  \item \textbf{笔试推导：}解释积分分位数 ES 如何处理分位点概率质量。
  \item \textbf{笔试推导：}推导二阶重估下反向压力的标量方程。
  \item \textbf{数值编程：}对 500 个样本比较 95\% 与 99\% 下的尾部观察数，并解释估计敏感性。
  \item \textbf{数值编程：}加入交叉 Gamma 或期权全重估，比较线性与非线性损失。
  \item \textbf{错误研究材料审计：}审查只报告四个损失样本 99\% ES 的报告。
  \item \textbf{研究判断：}为什么反向压力阈值不等于危机发生概率？
\end{enumerate}

四个样本估计 99\% ES 时，最坏 1\% 概率质量完全落在最大观察上。请说明这一计算在经验分布下的含义，以及为什么它没有提供未观察尾部的信息。
\subsection*{基础衔接：独立计算}
对同样四个损失，求 $\alpha=0.75$ 的 VaR 与 ES。


\chapter{市场微观结构与执行：从订单簿到事件仿真}
\label{ch:lower-microstructure-events}

\begin{itemize}
  \input{tex/generated/prerequisites/lower-microstructure-model}
  \item \textbf{学习产出：}从事件强度和 FIFO 订单簿出发，分解冲击与实施差额，求解执行和库存控制，并用可重放仿真比较策略。
  \item \textbf{研究边界：}公开逐笔成交不含完整队列状态；冲击系数、成交率和事件生成机制都需要估计，合成簿不能冒充真实成交概率。
\end{itemize}

限价单是否成交，取决于它前面的数量怎样减少。先把事件率与等待时间联系起来，再按价格和时间顺序逐笔扣减队列。


\section{事件时间不是把横轴换个名字}

令 $N_t$ 表示截至 $t$ 的事件数，历史信息为 $\mathcal F_{t^-}$。条件强度定义为
\[
 \lambda(t\mid\mathcal F_{t^-})
 =\lim_{h\downarrow0}\frac{\mathbb P(N_{t+h}-N_t=1\mid\mathcal F_{t^-})}{h}.
\]
它描述给定过去后下一瞬间的事件率，而不是“每秒平均几笔”的同义词。齐次 Poisson 过程要求独立平稳增量；等待时间 $\Delta_i$ 独立且服从参数 $\lambda$ 的指数分布。观察 $n$ 段等待时间、总暴露 $T=\sum_i\Delta_i$ 时，
\[
 \ell(\lambda)=n\log\lambda-\lambda T,
 \qquad \widehat\lambda=n/T.
\]
算例中 $n=T=4$，所以 $\widehat\lambda=1$、最大对数似然为 $-4$。这既有手算 解析参照，也由一维数值优化独立重建。

盘中 U 形季节性会让常强度模型产生系统残差。可写成
\[
 \lambda(t\mid\mathcal F_{t^-})=s(t)\widetilde\lambda(t),
\]
其中 $s(t)>0$ 是预先估计并在样本外事先确定的季节因子。若先用全日数据估计 $s(t)$ 再回测同日，会把未来活动率泄漏到早盘。

\section{Hawkes 自激与稳定边界}

指数核 Hawkes 过程写成
\[
 \lambda_t=\mu+\alpha\sum_{t_i<t}e^{-\beta(t-t_i)},
 \qquad \mu>0,\ \alpha\ge0,\ \beta>0.
\]
每个事件瞬时增加强度 $\alpha$，影响以速度 $\beta$ 衰减。分枝比 $\alpha/\beta<1$ 是平稳版本的稳定条件；它不是“拟合得好”的经验阈值。区间 $[0,T]$ 的对数似然为
\[
 \ell=\sum_{t_i\le T}\log\lambda_{t_i}
 -\mu T-\frac{\alpha}{\beta}\sum_{t_i\le T}\left(1-e^{-\beta(T-t_i)}\right).
\]
第一项奖励事件发生处的高条件强度，第二项是补偿子；漏掉补偿子会把强度无限推高。估计时还要报告窗口、初始历史、时间单位、参数约束和残差时间变换，而不能只给三个参数。

% Generated by tools/build_figures.py; edit figures/figure-specs.json instead.
\MFQFigure{tex/figures/generated/lower-ms-hawkes.pdf}{Hawkes 强度在事件到达后跳升并指数衰减；分支比接近一时，聚集持续时间与方差显著上升。}{lower-ms-hawkes}


事先确定事件时刻为 $(0.2,0.7,1.4,2.0)$，观察窗明确延伸到 $T=2.5$，参数 $(\mu,\alpha,\beta)=(0.8,0.3,1.2)$，初始激发为 $h_0=0.1$。此时上式强度还需加 $h_0e^{-\beta t}$，对数似然的补偿项还需减 $h_0(1-e^{-\beta T})/\beta$。递推实现与独立数值积分都得到对数似然 $-2.940341$；若把观察窗缩到最后事件之前，程序直接拒绝。季节调整例的基准强度为 $4/3$，变换后等待残差均值为 1。由此，观察窗、初始历史与季节暴露均不是只写在正文里的口头约定。

\begin{diagnostic}
Poisson、季节性 Poisson 与 Hawkes 是嵌套的研究假设。先检查季节性和时间尺度，再判断剩余聚集是否支持自激；不要用一个 Hawkes 名称同时解释开收盘季节性、新闻簇集和数据时间戳重复。
\end{diagnostic}

事件结构为微观结构路线 Capstone 提供可复核的队列、到达和成交假设，后续控制与仿真只能在这些边界内比较。

\section{订单流与价格变化必须在同一时间协议中}

令 $q_i\in\{-1,+1\}$ 表示第 $i$ 笔主动方向，$\Delta m_i$ 是在事先确定采样规则下的中价变化。最小联合模型为
\[
 \Delta m_i=a+bq_i+u_i,
 \qquad \mathbb E(u_i\mid q_i)=0.
\]
系数 $b$ 是给定对齐协议下的条件均值差，不自动等于结构性永久冲击。算例的 OLS 斜率为 $0.15$，透明协方差公式与带截距最小二乘一致。

真实轨采用 SEC 的公共领域订单位置摘要。它给出五个互斥位置的事件占比与 cancel-to-trade 比率\cite{sec2014orderplacement}；解析器要求占比和为 1、比率非负且所有数字有限。用 $1/(1+\text{ratio})$ 得到的只是“取消或成交”二选一近似下的隐含执行概率，加权值为 $0.044470$。合成样本 另行声明每秒 10 个终止事件的压力尺度，两者相乘才得到每秒耗尽强度；10 并非 SEC 估计。聚合摘要没有订单方向与中价变化，不能估计本节的联合回归系数。

\section{价格优先、时间优先与队列位置}

限价簿状态由订单 $(\mathrm{id},\mathrm{side},p,q,\mathrm{seq})$ 组成。被动簿至少满足：订单编号唯一、活动数量为正、最高买价低于最低卖价；同一价格按序号 FIFO。对订单 $j$，前方数量为
\[
 A_j=\sum_{i:\,p_i=p_j,\,\mathrm{seq}_i<\mathrm{seq}_j}q_i.
\]
若对手方耗尽数量在期限 $h$ 内近似 $D_h\sim\mathrm{Poisson}(\nu h)$，自己的 $Q$ 手全部成交概率是
\[
 \mathbb P(D_h\ge A_j+Q).
\]
算例 $A_j=2,Q=1,\nu h=3$，概率为 $0.576810$。这是明确假设下的队列概率，不是由公开成交快照识别出的现实概率。

% Generated by tools/build_figures.py; edit figures/figure-specs.json instead.
\MFQFigure{tex/figures/generated/lower-ms-order-book.pdf}{订单簿按价格优先、时间优先组织队列；市场单消耗可见深度并留下成交与剩余状态。}{lower-ms-order-book}


执行一笔 4 手市价卖单时，合成簿先吃掉第一张 3 手买单，再使第二张 2 手买单部分成交 1 手；随后撤去剩余 1 手，下一笔市价卖单得到未成交 1 手。这个小计算记录同时覆盖撤单、部分成交、流动性不足和状态不变量。

本单元向本路线的综合项目交付事件方向、队列状态、成交规则和可审计的时间协议\cite{hawkes1971spectra,cont2001empirical}。

\MFQTransition{时钟与标记}
事件除时刻外，还包含方向、数量与价格档。对多维买卖事件，Hawkes 稳定条件推广为核积分矩阵的谱半径小于 1；只拟合总事件数不能识别哪一侧流动性将被消耗。事件时间下相邻两条记录可能间隔极不相同，等待时间本身仍是信息。

\MFQTransition{订单流不平衡与短期价格响应}

简化 order-flow imbalance 可写为买侧新增/撤减与卖侧新增/撤减的净差。短期价格变化回归
\[
\Delta m_{t+h}=\beta\,\mathrm{OFI}_t+\varepsilon_{t+h}
\]
需要同一事件窗口、正确方向和无重叠标准误。$\beta$ 随深度和波动变化，不能从一段样本永久外推。

买卖价差补偿订单处理、库存和逆向选择。知情交易概率上升时，被动报价更容易在价格即将不利移动前成交。高成交率可能是坏消息，不能只把 fill rate 当执行质量。

\begin{diagnostic}
用聚合 K 线拟合“队列成交概率”通常不可识别：缺少订单 ID、队列前方量、撤单和到达顺序。此时应把概率标成压力假设，而不是经验估计。
\end{diagnostic}
\subsection{队列尾概率与 Hawkes 平均强度}

若每次到达恰耗尽一手，前方两手、自身一手，则全部成交要求 $N_h\geq3$。$N_h\sim\operatorname{Poisson}(3)$ 时
\[
 \mathbb P(N_h\geq3)=1-e^{-3}\left(1+3+\frac{3^2}{2}\right)\approx0.5768.
\]
若事件数量随机，就应研究标记数量之和，不能把事件笔数直接当股数。

指数 Hawkes 的平稳平均强度 $\bar\lambda$ 满足 $\bar\lambda=\mu+(\alpha/\beta)\bar\lambda$，因为单位时间的每个历史事件平均贡献核积分 $\alpha/\beta$。所以 $\bar\lambda=\mu/(1-\alpha/\beta)$。$\mu=0.8,\alpha=0.3,\beta=1.2$ 时为 $16/15$；有限窗口从空历史开始，平均强度未必立即达到平稳值。

notebook 用同价订单数量 $(3,2)$ 演示四手市价单先成交第一单 3 手、再成交第二单 1 手，并计算不同前方数量的成交概率。

\section{执行、做市与库存控制}
\label{ch:lower-microstructure-control}

\begin{itemize}
  % Generated from curriculum/manifest.json; do not edit by hand.
\item 直接先修：下册《事件与订单簿：从条件强度到队列成交》、上册第 13 章《凸优化、KKT 与对偶》、上册第 14 章《浮点数、数值线代与病态问题》、上册第 15 章《Monte Carlo、bootstrap 与方差缩减》。

\end{itemize}

立即成交会付出冲击，等待又承担价格与未完成风险。将剩余数量作为状态，可以把这两种代价写在同一个决策方程里。


\MFQLead{先定义实施差额的基准}

对买入母单，决策时到达价为 $S_0$，第 $k$ 期实际成交 $x_k$、成交价 $P_k$，实现差额为
\[
 \mathrm{IS}=\sum_k x_k(P_k-S_0)+R_K(S_K-S_0)+C,
\]
其中 $R_K$ 是终止时未完成数量，$C$ 是显式费用。若只统计已成交部分，应明确这是 partial-fill IS，不能把未完成任务静默删除。卖单符号相反，所有价格和费用必须统一为现金或基点口径。

教学冲击模型写成
\[
 P_k=S_0+\gamma X_{k-1}+\eta x_k,
\]
其中 $X_{k-1}=\sum_{j<k}x_j$，$\eta x_k$ 是本期临时冲击，$\gamma X_{k-1}$ 是累计永久冲击。永久冲击改变后续基准，临时冲击只改变当前成交价；把二者合并成一个系数会扭曲日程的跨期权衡\cite{almgren2001optimal}。

事先确定计划 $(3,3,3)$ 在完成率 $(1,0.5,0)$ 下实际成交 $(3,2,0)$。以 $S_0=100,\eta=0.2,\gamma=0.1$ 计算，partial-fill IS 为 $3.2$，剩余 4 手必须保留在状态中。

\MFQLead{随机完成率与停止规则}

设计划量 $u_k$，随机完成率 $F_k\in[0,1]$，实际成交为
\[
 x_k=\min\{R_k,\operatorname{round}(u_kF_k)\},
 \qquad R_{k+1}=R_k-x_k.
\]
完成率不是预测误差的装饰项：它决定剩余库存、后续冲击和最终违约状态。停止规则可以是固定期限、累计损失阈值、市场状态或风险限制；一旦触发，应输出 stopped、剩余数量和采用的处置规则。

动态规划的一般值函数记为 $V_t(x)$，表示时点 $t$、剩余库存 $x$ 下的最小未来目标；离散状态取 $(k,R_k)$。一个最小成本递推是
\[
 V_k(R)=\min_{0\le x\le\min(R,x_{\max})}
 \left\{\eta x^2+\gamma(I-R)x+\phi(R-x)^2+V_{k+1}(R-x)\right\},
\]
终端只允许 $R=0$。事先确定六手、三期问题的最优日程是 $(3,2,1)$，目标 $19.2$；完整枚举所有可行路径得到相同答案。DP 证明的是离散目标的最优性，不证明冲击参数、二次风险或硬完成约束适合现实委托。

\MFQLead{库存如何影响下一次报价}

令中价 $m_t$、半点差 $h$、库存 $q_t$、库存偏移系数 $\kappa$。最小双边报价为
\[
 r_t=m_t-\kappa q_t,
 \qquad b_t=r_t-h,\quad a_t=r_t+h.
\]
若 bid 成交，$q_{t+1}=q_t+1$、现金减少 $b_t$；若 ask 成交，库存减一、现金增加 $a_t$。下一轮必须用 $q_{t+1}$ 计算报价。算例从零库存开始，首轮 bid 成交后库存为 1，下一轮报价从 $99.5/100.5$ 下移为 $99.4/100.4$。双边同时下移，使买入变得不积极、卖出更容易，才形成真正反馈。

% Generated by tools/build_figures.py; edit figures/figure-specs.json instead.
\MFQFigure{tex/figures/generated/lower-ms-feedback.pdf}{每次成交改变现金与库存，新的库存影响下一次报价；库存限制也会改变后续可选动作。}{lower-ms-feedback}


成交改变库存，新的库存又影响下一次报价，这构成动态反馈。如果预先生成与库存无关的全天报价，再计算库存路径，研究的是另一种不含库存反馈的策略；如果报价使用了尚未发生的库存或事件，则超出了当时的信息。进一步的模型可以考虑成交强度不确定性、价格跳跃、库存边界和撤单延迟。

\MFQTransition{Almgren--Chriss 的风险--冲击折中}

离散执行把库存 $x_k$ 从 $X$ 降到 0，交易量 $n_k=x_{k-1}-x_k$。典型目标为期望成本加方差惩罚：
\[
\min_{n_1,\ldots,n_K}
\mathbb E[C]+\lambda\operatorname{Var}(C),
\qquad \sum_kn_k=X.
\]
临时冲击影响本次成交，永久冲击改变后续基准。$\lambda=0$ 更愿意等待低冲击，$\lambda$ 大则加快执行以降低价格风险。参数应从订单级数据估计并做容量压力，不能从公开日频量价唯一识别。

% Generated by tools/build_figures.py; edit figures/figure-specs.json instead.
\MFQFigure{tex/figures/generated/lower-ms-control.pdf}{执行与做市控制根据库存和剩余时间选择动作，再由随机成交改变下一状态；终端惩罚约束未完成风险。}{lower-ms-control}


参与率、最小成交单位、最大切片和收盘前强制完成使问题成为约束动态规划。未完成状态必须在终端施加机会成本或硬约束，否则最优策略可能选择永不交易。

\MFQTransition{做市的库存控制}

简化 Avellaneda--Stoikov 中，库存 $q$ 改变保留价格
\[
r_t=S_t-q\gamma\sigma^2(T-t).
\]
多头库存使报价中心下移，鼓励卖出；风险厌恶、波动和剩余时间放大调整。最优价差还取决于订单到达对距离的敏感度。模型假设连续 Brownian 中价和指数到达，真实跳跃、队列与延迟都可能破坏闭式。

报价状态机还要处理库存上限、单边停报、撤单确认、自成交防护和 stale quote。控制动作必须在收到事件后生成，并等待交易所确认；本地认为已撤单不代表市场上已撤。

\MFQTransition{部分可观测与鲁棒控制}

真实状态包含不可见的短期 alpha、对手意图和未来流动性，执行问题更接近 POMDP。滤波器提供 belief state，控制器基于 belief 行动。若状态估计和控制用同一未来数据平滑，会产生前视。

模型参数不确定时可做情景 DP、分布鲁棒或保守策略。最简单的鲁棒基线是限制参与率、库存和单步动作，并在成交率低于阈值时切换预定义应急日程。复杂 RL 只有在能恢复小型 DP 解析参照 后才值得使用。

\begin{diagnostic}
控制器的目标函数必须包含未完成、库存和成本。若奖励只含已实现价差，代理可以通过不成交或积累巨大库存获得虚假高分。
\end{diagnostic}

控制规则将在微观结构路线综合项目中与事件仿真、执行成本和库存压力一起复核。
\subsection{两期执行的风险权衡}

需要买入总量 $Q$，第一期买 $x$，第二期买 $Q-x$，暂设均可成交。临时冲击成本为 $\eta[x^2+(Q-x)^2]$，持有未完成数量的风险惩罚为 $\phi(Q-x)^2$。对 $x$ 求导得
\[
 2\eta x-2(\eta+\phi)(Q-x)=0,
 \qquad x^*=\frac{\eta+\phi}{2\eta+\phi}Q.
\]
$\phi=0$ 时均分；$\phi$ 增加时提前买入更多。$Q=6,\eta=1,\phi=2$ 时 $x^*=4.5$，第二期为 1.5；若只能买整数，比较 $x=4,5$，两者目标都为 28，形成并列最优。

报价反馈也可直接算现金。中价 100、半价差 0.5，从零库存买入一股于 99.5，现金为 $-99.5$、库存为 1，盯市财富为 0.5。若中价随即跌到 99，财富变成 $-0.5$；成交并不保证赚钱，逆向选择可以超过点差收益。notebook 枚举整数日程并显示每次成交后的库存与现金。

\section{事件时间仿真与研究项目}
\label{ch:lower-microstructure-simulation}

\begin{itemize}
  % Generated from curriculum/manifest.json; do not edit by hand.
\item 直接先修：下册《执行与做市：价格冲击与库存反馈》。

\end{itemize}

仿真将到达、报价、成交和库存连续串起来。比较策略时要重复生成市场路径，并保留同一路径上的结果差，才能区分策略作用与随机波动。


\MFQTransition{先确定生成机制，再比较策略}

% Generated by tools/build_figures.py; edit figures/figure-specs.json instead.
\MFQFigure{tex/figures/generated/lower-ms-execution.pdf}{执行轨迹在冲击成本与库存风险之间权衡；前置成交降低尾部库存，却提高早期冲击。}{lower-ms-execution}


一个最小事件生成器需要明确：等待时间、订单方向、价格变化、委托大小和终止条件。例如
\[
 \Delta t_i\sim\operatorname{Exp}(\lambda),\quad
 q_i\in\{-1,+1\},\quad
 \Delta m_i=\beta q_i+\sigma\sqrt{\Delta t_i}\varepsilon_i.
\]
这些方程是数据生成假设，不是对现实的事实陈述。若要引入季节或 Hawkes 聚集，应逐层增加并分别做残差诊断；模型复杂度不能替代可识别性。

比较静态报价与库存控制时，可以让两者使用相同的 $(\Delta t_i,q_i,\varepsilon_i)$ 和成交均匀数，但不能强制共享最终成交结果。每个策略用自己的报价计算成交阈值，再把同一个均匀数与阈值比较；因此报价能够改变成交、库存与收益，同时随机源仍可配对。共享底层随机数有助于减少比较的噪声；独立路径也能形成有效估计，只是需要正确计算两组的不确定性。共享随机源不等于强制两种策略具有相同成交。

\MFQTransition{共同随机数何时降低方差}

对同一路径上的结果 $Y_A,Y_B$，差值 $D=Y_A-Y_B$ 满足
\[
 \operatorname{Var}D=\operatorname{Var}Y_A+\operatorname{Var}Y_B
 -2\operatorname{Cov}(Y_A,Y_B).
\]
若方差都为 4、相关系数为 0.75，则差值方差为 $8-6=2$；独立抽样则为 8。100 条独立配对路径的均值差标准误为 $\sqrt{2/100}$，是独立比较的一半。若协方差为负，配对反而可能增加方差。

一个成交例子也能说明为什么不应共享最终成交：同一均匀数为 0.4，A 报价对应成交概率 0.3，B 对应 0.6，则 A 不成交、B 成交。底层随机源相同，策略后果不同。notebook 用这一规则比较静态与库存偏移报价，在多条独立路径上计算损益差及最大库存，而不把一条路径的胜负当作总体结论。

做市商的损益由现金与期末库存估值构成：主动买单在卖价成交时，做市商收到现金、库存减少；主动卖单在买价成交时方向相反。库存反馈规则先根据当前库存调整报价，报价又影响成交概率，成交后的库存再进入下一次决策。配套实验在多条路径上比较这些递推产生的损益与库存分布。

\MFQTransition{同时比较损益与库存风险}

把终点累计损益记为 $\operatorname{PnL}_T$。可同时考察
\[
 \text{PnL},\quad \max_t|q_t|,\quad q_T,\quad
 \text{fill rate},\quad \text{turnover},\quad \text{shortfall},
\]
这些量的路径分布描述平均表现以外的风险。若只看平均 PnL，策略可以通过无限库存或未完成母单获得虚假优势。做市 PnL 还应拆成点差捕获、持仓重估、费用与冲击；执行策略则应以相同到达价和母单任务比较。

动态规划策略与仿真统计承担不同证据：DP 在给定离散状态、转移与目标下给出最优控制；Monte Carlo 描述该控制在声明分布下的结果分布。DP 最优不意味着模型正确，仿真均值领先也不意味着总体或未来市场领先。

\MFQTransition{公开汇总数据能说明什么}

配套材料使用公开的 SEC 汇总统计设定一个低成交率的比较情形。把其数值作为模型输入后，实际模拟成交仍由每种策略自己的报价与成交规则决定。

公开汇总数据可以帮助设定低成交率等比较情形，但没有提供具体订单的队列位置、交易发起方或网络延迟。由撤单成交比率换算为执行概率，还额外假设每个终止事件只有取消或成交两类结果。因此，这种换算只能在所给简化模型内解释。

\MFQTransition{比较对象要一致}
两种策略从相同初始资金和库存出发，各自按报价决定成交。期末库存按同一价格估值，费用与未完成任务采用相同口径。这样才能把损益差解释为策略在所给模型中的差异\cite{hasbrouck2007empirical}。

\MFQTransition{事件先后顺序与策略可用信息}

事件时间模型按发生顺序推进：市场变化形成新的可见信息，策略据此报价或下单，订单经过延迟后才可能成交。把数据到达、决策、交易所确认和成交分开，才能表达动作发生后、成交之前的价格风险。

外生订单流、价格创新、延迟和成交抽样可以分别生成。使用共同随机数比较策略时，共享外生随机变化仍允许各策略采取不同动作；由此产生的成交和库存路径可以分叉。

\MFQTransition{校准与验证不是同一件事}

校准让模拟器匹配选定统计量，例如到达率、点差、深度、订单大小和自相关；验证检查未参与校准的统计量，例如价格响应、队列等待和极端流动性。若所有统计量都用于拟合，就没有留出证据。

一条仿真路径同时受到数量与时间关系约束：成交数量不能超过可成交委托，现金变化应与成交方向一致，未来事件不能影响过去的决策。均值、尾部和时间依赖则描述多条路径的分布性质。个别路径上的数量关系与整体分布是否合适，需要分别考察。

\MFQTransition{策略比较与置信区间}

对策略 A/B 使用相同外生随机路径，比较配对差 $D_b=Y_b^A-Y_b^B$，若两个结果正相关，其标准误可小于两组独立模拟。报告均值差、置信区间、尾部分位和失败率。若策略改变市场状态，后续事件可能分叉，但仍可共享底层随机数。

调参路径与最终评价路径分开。增加仿真次数只降低给定模拟器内的 Monte Carlo 误差，不降低模拟器错设。结论应写成“在这些生成假设下”，而不是“市场中”。

\MFQTransition{压力场景与代理交互}

压力场景包括到达率骤降、撤单激增、价格跳跃、延迟上升、单边流动性枯竭和交易所拒单。这些情形分别改变流动性、价格或信息延迟；比较它们有助于找出策略依赖最强的假设。

多代理仿真可加入做市者、噪声交易者、执行者和信息交易者。产生相似 stylized facts 不代表机制被识别：不同代理规则可能生成相同分布。应做消融，逐类移除代理并观察哪些统计量变化。

\MFQTransition{历史回放与反事实的局限}

历史消息回放可以重现已经发生的事件顺序；实时观察而不实际下单，可以检验信号能否在相应时点形成。两者都没有观察到假想订单真正提交后会产生的成交与市场影响，因此这部分仍依赖模型假设。

把仿真结果推广到市场，需要考虑订单对流动性和其他参与者行为的反馈。本书的综合项目保留这些限制，重点比较给定模型内的策略后果。

\begin{diagnostic}
某条路径上的领先可能来自随机波动。增加独立路径可以估计平均差异；更换流动性状态或生成机制，则用来研究结论对模型假设的依赖。两种比较回答不同的问题。
\end{diagnostic}

\section{本路线的综合项目}

选择一个报价或执行问题，把事件流、成交和库存递推连起来。用多条路径比较策略，给出损益差的置信区间，并观察延迟及未成交怎样改变结果。再说明哪些假设能够由历史消息检验，哪些涉及尚未观察到的订单后果。

\section{综合练习}
\begin{enumerate}[leftmargin=*]
  \item \textbf{口述概念：}条件强度与无条件平均事件率有什么差别？
  \item \textbf{笔试推导：}由指数等待时间似然推导 $\widehat\lambda=n/\sum_i\Delta_i$。
  \item \textbf{笔试推导：}解释 Hawkes 似然中补偿子为何不可省略。
  \item \textbf{数值编程：}扫描 $\alpha/\beta$ 并画出单个事件后的强度衰减。
  \item \textbf{订单簿：}手工推进两档 FIFO 簿，记录撤单、部分成交和未成交。
  \item \textbf{研究判断：}找出把 maker side 当主动方时报告会发生的错误。
  \item \textbf{错误研究材料：}某报告只用成交笔数拟合 Hawkes，却不处理盘中季节性；列出至少三项缺失证据。
  \item \textbf{开放任务：}设计一个带盘口深度、撤单和队列位置的数据协议，明确可知时点和失败状态。
\end{enumerate}

\MFQLead{执行、做市与库存控制题组}
\begin{enumerate}[leftmargin=*]
  \item \textbf{口述概念：}临时冲击与永久冲击怎样进入不同期的成交价？
  \item \textbf{笔试推导：}手算冻结日程 $(3,2,0)$ 的 partial-fill IS。
  \item \textbf{笔试推导：}写出 DP 终端条件，并解释为何剩余库存不能零成本消失。
  \item \textbf{数值编程：}枚举六手三期所有日程，与 DP 结果逐项比较。
  \item \textbf{控制实验：}提高库存惩罚，记录报价中心与期末库存如何变化。
  \item \textbf{研究判断：}识别“先生成报价、后计算库存”的前视错误。
  \item \textbf{错误研究包：}某执行报告只给平均成交价，没有到达价、未完成量和停止原因；写出修复清单。
  \item \textbf{开放任务：}加入随机完成率与最大损失停止，设计可复现的压力矩阵。
\end{enumerate}

\MFQTransition{事件仿真、策略比较与综合项目题组}
\begin{enumerate}[leftmargin=*]
  \item \textbf{口述概念：}为什么共同随机数能提高控制器差异比较的精度？
  \item \textbf{笔试推导：}写出静态半点差策略在给定方向与价格变化下的逐事件 PnL。
  \item \textbf{统计设计：}为 PnL、期末库存和最大库存设计置信区间。
  \item \textbf{数值编程：}把齐次 Poisson 换为日内季节强度，并验证时间变换残差。
  \item \textbf{反例实验：}重复记录一次成交或改变事件先后顺序，手算库存与现金会怎样出错，并指出被破坏的数量或信息关系。
  \item \textbf{研究判断：}审阅只报告单路径 PnL 的做市研究材料。
  \item \textbf{真实数据：}解释 SEC 聚合 cancel-to-trade 比率为何不能识别自己的队列成交概率。
  \item \textbf{综合项目：}建立报价、成交与库存相互影响的小模型，比较两种策略，并解释参数变化怎样影响结论。
\end{enumerate}

\subsection*{基础衔接：独立计算}
两策略在同一路径上的损益完全相同，配对差标准误为零能否证明市场盈利？



\appendix

\chapter{Brainteaser 分级提示}
\begin{enumerate}[leftmargin=*]
  \item 从最小存活人数逆推，先买下一轮收益为零者的票。
  \item 状态写成三者位置与轮到谁；在有限图上做胜负标记。
  \item 对合法岸边状态做 BFS；最短路层数就是下界。
  \item 看相邻生日星期增量是 $1$ 还是 $2$。
  \item 单调阈值来自单交叉性，但阈值依赖支付矩阵和信念。
  \item 一根两端同时点燃，另一根一端点燃。
  \item 异常状态有 $24$ 个；三次天平结果有 $27$ 个叶子。
  \item 用 $v_p(n!)=\sum_{k\ge1}\lfloor n/p^k\rfloor$。
  \item 第一轮五组，第二轮冠军组；随后只保留仍可能前三的名次格。
  \item 特征方程根为 $1,2$；注意递推的主导根。
  \item 投影宽度而非总体积或空间对角线控制可放入性。
  \item $0,1,2$ 必须同时出现在两枚方块上；把 $6$ 翻转作 $9$。
  \item 一次是/否只有两个叶节点；增加一次换门或第二问。
  \item 区分“消息已收到”与“发送方知道消息已收到”。
  \item 追踪红球数的奇偶性。
  \item 从左到右，前一位目标唯一决定下一次是否翻转。
  \item 期望相同不代表尾部、下行风险、流动性或效用相同。
  \item 任取一堆含 $k$ 枚，把这一堆全部翻面。
  \item 从“混合”标签袋抽，因为它一定不是混合。
  \item 把“不知道”视为排除一个世界，而非没有信息。
  \item 先算 $1+\cdots+12=78$，每块目标和若相等应为 $19.5$。
  \item 先求 $x+y$ 和 $x^2+y^2$，再得到 $xy$。
  \item 用二进制权重取样；读数编码缺失权重和。
  \item 最坏先取三种不同颜色；第二问分别构造下界。
  \item 握手图中每条边贡献两个度。
  \item 固定一人，他至少有三条同类关系。
  \item 碰撞时交换身份标签，几何轨迹不变。
  \item 为每枚物品分配 $\{-1,0,1\}^3$ 非零码。
  \item 第一人公布总状态的模校验，而不是自己的状态。
  \item 用 $10\equiv1\pmod9$ 展开十进制位值。
  \item 观察两种颜色计数差在模 $3$ 下怎样变化。
  \item 任取 $n$ 枚为一堆，再把这一堆全部翻面。
  \item 每次操作只会让碎片数增加一。
  \item 同向追及看相对速度 $u-v$，反向相遇看 $u+v$。
  \item 对 $x=\sqrt2^{\sqrt2}$ 是否有理分两种情况。
  \item 第一人的回答可用来公布总颜色和的模校验。
\end{enumerate}

\chapter{多因子模型、估计与研究提示}
\begin{enumerate}[leftmargin=*]
  \item 先写清被解释变量、横截面或时间轴以及条件信息集。
  \item 构造一个单调但非线性的关系，再加入极端值或并列值。
  \item 标准误来自斜率时间序列，而不是同一期资产数量。
  \item 第一遍沿时间估 beta，第二遍沿资产估 lambda；beta 是估计量。
  \item 先保持共同 alpha，再制造横截面 alpha 差异。
  \item BH 的保证依赖检验统计量之间的结构，不能只看阈值公式。
  \item 寻找可交易因子、beta 解释、零 alpha 与误差条件。
  \item 任何看过最终期后改变的规则都消耗了最终期的独立性。
\end{enumerate}

\MFQLead{估计与正则化提示}
\begin{enumerate}[leftmargin=*]
  \item 分别检查条件数、抽样协议和领域机制。
  \item 相关特征会争夺同一稀疏位置；做重采样稳定性检查。
  \item 对目标函数求梯度并令其为零。
  \item 固定其他坐标，把一维目标拆成二次项与绝对值项。
  \item 对齐截距、损失缩放和标准化，再比较系数。
  \item 零列无更新分母；近共线列表现为小奇异值。
  \item 把 lambda 选择写进滚动训练/选择协议。
  \item 逐项检查资产定义、可交易性、样本量和识别假设。
\end{enumerate}

\MFQTransition{组合实施与应用分析提示}
\begin{enumerate}[leftmargin=*]
  \item 沿数据流逐个写出数组形状、时点和单位。
  \item 等权控制组内集中度；市值权重改变暴露；重叠持有改变依赖。
  \item 两期毛收益分别为 $0.035$ 与 $0.030$，再处理换手。
  \item 有效权重是仍存续 vintage 的平均，收益因此重叠。
  \item 固定同一资产池和成本模型，只改变预先声明的一项。
  \item 每期毛收益由 $w_t^{\mathsf T}r_{t+1}$ 得到，成本由实际交易量计算。
  \item 恢复 500 项完整检验族；恢复不了就降级结论。
  \item 停牌、借券与成交限制会改变可持有资产和实际持仓。
  \item 将回归系数、分组权重和净收益联系起来，说明每种比较保持了什么、改变了什么。
\end{enumerate}

\chapter{时间序列、状态估计与统计套利提示}
\begin{enumerate}[leftmargin=*]
  \item 单位根零假设位于非平稳边界。
  \item 一个是短期线性共动，另一个是水平线性组合的积分阶。
  \item 使用 $\phi=e^{-\kappa\Delta}$。
  \item 将 $w_{t-1}$ 加减并整理滞后差分。
  \item 特别检查 $\phi\le0$ 与 $\phi\ge1$。
  \item 对齐确定性项与滞后阶后再比较。
  \item 配对搜索必须位于训练期，测试期只评一次。
\end{enumerate}

\chapter{估计与状态推断提示}
\begin{enumerate}[leftmargin=*]
  \item 分别是 $y_{1:t-1}$、$y_{1:t}$、$y_{1:T}$。
  \item 似然对状态编号置换不变。
  \item 从联合正态条件分布或最小均方误差投影出发。
  \item 先乘转移矩阵，再乘发射密度并归一化。
  \item $Q/R$ 越大，状态越愿意追随新观测。
  \item 对齐标签后再比较，不能只看编号。
  \item 在决策日检查估计使用的最后观测日。
\end{enumerate}

\chapter{统计套利研究提示}
\begin{enumerate}[leftmargin=*]
  \item 一个切断跨边界标签，另一个隔离已经参与选择的验证期。
  \item 订单可能部分成交、未成交或延迟。
  \item 从空仓开始逐项计算 $|w_t-w_{t-1}|$。
  \item 冻结仓位直到年龄达到 2。
  \item 让成交比例和原因码进入账本。
  \item 检查拟合截止与 fill fraction。
  \item 没有合法时点数据时，明确降级为方法演示。
\end{enumerate}

\chapter{ML 模型与表示提示}
\begin{enumerate}[leftmargin=*]
  \item shape 是语义轴，dtype/device 决定算子，梯度状态决定是否进入反向图。
  \item 相邻输入行不等于相邻特征值。
  \item 依次写仿射、激活、仿射和链式法则。
  \item 关注 $QK^{\mathsf T}$ 在同步置换下的变化。
  \item 均值应为零差，四种顺序模型应为正差。
  \item 同时冻结配置与参数字节。
  \item 查找是否有位置、因果 mask 和置换负例。
  \item 比较可训练参数、任务数据与基础表示是否更新。
\end{enumerate}

\chapter{现代人工智能方法提示}
\begin{enumerate}[leftmargin=*]
  \item shape 是语义轴，dtype/device 决定数值算子，梯度状态决定是否进入反向图。
  \item 比较局部感受野、递归状态和内容寻址；所有路径都要服从因果可见集。
  \item GNN 增加关系图；DRL 增加状态、动作、转移和长期奖励。
  \item 依次写仿射、激活、仿射和链式法则；残差提供恒等梯度路径。
  \item 输出是相邻差；检查填充是否让 $t$ 读取 $t+1$。
  \item 同步置换 $Q,K,V$ 只会重排输出；手算 $\operatorname{softmax}(0,1)$。
  \item 图边要有生效时间；RL 转移模型必须响应动作。
  \item 均值应不变，顺序模型应变化；两类 mask 分别处理未来和不存在位置。
  \item checkpoint 同时绑定配置、参数和输入 schema。
  \item 为每个决策日截断边表，而不是今天的全图。
  \item 查是否有可训练序列层、位置结构和置换负例。
  \item 比较可训练参数和基础表示是否更新；记录权重训练截止与语料版本。
  \item 先用两步、有限状态的动态规划给出解析策略。
  \item 项目只选一条主线，但必须保留树或线性基线并连接净收益。
\end{enumerate}

\chapter{ML 验证与监控提示}
\begin{enumerate}[leftmargin=*]
  \item purge 看标签跨度，embargo 看选择后的邻接污染。
  \item 外层测试评价整个选择程序。
  \item 代入 $7+2<10$ 与 $11+2<14$。
  \item 校准只改变概率刻度。
  \item 每折训练索引必须严格早于验证索引。
  \item 逐个 $k$ 比较集合交并比。
  \item 一旦反馈，测试集已经成为验证集。
  \item 响应需在监控前冻结。
\end{enumerate}

\MFQLead{Alpha 决策与研究包提示}
\begin{enumerate}[leftmargin=*]
  \item 点预测、次序和最终效用不是同一对象。
  \item 成交比例作用于订单增量。
  \item 第二期订单为 $[-2,1,1]$。
  \item 对分数放大倍数做尺度分析。
  \item 保留原因码和旧仓位。
  \item 三个时间都必须不晚于相应决策边界。
  \item 检查仓位、换手、成本、容量和未成交。
  \item 每条主张绑定命令、oracle 或明确限制。
\end{enumerate}

\chapter{随机分析与无套利提示}
\begin{enumerate}
  \item 先写清随机对象、极限模式和有限样本对象。
  \item 检查每个系数在哪个时点可测。
  \item 展开 $(W_{i+1}-W_i)^2+2W_i(W_{i+1}-W_i)$。
  \item 对 $f(s)=\log s$ 计算一、二阶导数。
  \item 聚合后再平方，不要从不同种子重采样。
  \item 先手算 $\theta=(\mu-r)/\sigma$。
  \item 有限常数的指数仍有限。
  \item “充分条件失败”与“结论为假”不是同一命题。
\end{enumerate}

% Source-mapped interview replacements.  The questions are independent Chinese
% adaptations; only the short source titles are retained for auditability.
\providecommand{\MFQInterviewFollowup}[1]{\par\noindent\textbf{面试官追问：}#1}

\providecommand{\MFQMappedUpperChFiveQuestions}{
\subsection*{笔试推导（来源映射替换）}
\begin{enumerate}[nosep,leftmargin=*]
\item \MFQInterviewMeta{笔试推导}{12 分钟}{中}{一阶条件与边界比较}{GB-3-02 Maximum and minimum}给定 $f(x)=x^4-4x^2+cx$，写出随参数 $c$ 变化的驻点判定流程；说明只解 $f'(x)=0$ 为什么不足以确定全局最小值。
\end{enumerate}
\subsection*{研究判断（来源映射替换）}
\begin{enumerate}[nosep,leftmargin=*]
\item \MFQInterviewMeta{研究判断}{12 分钟}{中}{Taylor 余项与误差控制}{GB-3-08 Taylor series}在 $x=0$ 展开 $\log(1+x)$ 到三阶，用余项给出 $|x|\le0.2$ 的统一误差界；据此审计“在所有收益率上都可安全替代对数收益”的报告。
\end{enumerate}}
\providecommand{\MFQMappedUpperChFiveHints}{\par\textbf{来源映射题提示：}先比较全部驻点、边界与无穷远行为；Taylor 题用四阶导数在区间上的上界控制 Lagrange 余项。}
\providecommand{\MFQMappedUpperChFiveSolutions}{
\section*{来源映射替换题}
\begin{enumerate}[leftmargin=*]
\item \textbf{GB-3-02。}$f'(x)=4x^3-8x+c$ 只给候选点；还需用 $f''(x)=12x^2-8$ 分类，并比较所有局部极小值。因首项系数为正，$f(x)\to\infty$，故全局最小一定在驻点中；若定义域有边界还要加入端点。参数追问应讨论三次方程重根处的分岔，而不能只给数值根。
\item \textbf{GB-3-08。}$\log(1+x)=x-x^2/2+x^3/3+R_4$，且 $f^{(4)}(u)=-6/(1+u)^4$。在 $|x|\le0.2$ 上 $|R_4|\le |x|^4/[4(0.8)^4]$。下一项为负，因此正 $x$ 附近三阶截断高估对数收益；负 $x$ 时应直接依余项而不是机械套交错级数。
\end{enumerate}\MFQInterviewFollowup{GB-3-02：定义域变成紧集时还要加哪些候选点？GB-3-08：如何用同一余项界选取收益率阈值？}}

\providecommand{\MFQMappedUpperChSixQuestions}{
\subsection*{笔试推导（来源映射替换）}
\begin{enumerate}[nosep,leftmargin=*]
\item \MFQInterviewMeta{笔试推导}{12 分钟}{中}{QR 与最小二乘}{GB-3-16 QR decomposition}对满列秩 $X=QR$，从残差正交性推导 $R\hat\beta=Q^Ty$，并解释为何不必形成 $X^TX$。
\end{enumerate}
\subsection*{口述概念（来源映射替换）}
\begin{enumerate}[nosep,leftmargin=*]
\item \MFQInterviewMeta{口述概念}{10 分钟}{中}{谱与经济方向}{GB-3-17 Determinant, eigenvalue and eigenvector}解释协方差矩阵的特征向量、特征值和行列式分别表达什么；零行列式为什么不等于“所有风险都为零”？
\end{enumerate}
\subsection*{研究判断（来源映射替换）}
\begin{enumerate}[nosep,leftmargin=*]
\item \MFQInterviewMeta{研究判断}{10 分钟}{中}{半正定与模型边界}{GB-3-18 Positive semidefinite and positive definite matrices}证明 $B^TB$ 半正定，再审计“将任意估计矩阵换成 $B^TB$ 就保留了原风险模型”的报告。
\end{enumerate}
\subsection*{数值编程（来源映射替换）}
\begin{enumerate}[nosep,leftmargin=*]
\item \MFQInterviewMeta{数值编程}{12 分钟}{中}{LU/Cholesky 分解与失败诊断}{GB-3-19 LU decomposition and Cholesky decomposition}对比 LU 与 Cholesky 的前提和计算结构；实现带对角门禁的 Cholesky，并用一个对称但不正定矩阵证明失败是模型证据而非“库坏了”。
\end{enumerate}}
\providecommand{\MFQMappedUpperChSixHints}{\par\textbf{来源映射题提示：}利用 $Q^TQ=I$；谱题区分零方向与全部方向；$x^TB^TBx=\|Bx\|^2$；Cholesky 的根号参数必须为正。}
\providecommand{\MFQMappedUpperChSixSolutions}{
\section*{来源映射替换题}
\begin{enumerate}[leftmargin=*]
\item \textbf{GB-3-16。}$\|y-QR\beta\|^2=\|Q^Ty-R\beta\|^2+\|(I-QQ^T)y\|^2$，故最小化第一项得到 $R\hat\beta=Q^Ty$。形成正规方程会把奇异值平方，使条件数平方。
\item \textbf{GB-3-17。}特征向量是风险二次型的正交主方向，特征值是这些方向的方差；行列式是特征值乘积，反映总体积尺度。行列式为零只说明至少一个线性组合方差为零，其他方向仍可有正风险。
\item \textbf{GB-3-18。}$x^TB^TBx=\|Bx\|_2^2\ge0$，所以半正定。其正定当且仅当 $Bx=0$ 只有零解，即 $B$ 满列秩。
\item \textbf{GB-3-19。}LU 面向一般方阵且通常需选主元；Cholesky 利用对称正定结构，将 $A$ 写成 $LL^T$。逐列更新时若 $a_{jj}-\sum_{k<j}l_{jk}^2\le0$，则正定假设失败。用 $\left(\begin{smallmatrix}1&2\\2&1\end{smallmatrix}\right)$ 时第二个枢轴为 $1-4<0$；应报告负特征方向和输入来源，不能加任意 jitter 后假装原矩阵有效。
\end{enumerate}\MFQInterviewFollowup{GB-3-16：列近似相关时 QR 为什么比正规方程稳定？GB-3-17 与 GB-3-18：零特征值对风险模型意味什么？GB-3-19：何时允许收缩而不是盲目加 jitter？}}

\providecommand{\MFQMappedUpperChSevenQuestions}{
\subsection*{笔试推导（来源映射替换）}
\begin{enumerate}[nosep,leftmargin=*]
\item \MFQInterviewMeta{笔试推导}{10 分钟}{中}{置换对称}{GB-4-03 Drunk passenger}第一位旅客随机占一个座位，后续旅客若自己的座位被占就随机选空座。证明最后一位坐到自己座位的概率，并指出递推中的最小状态。
\item \MFQInterviewMeta{笔试推导}{10 分钟}{中}{停止时与对称随机游走}{GB-4-01 Coin toss game I}公平硬币重复抛掷，正面累计领先反面两次时停止。写出吸收随机游走模型，并求从零开始先到 $+2$ 而非 $-2$ 的概率。
\item \MFQInterviewMeta{笔试推导}{12 分钟}{中}{组合计数}{GB-4-05 Poker hands}从标准扑克牌抽五张，推导恰为两对的概率；分母、点数组合、两对花色和第五张牌必须分别计数。
\item \MFQInterviewMeta{笔试推导}{10 分钟}{中}{卷积与独立性}{GB-4-37 Sum of random variables}独立 $X,Y\sim U(0,1)$，推导 $S=X+Y$ 的分段密度，并指出若只知道边际分布而不知道联合分布，为何答案不唯一。
\end{enumerate}
\subsection*{数值编程（来源映射替换）}
\begin{enumerate}[nosep,leftmargin=*]
\item \MFQInterviewMeta{数值编程}{12 分钟}{中}{指示变量与期望}{GB-4-38 Coupon collection}从 $m$ 种等概率标签中独立抽取，推导集齐全部标签的期望抽数，并写程序对比精确值与 Monte Carlo 结果。
\item \MFQInterviewMeta{数值编程}{12 分钟}{中}{碰撞概率与数值稳定}{GB-4-10 Birthday problem II}推导 $n$ 人生日至少一次重复的精确概率，并实现对数域计算；比较常用指数近似在 $n=23,50,100$ 时的误差。
\item \MFQInterviewMeta{数值编程}{12 分钟}{中}{Poisson 独立增量}{GB-4-31 Poisson process property}模拟强度 $\lambda$ 的 Poisson 过程，验证不相交区间计数近似独立且均值等于区间长度乘 $\lambda$；说明有限模拟不能证明独立性。
\end{enumerate}
\subsection*{研究判断（来源映射替换）}
\begin{enumerate}[nosep,leftmargin=*]
\item \MFQInterviewMeta{研究判断}{10 分钟}{中}{几何概率与分布假设}{GB-4-29 Meeting probability}两人在一小时内独立均匀到达、各等待十二分钟时计算相遇概率；再审计直接将该结果用于日内非均匀订单到达的报告。
\item \MFQInterviewMeta{研究判断}{12 分钟}{难}{相关违约与尾部依赖}{GB-4-39 Joint default}两个主体一年违约概率均为 $p$，相关系数给定。判断这些信息是否足以唯一确定联合违约概率；写出 Bernoulli 情形的可行界，并说明相关系数为何不能替代 copula 或条件违约模型。
\end{enumerate}}
\providecommand{\MFQMappedUpperChSevenHints}{\par\textbf{来源映射题提示：}旅客过程只追踪两个关键座位；随机游走利用反射对称；扑克牌按点数与花色分层计数；卷积看单位正方形截线长度；生日题使用对数乘积；Poisson 题区分理论性质与模拟证据；联合违约先写 Fr\'echet 界。}
\providecommand{\MFQMappedUpperChSevenSolutions}{
\section*{来源映射替换题}
\begin{enumerate}[leftmargin=*]
\item \textbf{GB-4-03。}中间被迫选择时，尚未解决的关键座位只有第一人和最后一人；先选到前者则链终止且最后一人成功，先选到后者则失败。两者在对称候选中等可能，故概率 $1/2$。逐位枚举不是必要状态。
\item \textbf{GB-4-01。}状态是领先数 $D_n\in\{-2,-1,0,1,2\}$，边界 $\pm2$ 吸收。令 $h(d)$ 为先到 $+2$ 的概率，则 $h(d)=[h(d-1)+h(d+1)]/2$、$h(-2)=0,h(2)=1$，线性解给 $h(0)=1/2$。若硬币有偏，差分方程的系数必须改变。
\item \textbf{GB-4-05。}总手数 $\binom{52}{5}$。选两种对子点数 $\binom{13}{2}$，各选两张花色 $\binom42^2$，第五张从剩余十一种点数与四种花色选 $11\cdot4$，故概率为 $\binom{13}{2}\binom42^2(44)/\binom{52}{5}$。第五张不能与对子同点，否则牌型改变。
\item \textbf{GB-4-37。}$f_S(s)=\int f_X(x)f_Y(s-x)\,dx$，积分区间长度在 $0\le s\le1$ 为 $s$，在 $1<s\le2$ 为 $2-s$，其余为零。卷积使用独立性；相同均匀边际可以通过不同 copula 产生不同和分布。
\item \textbf{GB-4-29。}令到达时刻缩放到 $[0,1]^2$，等待十二分钟即 $0.2$。不相遇区域是对角线两侧两个边长 $0.8$ 的直角三角形，总面积 $0.8^2$，故相遇概率 $1-0.64=0.36$。
\item \textbf{GB-4-38。}已有 $k$ 种时，新标签概率为 $(m-k)/m$，等待期望 $m/(m-k)$。总期望为 $m\sum_{j=1}^m1/j=mH_m\sim m\log m+\gamma m$。线性期望不要求各阶段等待独立。
\item \textbf{GB-4-10。}无重复概率为 $\prod_{k=0}^{n-1}(365-k)/365$，故重复概率为一减该乘积；实现时计算 $\sum_k\log(1-k/365)$ 再取指数。近似 $1-\exp[-n(n-1)/(730)]$ 忽略高阶碰撞项，$n$ 增大时必须报告误差。
\item \textbf{GB-4-31。}定义要求 $N(t)-N(s)\sim\mathrm{Poisson}(\lambda(t-s))$，且不相交区间增量独立。模拟可比较样本均值、方差和交叉协方差是否接近理论值，但“协方差接近零”不证明独立，且罕见尾部需要更大样本。
\item \textbf{GB-4-39。}联合概率 $q=P(D_1=D_2=1)$ 满足 $\max(0,2p-1)\le q\le p$；Bernoulli 相关系数若精确给定且两边际相同，可由 $q=p^2+\rho p(1-p)$ 恢复，但还必须落在可行界内。真实信用组合还需期限结构、共同因子和尾部依赖，单一相关系数不能跨状态外推。
\end{enumerate}\MFQInterviewFollowup{GB-4-01：有偏硬币时命中概率如何变化？GB-4-03：首位旅客按非均匀分布选座时结论如何改？GB-4-05：允许 joker 后计数如何重做？GB-4-10：闰日和非均匀出生率会破坏哪一步？GB-4-29：到达分布非均匀时如何积分？GB-4-31：非齐次强度怎样改变增量分布？GB-4-37：怎样用特征函数复核卷积？GB-4-38：非等概标签的瓶颈是什么？GB-4-39：怎样用单因子模型产生尾部相关？}}

\providecommand{\MFQMappedUpperChEightQuestions}{
\subsection*{口述概念（来源映射替换）}
\begin{enumerate}[nosep,leftmargin=*]
\item \MFQInterviewMeta{口述概念}{8 分钟}{中}{随机化提纯}{GB-4-16 Fair coin from an unfair coin}一枚硬币正面概率未知但每次独立且不变。构造无偏比特并计算每个输出比特的期望抛掷次数。
\item \MFQInterviewMeta{口述概念}{8 分钟}{中}{可辨识性与似然}{GB-4-15 Unfair coin}只观察一枚硬币的有限次结果，解释为什么不能“证明”它公平；写出估计正面概率的似然、置信区间与顺序检验所需的停止规则。
\end{enumerate}
\subsection*{笔试推导（来源映射替换）}
\begin{enumerate}[nosep,leftmargin=*]
\item \MFQInterviewMeta{笔试推导}{10 分钟}{中}{条件概率与无放回抽样}{GB-4-25 Aces}从洗匀牌堆依次翻牌，已知前两张至少有一张 A，求两张都是 A 的概率；分别比较“至少一张”与“指定第一张是 A”的条件事件。
\end{enumerate}
\subsection*{数值编程（来源映射替换）}
\begin{enumerate}[nosep,leftmargin=*]
\item \MFQInterviewMeta{数值编程}{12 分钟}{中}{Poisson 流量与暴露时间}{GB-4-28 Cars on road}假设车辆按强度 $\lambda$ 的 Poisson 过程到达长度为 $L$ 的路段、速度固定为 $v$。推导任意时刻路段车辆数分布并模拟核验；再指出拥堵使速度依赖密度时为何模型失效。
\end{enumerate}
\subsection*{研究判断（来源映射替换）}
\begin{enumerate}[nosep,leftmargin=*]
\item \MFQInterviewMeta{研究判断}{8 分钟}{中}{条件概率与信息更新}{GB-4-20 Monty Hall}主持人知道奖品位置且总会打开一扇空门时证明换门胜率；再审计“主持人随机开门也可直接套用 $2/3$”的结论。
\item \MFQInterviewMeta{研究判断}{10 分钟}{中}{条件生存与协议差异}{GB-4-24 Russian roulette}左轮有六个弹膛和相邻两发子弹。已知第一次扣动为空，比较“直接再扣”与“重新随机旋转后再扣”的风险；明确枪膛是否顺序前进这一机制条件。
\end{enumerate}}
\providecommand{\MFQMappedUpperChEightHints}{\par\textbf{来源映射题提示：}无偏提纯只接受 HT/TH；有限样本不能证明点假设；A 题先写清条件事件；路段车辆数利用到达率乘暴露时间；Monty Hall 与轮盘题都要把主持或机械协议写进似然。}
\providecommand{\MFQMappedUpperChEightSolutions}{
\section*{来源映射替换题}
\begin{enumerate}[leftmargin=*]
\item \textbf{GB-4-16。}成对抛掷，HT 输出 1、TH 输出 0，HH/TT 丢弃。两种接受序列概率同为 $p(1-p)$，故无偏；每对接受概率 $2p(1-p)$，期望对数为其倒数，期望抛数 $1/[p(1-p)]$。独立且固定 $p$ 是必要条件。
\item \textbf{GB-4-15。}若观察 $h$ 次正面、$t$ 次反面，似然为 $L(p)=p^h(1-p)^t$，极大似然估计 $\hat p=h/(h+t)$。任何有限序列在所有 $0<p<1$ 下都有正概率，故不能逻辑证明 $p=1/2$；只能给区间或在预先规定的显著性与停止规则下检验。边看边停会改变错误率。
\item \textbf{GB-4-25。}两张无序抽样中，条件事件“至少一张 A”有 $\binom{52}{2}-\binom{48}{2}=198$ 种，二者均 A 有 $\binom42=6$ 种，故概率 $6/198=1/33$。若指定第一张是 A，则第二张为 A 的概率为 $3/51=1/17$；两个条件事件不能混用。
\item \textbf{GB-4-28。}\par 一辆车停留时间为 $L/v$。稳态下，路段车辆数等于过去 $L/v$ 时间内的到达数，故服从 $\mathrm{Poisson}(\lambda L/v)$。模拟应从稳态开始或丢弃预热期。若速度随密度变化，停留时间与车辆数相互依赖，简单的无限服务器结论失效。
\item \textbf{GB-4-20。}初选中奖概率 $1/3$，未选两门合计 $2/3$；知情主持人必排除其中空门，所以换到剩余门继承 $2/3$。若主持人随机开门且可能揭奖，观察到空门的似然随奖品位置不同，必须用 Bayes 公式重新条件化。
\item \textbf{GB-4-24。}两发相邻时共有四个空膛位置。已知当前为空后，这四种位置等可能，其中只有一个位置的下一膛是子弹，故直接再扣风险为 $1/4$；重新均匀旋转后风险回到 $2/6=1/3$，所以不旋转更安全。若弹膛不是确定前进一步或初始旋转不均匀，需重写条件样本空间。
\end{enumerate}\MFQInterviewFollowup{GB-4-15：序贯检验如何控制可选停止？GB-4-16：偏置 $p$ 随时间漂移会破坏哪一步？GB-4-20：主持人策略未知时需要估计什么似然？GB-4-24：三发相邻时结论如何变化？GB-4-25：若抽样信息来自“随机告诉你一张 A”会怎样？GB-4-28：随机速度如何改变占用分布？}}

\providecommand{\MFQMappedUpperChNineQuestions}{
\subsection*{研究判断（来源映射替换）}
\begin{enumerate}[nosep,leftmargin=*]
\item \MFQInterviewMeta{研究判断}{12 分钟}{中}{次序统计量与假设边界}{GB-4-40 Expected maximum and minimum}若 $X_1,\ldots,X_n$ 独立均匀于 $[0,1]$，推导最大值与最小值的分布和期望；再说明样本非独立时为何不能直接复用。
\item \MFQInterviewMeta{笔试推导}{10 分钟}{中}{正态矩与递推}{GB-4-32 Normal moments}对 $Z\sim N(0,1)$，用分部积分证明 $E[Z^{2k}]=(2k-1)E[Z^{2k-2}]$，并解释奇数阶矩为何为零。
\item \MFQInterviewMeta{研究判断}{12 分钟}{难}{相关界与可达性}{GB-4-41 Correlation maximum and minimum}给定随机变量 $X$ 的边际分布，讨论 $\mathrm{Corr}(X,g(X))$ 的上下界；说明 Cauchy--Schwarz 的 $[-1,1]$ 只是普适界，何时才能取到等号。
\end{enumerate}}
\providecommand{\MFQMappedUpperChNineHints}{\par\textbf{来源映射题提示：}$P(M\le x)=x^n$；正态矩把 $z\phi(z)$ 写成 $-\phi'(z)$；相关取等要求中心化变量几乎处处线性相关。}
\providecommand{\MFQMappedUpperChNineSolutions}{
\section*{来源映射替换题}
\begin{enumerate}[leftmargin=*]
\item \textbf{GB-4-40。}$F_M(x)=x^n$，故密度 $nx^{n-1}$、$E[M]=n/(n+1)$。$P(m>x)=(1-x)^n$，所以 $E[m]=\int_0^1(1-x)^n dx=1/(n+1)$。仿射变换给 $a+(b-a)$ 乘相应结果；独立同分布是推导前提。
\item \textbf{GB-4-32。}标准正态密度满足 $\phi'(z)=-z\phi(z)$。分部积分得 $E[Z^{2k}]=(2k-1)E[Z^{2k-2}]$，从 $E[1]=1$ 递推为 $(2k-1)!!$；奇次幂乘偶密度是奇函数，积分为零。非中心正态应先标准化或用矩母函数。
\item \textbf{GB-4-41。}Cauchy--Schwarz 给 $|\mathrm{Corr}(X,g(X))|\le1$；等号当且仅当 $g(X)-Eg(X)=a(X-EX)$ 几乎处处。若 $g$ 被限制为非线性、单调或离散决策类，可达界通常更窄，必须在允许函数类与边际分布下重新优化，不能把普适界当成可实现值。
\end{enumerate}\MFQInterviewFollowup{GB-4-32：如何由矩母函数复核偶数矩？GB-4-40：样本有厚尾时，极值渐近要改用哪类条件？GB-4-41：限制 $g$ 为指标函数后如何求最优阈值？}}

\providecommand{\MFQMappedUpperChElevenQuestions}{
\subsection*{研究判断（来源映射替换）}
\begin{enumerate}[nosep,leftmargin=*]
\item \MFQInterviewMeta{研究判断}{15 分钟}{难}{模式等待与 Markov 状态}{GB-5-03 Coin triplets}反复抛公平硬币，比较模式 HHT 与 HTH 的首次出现等待时间；构造前缀状态，并审计“长度相同所以等待时间相同”的结论。
\item \MFQInterviewMeta{笔试推导}{12 分钟}{中}{吸收概率与差分方程}{GB-5-01 Gambler's ruin II}资本在 $0$ 与 $N$ 之间作步长 $\pm1$ 的随机游走，向上概率为 $p$。推导从 $i$ 出发先到 $N$ 的概率，并单独处理 $p=1/2$。
\item \MFQInterviewMeta{数值编程}{12 分钟}{中}{首达时间与状态截断}{GB-5-05 Drunk man}随机游走者从整数点 $i$ 出发，在 $0,N$ 吸收。建立期望吸收时间的线性方程并用三对角求解器计算；与 Monte Carlo 比较时报告置信误差。
\item \MFQInterviewMeta{研究判断}{12 分钟}{中}{排队状态与独立假设}{GB-5-07 Ticket line}售票队列中到达与服务均被建模为 Poisson/指数。写出出生--死亡状态，并说明只匹配均值为何不足以支持等待时间结论。
\item \MFQInterviewMeta{笔试推导}{12 分钟}{中}{动态规划与系列赛}{GB-5-11 World series}两队进行先赢四场的系列赛，每场 A 胜率固定为 $p$ 且独立。写出从比分 $(a,b)$ 出发 A 最终夺冠的递推和边界，并给出 $p=1/2$ 时开赛前概率。
\end{enumerate}}
\providecommand{\MFQMappedUpperChElevenHints}{\par\textbf{来源映射题提示：}模式状态取仍可延伸目标的最长后缀；破产概率解二阶差分；吸收时间右端多一个常数项；排队先核验 Markov/无记忆假设；系列赛以当前胜场为状态。}
\providecommand{\MFQMappedUpperChElevenSolutions}{
\section*{来源映射替换题}
\begin{enumerate}[leftmargin=*]
\item \textbf{GB-5-03。}对每个模式建立前缀自动机，状态为已匹配长度 $0,1,2,3$，吸收态为 3；对 $i<3$ 写 $E_i=1+(E_{T(i,H)}+E_{T(i,T)})/2$。求解得到 HHT 的期望等待 $8$，HTH 的期望等待 $10$。差异来自 HTH 失败后保留的前缀结构，不能用“长度都是三”断言相同。
\item \textbf{GB-5-01。}令 $h_i=P_i(\tau_N<\tau_0)$，则 $h_i=ph_{i+1}+(1-p)h_{i-1}$，$h_0=0,h_N=1$。若 $p\ne1/2$，$h_i=[1-((1-p)/p)^i]/[1-((1-p)/p)^N]$；公平时取极限得 $i/N$。边界和步长改变时不能套用闭式。
\item \textbf{GB-5-05。}期望吸收时间满足 $e_i=1+pe_{i+1}+(1-p)e_{i-1}$、$e_0=e_N=0$，形成三对角线性系统。公平时闭式为 $e_i=i(N-i)$。模拟均值应附标准误或区间；厚尾的停止时间会使有限样本收敛看起来很慢。
\item \textbf{GB-5-07。}状态可取系统人数 $n$，转移率为 $n\to n+1$ 的 $\lambda$ 与 $n\to n-1$ 的 $\mu$（$n>0$）。稳态要求 $\rho=\lambda/\mu<1$。若到达有聚集、服务时长非指数或服务能力随时段变化，均值相同也会产生不同等待尾部。
\item \textbf{GB-5-11。}令 $V(a,b)$ 为 A 夺冠概率，边界 $V(4,b)=1,V(a,4)=0$，递推 $V(a,b)=pV(a+1,b)+(1-p)V(a,b+1)$。$p=1/2$ 且开赛前双方对称，故 $V(0,0)=1/2$。若胜率随主客场或比分变化，状态必须扩充。
\end{enumerate}\MFQInterviewFollowup{GB-5-01：加入每步不同赌注后状态如何改变？GB-5-03：如何由前缀自动机直接生成转移矩阵并比较模式竞赛胜率？GB-5-05：如何避免无限循环模拟拖垮测试？GB-5-07：怎样从数据检验指数服务假设？GB-5-11：七局赛改成动态胜率时如何后向计算？}}

\providecommand{\MFQMappedUpperThirteenQuestions}{
\subsection*{研究判断（来源映射替换）}
\begin{enumerate}[nosep,leftmargin=*]
\item \MFQInterviewMeta{研究判断}{12 分钟}{中}{约束与乘子解释}{GB-3-10 Lagrange multipliers}最小化 $w^T\Sigma w$，约束 $\mathbf1^Tw=1$ 与 $\mu^Tw=r_0$。写出 KKT 线性系统，并审计“求解器成功即表示乘子和经济约束都可识别”的结论。
\end{enumerate}}
\providecommand{\MFQMappedUpperThirteenHints}{\par\textbf{来源映射题提示：}把两个等式约束矩阵记为 $A^Tw=b$，联立 $2\Sigma w+A\lambda=0$。}
\providecommand{\MFQMappedUpperThirteenSolutions}{
\section*{来源映射替换题}
\begin{enumerate}[leftmargin=*]
\item \textbf{GB-3-10。}KKT 系统为 $\left(\begin{smallmatrix}2\Sigma&A\\A^T&0\end{smallmatrix}\right)(w,\lambda)^T=(0,b)^T$，其中 $A=[\mathbf1,\mu]$、$b=(1,r_0)^T$。若 $A$ 列相关或与 $\Sigma$ 的零空间形成退化，约束资格或严格凸性失效；程序返回一个解不证明经济约束可识别。
\end{enumerate}\MFQInterviewFollowup{GB-3-10：加入不等式限制后哪些 KKT 条件需要新增？}}

\providecommand{\MFQMappedUpperFourteenQuestions}{
\subsection*{数值编程与研究判断（来源映射替换）}
\begin{enumerate}[nosep,leftmargin=*]
\item \MFQInterviewMeta{数值编程}{12 分钟}{中}{Newton 收敛域}{GB-3-09 Newton's method}用 Newton 法求 $x^3-2x-2=0$，记录不同初值的迭代；给出停止准则，并构造一个进入循环或远离根的初值。
\item \MFQInterviewMeta{笔试推导}{12 分钟}{中}{有限差分误差}{GB-7-15 Finite difference}从 Taylor 展开推导中心差分的一阶导数公式与 $O(h^2)$ 截断误差，再解释 $h$ 过小时舍入误差为何上升。
\item \MFQInterviewMeta{数值编程}{15 分钟}{中}{Monte Carlo 误差预算}{GB-7-14 Monte Carlo}用 Monte Carlo 估计 $E[e^Z]$（$Z\sim N(0,1)$），给出独立理论值、标准误和置信区间；再用 antithetic variates 比较方差而不是只比较单次误差。
\end{enumerate}}
\providecommand{\MFQMappedUpperFourteenHints}{\par\textbf{来源映射题提示：}Newton 需同时检查残差与步长；中心差分把两侧 Taylor 展开相减；Monte Carlo 先用矩母函数给独立 oracle，再比较配对样本的方差。}
\providecommand{\MFQMappedUpperFourteenSolutions}{
\section*{来源映射替换题}
\begin{enumerate}[leftmargin=*]
\item \textbf{GB-3-09。}$x_{k+1}=x_k-(x_k^3-2x_k-2)/(3x_k^2-2)$。停止应同时要求 $|f(x_k)|$ 和相对步长足够小，并设最大迭代；导数接近零或初值落在非吸引区域时可能跳远/循环。面试追问不应以某个库的 success 标志替代理论收敛域。
\item \textbf{GB-7-15。}两式相减得 $[f(x+h)-f(x-h)]/(2h)=f'(x)+h^2f^{(3)}(\xi)/6$。截断误差随 $h^2$ 降低，但两近数相减的浮点绝对误差约为 $\epsilon|f|$，除以 $h$ 后放大为 $O(\epsilon/h)$，因此存在最优中间尺度。
\item \textbf{GB-7-14。}正态矩母函数给 $E[e^Z]=e^{1/2}$，且 $\mathrm{Var}(e^Z)=e^2-e$。独立样本均值的标准误为样本标准差除以 $\sqrt n$。反变量配对使用 $[e^Z+e^{-Z}]/2$；应在相同随机预算下比较配对估计量的样本方差和区间宽度，而不是挑一次更接近真值的结果。
\end{enumerate}\MFQInterviewFollowup{GB-3-09：如何用阻尼或信赖域避免跳远？GB-7-14：控制变量与重要性抽样分别需要什么额外结构？GB-7-15：如何从截断与舍入两项估计最佳 $h$？}}

\providecommand{\MFQMappedDerivativesStochasticQuestions}{
\section*{研究判断（来源映射替换）}
\begin{enumerate}[leftmargin=*]
\item \MFQInterviewMeta{研究判断}{15 分钟}{难}{停止时与反射原理}{GB-5-15 Stopping and first passage}对标准 Brownian 运动推导首次越过 $a>0$ 在 $T$ 前发生的概率，并审计直接将无界停止时代入鞅等式的证明。
\end{enumerate}}
\providecommand{\MFQMappedDerivativesStochasticHints}{\par\textbf{来源映射题提示：}用反射原理把越界后终点低于 $a$ 的路径映到终点高于 $a$。}
\providecommand{\MFQMappedDerivativesStochasticSolutions}{\section*{来源映射替换题}\begin{enumerate}[leftmargin=*]\item \textbf{GB-5-15。}反射原理给 $P(\tau_a\le T)=2P(W_T\ge a)=2[1-\Phi(a/\sqrt T)]$。若改用可选停止证明，必须先对停止时截断并验证一致可积或相应有界条件；直接把无界停止时代入鞅等式是常见错误。\end{enumerate}\MFQInterviewFollowup{GB-5-15：若边界改为随时间变化的 $a(t)$，反射证明的哪一步不再直接成立？}}

\providecommand{\MFQMappedDerivativesNumericsQuestions}{
\section*{研究判断（来源映射替换）}
\begin{enumerate}[leftmargin=*]
\item \MFQInterviewMeta{研究判断}{10 分钟}{中}{无套利复制}{GB-6-02 Put-call parity}在含连续股息率 $q$ 的市场推导欧式看涨--看跌平价，并审计一个忽略 bid/ask、融资、借券与结算时点便声称可套利的报告。
\end{enumerate}}
\providecommand{\MFQMappedDerivativesNumericsHints}{\par\textbf{来源映射题提示：}比较 $C+Ke^{-rT}$ 与 $P+S_0e^{-qT}$ 的终值。}
\providecommand{\MFQMappedDerivativesNumericsSolutions}{\section*{来源映射替换题}\begin{enumerate}[leftmargin=*]\item \textbf{GB-6-02。}两组合到期均为 $\max(S_T,K)$，故 $C-P=S_0e^{-qT}-Ke^{-rT}$。若左侧过高则卖出昂贵组合、买入便宜组合；真实可执行性还受 bid/ask、融资、股息预测、借券、行权与结算时点约束。\end{enumerate}\MFQInterviewFollowup{GB-6-02：美式期权或随机股息下，复制组合的哪个等式需要重写？}}

\providecommand{\MFQMappedPortfolioOptimizationQuestions}{
\section*{来源映射替换题}
\begin{enumerate}[leftmargin=*]
\item \MFQInterviewMeta{笔试推导}{15 分钟}{难}{均值方差与估计误差}{GB-6-14 Portfolio optimization}推导预算约束下的最小方差权重；随后说明均值方差解为何会放大预期收益估计误差，并提出至少两种稳健化方案。
\end{enumerate}}
\providecommand{\MFQMappedPortfolioOptimizationHints}{\par\textbf{来源映射题提示：}先解 $\min w^T\Sigma w$、$\mathbf1^Tw=1$；含均值时观察 $\Sigma^{-1}\hat\mu$。}
\providecommand{\MFQMappedPortfolioOptimizationSolutions}{\section*{来源映射替换题}\begin{enumerate}[leftmargin=*]\item \textbf{GB-6-14。}$w^*=\Sigma^{-1}\mathbf1/(\mathbf1^T\Sigma^{-1}\mathbf1)$。含预期收益的解依赖 $\Sigma^{-1}\hat\mu$，小特征方向和噪声均值会被放大。可用协方差收缩、均值不确定集、权重/换手约束或 Bayesian shrinkage，并必须重新比较样本外风险和成本。\end{enumerate}\MFQInterviewFollowup{GB-6-14：禁止卖空后闭式解为何失效，应怎样识别活跃约束？}}

\providecommand{\MFQMappedPortfolioTailQuestions}{
\section*{来源映射替换题}
\begin{enumerate}[leftmargin=*]
\item \MFQInterviewMeta{口述概念}{10 分钟}{中}{VaR 的性质与反例}{GB-6-15 Value at Risk}解释历史 VaR、参数 VaR 与 ES 的估计对象；构造一个 VaR 不次可加的简单离散损失例，并说明这对风险限额意味着什么。
\end{enumerate}}
\providecommand{\MFQMappedPortfolioTailHints}{\par\textbf{来源映射题提示：}让两个头寸各自以小概率损失、且损失事件互斥；单项 VaR 可为零而组合 VaR 为正。}
\providecommand{\MFQMappedPortfolioTailSolutions}{\section*{来源映射替换题}\begin{enumerate}[leftmargin=*]\item \textbf{GB-6-15。}历史法取经验损失分位数，参数法在分布假设下求分位数，ES 是越过分位点的尾部平均。令两个头寸各以 $4\%$ 概率损失 1，事件互斥；在 $95\%$ 水平各自 VaR 为 0，而组合有 $8\%$ 概率损失 1，VaR 为 1，故不次可加。限额系统不能把分散收益当作由 VaR 自动保证。\end{enumerate}\MFQInterviewFollowup{GB-6-15：把置信度提到 $99\%$ 时，需要多少尾部观测才能支撑可辩护的 ES？}}

\MFQMappedDerivativesStochasticHints

\chapter{定价、校准与对冲提示}
\begin{enumerate}
  \item 按时间、空间、边界、采样、求根和模型错设分类。
  \item 数一数自由参数：逐点反演每个节点有一个参数。
  \item 从 $V_\tau=LV$ 写 $I-\Delta tL_h$。
  \item 比较相邻割线斜率，而非默认等距。
  \item 一次只改变一个网格维度。
  \item 同时报告参数误差与重定价误差。
  \item 检查分红、利率、远期与贴现归一化。
  \item par yield 还不是完整零息曲线。
\end{enumerate}

% Source-mapped interview replacements.  The questions are independent Chinese
% adaptations; only the short source titles are retained for auditability.
\providecommand{\MFQInterviewFollowup}[1]{\par\noindent\textbf{面试官追问：}#1}

\providecommand{\MFQMappedUpperChFiveQuestions}{
\subsection*{笔试推导（来源映射替换）}
\begin{enumerate}[nosep,leftmargin=*]
\item \MFQInterviewMeta{笔试推导}{12 分钟}{中}{一阶条件与边界比较}{GB-3-02 Maximum and minimum}给定 $f(x)=x^4-4x^2+cx$，写出随参数 $c$ 变化的驻点判定流程；说明只解 $f'(x)=0$ 为什么不足以确定全局最小值。
\end{enumerate}
\subsection*{研究判断（来源映射替换）}
\begin{enumerate}[nosep,leftmargin=*]
\item \MFQInterviewMeta{研究判断}{12 分钟}{中}{Taylor 余项与误差控制}{GB-3-08 Taylor series}在 $x=0$ 展开 $\log(1+x)$ 到三阶，用余项给出 $|x|\le0.2$ 的统一误差界；据此审计“在所有收益率上都可安全替代对数收益”的报告。
\end{enumerate}}
\providecommand{\MFQMappedUpperChFiveHints}{\par\textbf{来源映射题提示：}先比较全部驻点、边界与无穷远行为；Taylor 题用四阶导数在区间上的上界控制 Lagrange 余项。}
\providecommand{\MFQMappedUpperChFiveSolutions}{
\section*{来源映射替换题}
\begin{enumerate}[leftmargin=*]
\item \textbf{GB-3-02。}$f'(x)=4x^3-8x+c$ 只给候选点；还需用 $f''(x)=12x^2-8$ 分类，并比较所有局部极小值。因首项系数为正，$f(x)\to\infty$，故全局最小一定在驻点中；若定义域有边界还要加入端点。参数追问应讨论三次方程重根处的分岔，而不能只给数值根。
\item \textbf{GB-3-08。}$\log(1+x)=x-x^2/2+x^3/3+R_4$，且 $f^{(4)}(u)=-6/(1+u)^4$。在 $|x|\le0.2$ 上 $|R_4|\le |x|^4/[4(0.8)^4]$。下一项为负，因此正 $x$ 附近三阶截断高估对数收益；负 $x$ 时应直接依余项而不是机械套交错级数。
\end{enumerate}\MFQInterviewFollowup{GB-3-02：定义域变成紧集时还要加哪些候选点？GB-3-08：如何用同一余项界选取收益率阈值？}}

\providecommand{\MFQMappedUpperChSixQuestions}{
\subsection*{笔试推导（来源映射替换）}
\begin{enumerate}[nosep,leftmargin=*]
\item \MFQInterviewMeta{笔试推导}{12 分钟}{中}{QR 与最小二乘}{GB-3-16 QR decomposition}对满列秩 $X=QR$，从残差正交性推导 $R\hat\beta=Q^Ty$，并解释为何不必形成 $X^TX$。
\end{enumerate}
\subsection*{口述概念（来源映射替换）}
\begin{enumerate}[nosep,leftmargin=*]
\item \MFQInterviewMeta{口述概念}{10 分钟}{中}{谱与经济方向}{GB-3-17 Determinant, eigenvalue and eigenvector}解释协方差矩阵的特征向量、特征值和行列式分别表达什么；零行列式为什么不等于“所有风险都为零”？
\end{enumerate}
\subsection*{研究判断（来源映射替换）}
\begin{enumerate}[nosep,leftmargin=*]
\item \MFQInterviewMeta{研究判断}{10 分钟}{中}{半正定与模型边界}{GB-3-18 Positive semidefinite and positive definite matrices}证明 $B^TB$ 半正定，再审计“将任意估计矩阵换成 $B^TB$ 就保留了原风险模型”的报告。
\end{enumerate}
\subsection*{数值编程（来源映射替换）}
\begin{enumerate}[nosep,leftmargin=*]
\item \MFQInterviewMeta{数值编程}{12 分钟}{中}{LU/Cholesky 分解与失败诊断}{GB-3-19 LU decomposition and Cholesky decomposition}对比 LU 与 Cholesky 的前提和计算结构；实现带对角门禁的 Cholesky，并用一个对称但不正定矩阵证明失败是模型证据而非“库坏了”。
\end{enumerate}}
\providecommand{\MFQMappedUpperChSixHints}{\par\textbf{来源映射题提示：}利用 $Q^TQ=I$；谱题区分零方向与全部方向；$x^TB^TBx=\|Bx\|^2$；Cholesky 的根号参数必须为正。}
\providecommand{\MFQMappedUpperChSixSolutions}{
\section*{来源映射替换题}
\begin{enumerate}[leftmargin=*]
\item \textbf{GB-3-16。}$\|y-QR\beta\|^2=\|Q^Ty-R\beta\|^2+\|(I-QQ^T)y\|^2$，故最小化第一项得到 $R\hat\beta=Q^Ty$。形成正规方程会把奇异值平方，使条件数平方。
\item \textbf{GB-3-17。}特征向量是风险二次型的正交主方向，特征值是这些方向的方差；行列式是特征值乘积，反映总体积尺度。行列式为零只说明至少一个线性组合方差为零，其他方向仍可有正风险。
\item \textbf{GB-3-18。}$x^TB^TBx=\|Bx\|_2^2\ge0$，所以半正定。其正定当且仅当 $Bx=0$ 只有零解，即 $B$ 满列秩。
\item \textbf{GB-3-19。}LU 面向一般方阵且通常需选主元；Cholesky 利用对称正定结构，将 $A$ 写成 $LL^T$。逐列更新时若 $a_{jj}-\sum_{k<j}l_{jk}^2\le0$，则正定假设失败。用 $\left(\begin{smallmatrix}1&2\\2&1\end{smallmatrix}\right)$ 时第二个枢轴为 $1-4<0$；应报告负特征方向和输入来源，不能加任意 jitter 后假装原矩阵有效。
\end{enumerate}\MFQInterviewFollowup{GB-3-16：列近似相关时 QR 为什么比正规方程稳定？GB-3-17 与 GB-3-18：零特征值对风险模型意味什么？GB-3-19：何时允许收缩而不是盲目加 jitter？}}

\providecommand{\MFQMappedUpperChSevenQuestions}{
\subsection*{笔试推导（来源映射替换）}
\begin{enumerate}[nosep,leftmargin=*]
\item \MFQInterviewMeta{笔试推导}{10 分钟}{中}{置换对称}{GB-4-03 Drunk passenger}第一位旅客随机占一个座位，后续旅客若自己的座位被占就随机选空座。证明最后一位坐到自己座位的概率，并指出递推中的最小状态。
\item \MFQInterviewMeta{笔试推导}{10 分钟}{中}{停止时与对称随机游走}{GB-4-01 Coin toss game I}公平硬币重复抛掷，正面累计领先反面两次时停止。写出吸收随机游走模型，并求从零开始先到 $+2$ 而非 $-2$ 的概率。
\item \MFQInterviewMeta{笔试推导}{12 分钟}{中}{组合计数}{GB-4-05 Poker hands}从标准扑克牌抽五张，推导恰为两对的概率；分母、点数组合、两对花色和第五张牌必须分别计数。
\item \MFQInterviewMeta{笔试推导}{10 分钟}{中}{卷积与独立性}{GB-4-37 Sum of random variables}独立 $X,Y\sim U(0,1)$，推导 $S=X+Y$ 的分段密度，并指出若只知道边际分布而不知道联合分布，为何答案不唯一。
\end{enumerate}
\subsection*{数值编程（来源映射替换）}
\begin{enumerate}[nosep,leftmargin=*]
\item \MFQInterviewMeta{数值编程}{12 分钟}{中}{指示变量与期望}{GB-4-38 Coupon collection}从 $m$ 种等概率标签中独立抽取，推导集齐全部标签的期望抽数，并写程序对比精确值与 Monte Carlo 结果。
\item \MFQInterviewMeta{数值编程}{12 分钟}{中}{碰撞概率与数值稳定}{GB-4-10 Birthday problem II}推导 $n$ 人生日至少一次重复的精确概率，并实现对数域计算；比较常用指数近似在 $n=23,50,100$ 时的误差。
\item \MFQInterviewMeta{数值编程}{12 分钟}{中}{Poisson 独立增量}{GB-4-31 Poisson process property}模拟强度 $\lambda$ 的 Poisson 过程，验证不相交区间计数近似独立且均值等于区间长度乘 $\lambda$；说明有限模拟不能证明独立性。
\end{enumerate}
\subsection*{研究判断（来源映射替换）}
\begin{enumerate}[nosep,leftmargin=*]
\item \MFQInterviewMeta{研究判断}{10 分钟}{中}{几何概率与分布假设}{GB-4-29 Meeting probability}两人在一小时内独立均匀到达、各等待十二分钟时计算相遇概率；再审计直接将该结果用于日内非均匀订单到达的报告。
\item \MFQInterviewMeta{研究判断}{12 分钟}{难}{相关违约与尾部依赖}{GB-4-39 Joint default}两个主体一年违约概率均为 $p$，相关系数给定。判断这些信息是否足以唯一确定联合违约概率；写出 Bernoulli 情形的可行界，并说明相关系数为何不能替代 copula 或条件违约模型。
\end{enumerate}}
\providecommand{\MFQMappedUpperChSevenHints}{\par\textbf{来源映射题提示：}旅客过程只追踪两个关键座位；随机游走利用反射对称；扑克牌按点数与花色分层计数；卷积看单位正方形截线长度；生日题使用对数乘积；Poisson 题区分理论性质与模拟证据；联合违约先写 Fr\'echet 界。}
\providecommand{\MFQMappedUpperChSevenSolutions}{
\section*{来源映射替换题}
\begin{enumerate}[leftmargin=*]
\item \textbf{GB-4-03。}中间被迫选择时，尚未解决的关键座位只有第一人和最后一人；先选到前者则链终止且最后一人成功，先选到后者则失败。两者在对称候选中等可能，故概率 $1/2$。逐位枚举不是必要状态。
\item \textbf{GB-4-01。}状态是领先数 $D_n\in\{-2,-1,0,1,2\}$，边界 $\pm2$ 吸收。令 $h(d)$ 为先到 $+2$ 的概率，则 $h(d)=[h(d-1)+h(d+1)]/2$、$h(-2)=0,h(2)=1$，线性解给 $h(0)=1/2$。若硬币有偏，差分方程的系数必须改变。
\item \textbf{GB-4-05。}总手数 $\binom{52}{5}$。选两种对子点数 $\binom{13}{2}$，各选两张花色 $\binom42^2$，第五张从剩余十一种点数与四种花色选 $11\cdot4$，故概率为 $\binom{13}{2}\binom42^2(44)/\binom{52}{5}$。第五张不能与对子同点，否则牌型改变。
\item \textbf{GB-4-37。}$f_S(s)=\int f_X(x)f_Y(s-x)\,dx$，积分区间长度在 $0\le s\le1$ 为 $s$，在 $1<s\le2$ 为 $2-s$，其余为零。卷积使用独立性；相同均匀边际可以通过不同 copula 产生不同和分布。
\item \textbf{GB-4-29。}令到达时刻缩放到 $[0,1]^2$，等待十二分钟即 $0.2$。不相遇区域是对角线两侧两个边长 $0.8$ 的直角三角形，总面积 $0.8^2$，故相遇概率 $1-0.64=0.36$。
\item \textbf{GB-4-38。}已有 $k$ 种时，新标签概率为 $(m-k)/m$，等待期望 $m/(m-k)$。总期望为 $m\sum_{j=1}^m1/j=mH_m\sim m\log m+\gamma m$。线性期望不要求各阶段等待独立。
\item \textbf{GB-4-10。}无重复概率为 $\prod_{k=0}^{n-1}(365-k)/365$，故重复概率为一减该乘积；实现时计算 $\sum_k\log(1-k/365)$ 再取指数。近似 $1-\exp[-n(n-1)/(730)]$ 忽略高阶碰撞项，$n$ 增大时必须报告误差。
\item \textbf{GB-4-31。}定义要求 $N(t)-N(s)\sim\mathrm{Poisson}(\lambda(t-s))$，且不相交区间增量独立。模拟可比较样本均值、方差和交叉协方差是否接近理论值，但“协方差接近零”不证明独立，且罕见尾部需要更大样本。
\item \textbf{GB-4-39。}联合概率 $q=P(D_1=D_2=1)$ 满足 $\max(0,2p-1)\le q\le p$；Bernoulli 相关系数若精确给定且两边际相同，可由 $q=p^2+\rho p(1-p)$ 恢复，但还必须落在可行界内。真实信用组合还需期限结构、共同因子和尾部依赖，单一相关系数不能跨状态外推。
\end{enumerate}\MFQInterviewFollowup{GB-4-01：有偏硬币时命中概率如何变化？GB-4-03：首位旅客按非均匀分布选座时结论如何改？GB-4-05：允许 joker 后计数如何重做？GB-4-10：闰日和非均匀出生率会破坏哪一步？GB-4-29：到达分布非均匀时如何积分？GB-4-31：非齐次强度怎样改变增量分布？GB-4-37：怎样用特征函数复核卷积？GB-4-38：非等概标签的瓶颈是什么？GB-4-39：怎样用单因子模型产生尾部相关？}}

\providecommand{\MFQMappedUpperChEightQuestions}{
\subsection*{口述概念（来源映射替换）}
\begin{enumerate}[nosep,leftmargin=*]
\item \MFQInterviewMeta{口述概念}{8 分钟}{中}{随机化提纯}{GB-4-16 Fair coin from an unfair coin}一枚硬币正面概率未知但每次独立且不变。构造无偏比特并计算每个输出比特的期望抛掷次数。
\item \MFQInterviewMeta{口述概念}{8 分钟}{中}{可辨识性与似然}{GB-4-15 Unfair coin}只观察一枚硬币的有限次结果，解释为什么不能“证明”它公平；写出估计正面概率的似然、置信区间与顺序检验所需的停止规则。
\end{enumerate}
\subsection*{笔试推导（来源映射替换）}
\begin{enumerate}[nosep,leftmargin=*]
\item \MFQInterviewMeta{笔试推导}{10 分钟}{中}{条件概率与无放回抽样}{GB-4-25 Aces}从洗匀牌堆依次翻牌，已知前两张至少有一张 A，求两张都是 A 的概率；分别比较“至少一张”与“指定第一张是 A”的条件事件。
\end{enumerate}
\subsection*{数值编程（来源映射替换）}
\begin{enumerate}[nosep,leftmargin=*]
\item \MFQInterviewMeta{数值编程}{12 分钟}{中}{Poisson 流量与暴露时间}{GB-4-28 Cars on road}假设车辆按强度 $\lambda$ 的 Poisson 过程到达长度为 $L$ 的路段、速度固定为 $v$。推导任意时刻路段车辆数分布并模拟核验；再指出拥堵使速度依赖密度时为何模型失效。
\end{enumerate}
\subsection*{研究判断（来源映射替换）}
\begin{enumerate}[nosep,leftmargin=*]
\item \MFQInterviewMeta{研究判断}{8 分钟}{中}{条件概率与信息更新}{GB-4-20 Monty Hall}主持人知道奖品位置且总会打开一扇空门时证明换门胜率；再审计“主持人随机开门也可直接套用 $2/3$”的结论。
\item \MFQInterviewMeta{研究判断}{10 分钟}{中}{条件生存与协议差异}{GB-4-24 Russian roulette}左轮有六个弹膛和相邻两发子弹。已知第一次扣动为空，比较“直接再扣”与“重新随机旋转后再扣”的风险；明确枪膛是否顺序前进这一机制条件。
\end{enumerate}}
\providecommand{\MFQMappedUpperChEightHints}{\par\textbf{来源映射题提示：}无偏提纯只接受 HT/TH；有限样本不能证明点假设；A 题先写清条件事件；路段车辆数利用到达率乘暴露时间；Monty Hall 与轮盘题都要把主持或机械协议写进似然。}
\providecommand{\MFQMappedUpperChEightSolutions}{
\section*{来源映射替换题}
\begin{enumerate}[leftmargin=*]
\item \textbf{GB-4-16。}成对抛掷，HT 输出 1、TH 输出 0，HH/TT 丢弃。两种接受序列概率同为 $p(1-p)$，故无偏；每对接受概率 $2p(1-p)$，期望对数为其倒数，期望抛数 $1/[p(1-p)]$。独立且固定 $p$ 是必要条件。
\item \textbf{GB-4-15。}若观察 $h$ 次正面、$t$ 次反面，似然为 $L(p)=p^h(1-p)^t$，极大似然估计 $\hat p=h/(h+t)$。任何有限序列在所有 $0<p<1$ 下都有正概率，故不能逻辑证明 $p=1/2$；只能给区间或在预先规定的显著性与停止规则下检验。边看边停会改变错误率。
\item \textbf{GB-4-25。}两张无序抽样中，条件事件“至少一张 A”有 $\binom{52}{2}-\binom{48}{2}=198$ 种，二者均 A 有 $\binom42=6$ 种，故概率 $6/198=1/33$。若指定第一张是 A，则第二张为 A 的概率为 $3/51=1/17$；两个条件事件不能混用。
\item \textbf{GB-4-28。}\par 一辆车停留时间为 $L/v$。稳态下，路段车辆数等于过去 $L/v$ 时间内的到达数，故服从 $\mathrm{Poisson}(\lambda L/v)$。模拟应从稳态开始或丢弃预热期。若速度随密度变化，停留时间与车辆数相互依赖，简单的无限服务器结论失效。
\item \textbf{GB-4-20。}初选中奖概率 $1/3$，未选两门合计 $2/3$；知情主持人必排除其中空门，所以换到剩余门继承 $2/3$。若主持人随机开门且可能揭奖，观察到空门的似然随奖品位置不同，必须用 Bayes 公式重新条件化。
\item \textbf{GB-4-24。}两发相邻时共有四个空膛位置。已知当前为空后，这四种位置等可能，其中只有一个位置的下一膛是子弹，故直接再扣风险为 $1/4$；重新均匀旋转后风险回到 $2/6=1/3$，所以不旋转更安全。若弹膛不是确定前进一步或初始旋转不均匀，需重写条件样本空间。
\end{enumerate}\MFQInterviewFollowup{GB-4-15：序贯检验如何控制可选停止？GB-4-16：偏置 $p$ 随时间漂移会破坏哪一步？GB-4-20：主持人策略未知时需要估计什么似然？GB-4-24：三发相邻时结论如何变化？GB-4-25：若抽样信息来自“随机告诉你一张 A”会怎样？GB-4-28：随机速度如何改变占用分布？}}

\providecommand{\MFQMappedUpperChNineQuestions}{
\subsection*{研究判断（来源映射替换）}
\begin{enumerate}[nosep,leftmargin=*]
\item \MFQInterviewMeta{研究判断}{12 分钟}{中}{次序统计量与假设边界}{GB-4-40 Expected maximum and minimum}若 $X_1,\ldots,X_n$ 独立均匀于 $[0,1]$，推导最大值与最小值的分布和期望；再说明样本非独立时为何不能直接复用。
\item \MFQInterviewMeta{笔试推导}{10 分钟}{中}{正态矩与递推}{GB-4-32 Normal moments}对 $Z\sim N(0,1)$，用分部积分证明 $E[Z^{2k}]=(2k-1)E[Z^{2k-2}]$，并解释奇数阶矩为何为零。
\item \MFQInterviewMeta{研究判断}{12 分钟}{难}{相关界与可达性}{GB-4-41 Correlation maximum and minimum}给定随机变量 $X$ 的边际分布，讨论 $\mathrm{Corr}(X,g(X))$ 的上下界；说明 Cauchy--Schwarz 的 $[-1,1]$ 只是普适界，何时才能取到等号。
\end{enumerate}}
\providecommand{\MFQMappedUpperChNineHints}{\par\textbf{来源映射题提示：}$P(M\le x)=x^n$；正态矩把 $z\phi(z)$ 写成 $-\phi'(z)$；相关取等要求中心化变量几乎处处线性相关。}
\providecommand{\MFQMappedUpperChNineSolutions}{
\section*{来源映射替换题}
\begin{enumerate}[leftmargin=*]
\item \textbf{GB-4-40。}$F_M(x)=x^n$，故密度 $nx^{n-1}$、$E[M]=n/(n+1)$。$P(m>x)=(1-x)^n$，所以 $E[m]=\int_0^1(1-x)^n dx=1/(n+1)$。仿射变换给 $a+(b-a)$ 乘相应结果；独立同分布是推导前提。
\item \textbf{GB-4-32。}标准正态密度满足 $\phi'(z)=-z\phi(z)$。分部积分得 $E[Z^{2k}]=(2k-1)E[Z^{2k-2}]$，从 $E[1]=1$ 递推为 $(2k-1)!!$；奇次幂乘偶密度是奇函数，积分为零。非中心正态应先标准化或用矩母函数。
\item \textbf{GB-4-41。}Cauchy--Schwarz 给 $|\mathrm{Corr}(X,g(X))|\le1$；等号当且仅当 $g(X)-Eg(X)=a(X-EX)$ 几乎处处。若 $g$ 被限制为非线性、单调或离散决策类，可达界通常更窄，必须在允许函数类与边际分布下重新优化，不能把普适界当成可实现值。
\end{enumerate}\MFQInterviewFollowup{GB-4-32：如何由矩母函数复核偶数矩？GB-4-40：样本有厚尾时，极值渐近要改用哪类条件？GB-4-41：限制 $g$ 为指标函数后如何求最优阈值？}}

\providecommand{\MFQMappedUpperChElevenQuestions}{
\subsection*{研究判断（来源映射替换）}
\begin{enumerate}[nosep,leftmargin=*]
\item \MFQInterviewMeta{研究判断}{15 分钟}{难}{模式等待与 Markov 状态}{GB-5-03 Coin triplets}反复抛公平硬币，比较模式 HHT 与 HTH 的首次出现等待时间；构造前缀状态，并审计“长度相同所以等待时间相同”的结论。
\item \MFQInterviewMeta{笔试推导}{12 分钟}{中}{吸收概率与差分方程}{GB-5-01 Gambler's ruin II}资本在 $0$ 与 $N$ 之间作步长 $\pm1$ 的随机游走，向上概率为 $p$。推导从 $i$ 出发先到 $N$ 的概率，并单独处理 $p=1/2$。
\item \MFQInterviewMeta{数值编程}{12 分钟}{中}{首达时间与状态截断}{GB-5-05 Drunk man}随机游走者从整数点 $i$ 出发，在 $0,N$ 吸收。建立期望吸收时间的线性方程并用三对角求解器计算；与 Monte Carlo 比较时报告置信误差。
\item \MFQInterviewMeta{研究判断}{12 分钟}{中}{排队状态与独立假设}{GB-5-07 Ticket line}售票队列中到达与服务均被建模为 Poisson/指数。写出出生--死亡状态，并说明只匹配均值为何不足以支持等待时间结论。
\item \MFQInterviewMeta{笔试推导}{12 分钟}{中}{动态规划与系列赛}{GB-5-11 World series}两队进行先赢四场的系列赛，每场 A 胜率固定为 $p$ 且独立。写出从比分 $(a,b)$ 出发 A 最终夺冠的递推和边界，并给出 $p=1/2$ 时开赛前概率。
\end{enumerate}}
\providecommand{\MFQMappedUpperChElevenHints}{\par\textbf{来源映射题提示：}模式状态取仍可延伸目标的最长后缀；破产概率解二阶差分；吸收时间右端多一个常数项；排队先核验 Markov/无记忆假设；系列赛以当前胜场为状态。}
\providecommand{\MFQMappedUpperChElevenSolutions}{
\section*{来源映射替换题}
\begin{enumerate}[leftmargin=*]
\item \textbf{GB-5-03。}对每个模式建立前缀自动机，状态为已匹配长度 $0,1,2,3$，吸收态为 3；对 $i<3$ 写 $E_i=1+(E_{T(i,H)}+E_{T(i,T)})/2$。求解得到 HHT 的期望等待 $8$，HTH 的期望等待 $10$。差异来自 HTH 失败后保留的前缀结构，不能用“长度都是三”断言相同。
\item \textbf{GB-5-01。}令 $h_i=P_i(\tau_N<\tau_0)$，则 $h_i=ph_{i+1}+(1-p)h_{i-1}$，$h_0=0,h_N=1$。若 $p\ne1/2$，$h_i=[1-((1-p)/p)^i]/[1-((1-p)/p)^N]$；公平时取极限得 $i/N$。边界和步长改变时不能套用闭式。
\item \textbf{GB-5-05。}期望吸收时间满足 $e_i=1+pe_{i+1}+(1-p)e_{i-1}$、$e_0=e_N=0$，形成三对角线性系统。公平时闭式为 $e_i=i(N-i)$。模拟均值应附标准误或区间；厚尾的停止时间会使有限样本收敛看起来很慢。
\item \textbf{GB-5-07。}状态可取系统人数 $n$，转移率为 $n\to n+1$ 的 $\lambda$ 与 $n\to n-1$ 的 $\mu$（$n>0$）。稳态要求 $\rho=\lambda/\mu<1$。若到达有聚集、服务时长非指数或服务能力随时段变化，均值相同也会产生不同等待尾部。
\item \textbf{GB-5-11。}令 $V(a,b)$ 为 A 夺冠概率，边界 $V(4,b)=1,V(a,4)=0$，递推 $V(a,b)=pV(a+1,b)+(1-p)V(a,b+1)$。$p=1/2$ 且开赛前双方对称，故 $V(0,0)=1/2$。若胜率随主客场或比分变化，状态必须扩充。
\end{enumerate}\MFQInterviewFollowup{GB-5-01：加入每步不同赌注后状态如何改变？GB-5-03：如何由前缀自动机直接生成转移矩阵并比较模式竞赛胜率？GB-5-05：如何避免无限循环模拟拖垮测试？GB-5-07：怎样从数据检验指数服务假设？GB-5-11：七局赛改成动态胜率时如何后向计算？}}

\providecommand{\MFQMappedUpperThirteenQuestions}{
\subsection*{研究判断（来源映射替换）}
\begin{enumerate}[nosep,leftmargin=*]
\item \MFQInterviewMeta{研究判断}{12 分钟}{中}{约束与乘子解释}{GB-3-10 Lagrange multipliers}最小化 $w^T\Sigma w$，约束 $\mathbf1^Tw=1$ 与 $\mu^Tw=r_0$。写出 KKT 线性系统，并审计“求解器成功即表示乘子和经济约束都可识别”的结论。
\end{enumerate}}
\providecommand{\MFQMappedUpperThirteenHints}{\par\textbf{来源映射题提示：}把两个等式约束矩阵记为 $A^Tw=b$，联立 $2\Sigma w+A\lambda=0$。}
\providecommand{\MFQMappedUpperThirteenSolutions}{
\section*{来源映射替换题}
\begin{enumerate}[leftmargin=*]
\item \textbf{GB-3-10。}KKT 系统为 $\left(\begin{smallmatrix}2\Sigma&A\\A^T&0\end{smallmatrix}\right)(w,\lambda)^T=(0,b)^T$，其中 $A=[\mathbf1,\mu]$、$b=(1,r_0)^T$。若 $A$ 列相关或与 $\Sigma$ 的零空间形成退化，约束资格或严格凸性失效；程序返回一个解不证明经济约束可识别。
\end{enumerate}\MFQInterviewFollowup{GB-3-10：加入不等式限制后哪些 KKT 条件需要新增？}}

\providecommand{\MFQMappedUpperFourteenQuestions}{
\subsection*{数值编程与研究判断（来源映射替换）}
\begin{enumerate}[nosep,leftmargin=*]
\item \MFQInterviewMeta{数值编程}{12 分钟}{中}{Newton 收敛域}{GB-3-09 Newton's method}用 Newton 法求 $x^3-2x-2=0$，记录不同初值的迭代；给出停止准则，并构造一个进入循环或远离根的初值。
\item \MFQInterviewMeta{笔试推导}{12 分钟}{中}{有限差分误差}{GB-7-15 Finite difference}从 Taylor 展开推导中心差分的一阶导数公式与 $O(h^2)$ 截断误差，再解释 $h$ 过小时舍入误差为何上升。
\item \MFQInterviewMeta{数值编程}{15 分钟}{中}{Monte Carlo 误差预算}{GB-7-14 Monte Carlo}用 Monte Carlo 估计 $E[e^Z]$（$Z\sim N(0,1)$），给出独立理论值、标准误和置信区间；再用 antithetic variates 比较方差而不是只比较单次误差。
\end{enumerate}}
\providecommand{\MFQMappedUpperFourteenHints}{\par\textbf{来源映射题提示：}Newton 需同时检查残差与步长；中心差分把两侧 Taylor 展开相减；Monte Carlo 先用矩母函数给独立 oracle，再比较配对样本的方差。}
\providecommand{\MFQMappedUpperFourteenSolutions}{
\section*{来源映射替换题}
\begin{enumerate}[leftmargin=*]
\item \textbf{GB-3-09。}$x_{k+1}=x_k-(x_k^3-2x_k-2)/(3x_k^2-2)$。停止应同时要求 $|f(x_k)|$ 和相对步长足够小，并设最大迭代；导数接近零或初值落在非吸引区域时可能跳远/循环。面试追问不应以某个库的 success 标志替代理论收敛域。
\item \textbf{GB-7-15。}两式相减得 $[f(x+h)-f(x-h)]/(2h)=f'(x)+h^2f^{(3)}(\xi)/6$。截断误差随 $h^2$ 降低，但两近数相减的浮点绝对误差约为 $\epsilon|f|$，除以 $h$ 后放大为 $O(\epsilon/h)$，因此存在最优中间尺度。
\item \textbf{GB-7-14。}正态矩母函数给 $E[e^Z]=e^{1/2}$，且 $\mathrm{Var}(e^Z)=e^2-e$。独立样本均值的标准误为样本标准差除以 $\sqrt n$。反变量配对使用 $[e^Z+e^{-Z}]/2$；应在相同随机预算下比较配对估计量的样本方差和区间宽度，而不是挑一次更接近真值的结果。
\end{enumerate}\MFQInterviewFollowup{GB-3-09：如何用阻尼或信赖域避免跳远？GB-7-14：控制变量与重要性抽样分别需要什么额外结构？GB-7-15：如何从截断与舍入两项估计最佳 $h$？}}

\providecommand{\MFQMappedDerivativesStochasticQuestions}{
\section*{研究判断（来源映射替换）}
\begin{enumerate}[leftmargin=*]
\item \MFQInterviewMeta{研究判断}{15 分钟}{难}{停止时与反射原理}{GB-5-15 Stopping and first passage}对标准 Brownian 运动推导首次越过 $a>0$ 在 $T$ 前发生的概率，并审计直接将无界停止时代入鞅等式的证明。
\end{enumerate}}
\providecommand{\MFQMappedDerivativesStochasticHints}{\par\textbf{来源映射题提示：}用反射原理把越界后终点低于 $a$ 的路径映到终点高于 $a$。}
\providecommand{\MFQMappedDerivativesStochasticSolutions}{\section*{来源映射替换题}\begin{enumerate}[leftmargin=*]\item \textbf{GB-5-15。}反射原理给 $P(\tau_a\le T)=2P(W_T\ge a)=2[1-\Phi(a/\sqrt T)]$。若改用可选停止证明，必须先对停止时截断并验证一致可积或相应有界条件；直接把无界停止时代入鞅等式是常见错误。\end{enumerate}\MFQInterviewFollowup{GB-5-15：若边界改为随时间变化的 $a(t)$，反射证明的哪一步不再直接成立？}}

\providecommand{\MFQMappedDerivativesNumericsQuestions}{
\section*{研究判断（来源映射替换）}
\begin{enumerate}[leftmargin=*]
\item \MFQInterviewMeta{研究判断}{10 分钟}{中}{无套利复制}{GB-6-02 Put-call parity}在含连续股息率 $q$ 的市场推导欧式看涨--看跌平价，并审计一个忽略 bid/ask、融资、借券与结算时点便声称可套利的报告。
\end{enumerate}}
\providecommand{\MFQMappedDerivativesNumericsHints}{\par\textbf{来源映射题提示：}比较 $C+Ke^{-rT}$ 与 $P+S_0e^{-qT}$ 的终值。}
\providecommand{\MFQMappedDerivativesNumericsSolutions}{\section*{来源映射替换题}\begin{enumerate}[leftmargin=*]\item \textbf{GB-6-02。}两组合到期均为 $\max(S_T,K)$，故 $C-P=S_0e^{-qT}-Ke^{-rT}$。若左侧过高则卖出昂贵组合、买入便宜组合；真实可执行性还受 bid/ask、融资、股息预测、借券、行权与结算时点约束。\end{enumerate}\MFQInterviewFollowup{GB-6-02：美式期权或随机股息下，复制组合的哪个等式需要重写？}}

\providecommand{\MFQMappedPortfolioOptimizationQuestions}{
\section*{来源映射替换题}
\begin{enumerate}[leftmargin=*]
\item \MFQInterviewMeta{笔试推导}{15 分钟}{难}{均值方差与估计误差}{GB-6-14 Portfolio optimization}推导预算约束下的最小方差权重；随后说明均值方差解为何会放大预期收益估计误差，并提出至少两种稳健化方案。
\end{enumerate}}
\providecommand{\MFQMappedPortfolioOptimizationHints}{\par\textbf{来源映射题提示：}先解 $\min w^T\Sigma w$、$\mathbf1^Tw=1$；含均值时观察 $\Sigma^{-1}\hat\mu$。}
\providecommand{\MFQMappedPortfolioOptimizationSolutions}{\section*{来源映射替换题}\begin{enumerate}[leftmargin=*]\item \textbf{GB-6-14。}$w^*=\Sigma^{-1}\mathbf1/(\mathbf1^T\Sigma^{-1}\mathbf1)$。含预期收益的解依赖 $\Sigma^{-1}\hat\mu$，小特征方向和噪声均值会被放大。可用协方差收缩、均值不确定集、权重/换手约束或 Bayesian shrinkage，并必须重新比较样本外风险和成本。\end{enumerate}\MFQInterviewFollowup{GB-6-14：禁止卖空后闭式解为何失效，应怎样识别活跃约束？}}

\providecommand{\MFQMappedPortfolioTailQuestions}{
\section*{来源映射替换题}
\begin{enumerate}[leftmargin=*]
\item \MFQInterviewMeta{口述概念}{10 分钟}{中}{VaR 的性质与反例}{GB-6-15 Value at Risk}解释历史 VaR、参数 VaR 与 ES 的估计对象；构造一个 VaR 不次可加的简单离散损失例，并说明这对风险限额意味着什么。
\end{enumerate}}
\providecommand{\MFQMappedPortfolioTailHints}{\par\textbf{来源映射题提示：}让两个头寸各自以小概率损失、且损失事件互斥；单项 VaR 可为零而组合 VaR 为正。}
\providecommand{\MFQMappedPortfolioTailSolutions}{\section*{来源映射替换题}\begin{enumerate}[leftmargin=*]\item \textbf{GB-6-15。}历史法取经验损失分位数，参数法在分布假设下求分位数，ES 是越过分位点的尾部平均。令两个头寸各以 $4\%$ 概率损失 1，事件互斥；在 $95\%$ 水平各自 VaR 为 0，而组合有 $8\%$ 概率损失 1，VaR 为 1，故不次可加。限额系统不能把分散收益当作由 VaR 自动保证。\end{enumerate}\MFQInterviewFollowup{GB-6-15：把置信度提到 $99\%$ 时，需要多少尾部观测才能支撑可辩护的 ES？}}

\MFQMappedDerivativesNumericsHints

\MFQLead{对冲失效与应用分析提示}
\begin{enumerate}
  \item 从 P\&L 的二阶展开读剩余项。
  \item 一条路径没有抽样分布。
  \item 每期先融资，再交易，并把成本从现金扣除。
  \item bias 看方向，RMSE 看平方损失，分位数看尾部形状。
  \item 固定随机数后改变调仓频率。
  \item 在同一个时间步同时修改现货与收益。
  \item 比较一条路径的误差与多条随机路径的误差分布，并考虑是否挑选了结果。
  \item 区分价格模型、报价误差、实际成交与成本各自的作用。
\end{enumerate}

\chapter{投资组合与风险管理提示}
\begin{enumerate}[leftmargin=*]
  \item 区分“满足代数定义”与“对输入扰动稳定”。
  \item 展开 $\sum_iw_i(\Sigma w)_i$。
  \item 对角情形令每个 $w_i^2\sigma_i^2$ 相等。
  \item 固定收益和目标矩阵，只改变 $\lambda$。
  \item 每次抽取连续四行，而非逐行独立抽样。
  \item 让不同协方差元素使用不同缺失模式。
  \item 比较降维收益与因子遗漏、暴露漂移。
  \item 先写“何时可知”，再写“如何计算”。
\end{enumerate}

\chapter{组合优化提示}
\begin{enumerate}[leftmargin=*]
  \item $P$ 选择观点，$q$ 给观点数值，$\Omega$ 给观点误差，$\tau$ 缩放先验均值不确定性。
  \item 把先验和似然的二次型合并并完成平方。
  \item 用 $u_s\ge L_s-z$、$u_s\ge0$ 替代正部函数。
  \item 保留每个情景损失和最坏尾部集合。
  \item 上限进入 LP bounds，也进入枚举 oracle。
  \item 检查满仓、不可交易性和 KKT 条件。
  \item 三项分别约束均值误差、风险矩阵和实施摩擦。
  \item 订单层需要目标、成交、剩余量和现金四张表。
\end{enumerate}

\chapter{尾部与压力风险提示}
\begin{enumerate}[leftmargin=*]
  \item 先判断问题是概率尾部还是条件冲击。
  \item 把尾部积分拆成完整观察与边界概率质量。
  \item 令冲击为 $sx$，把二阶损失写成 $s$ 的二次式。
  \item 直接计算 $n(1-\alpha)$，不要只看总样本量。
  \item 保持同一风险因子和估值时点，再比较重估层级。
  \item 99\% 的四点样本只有 0.04 个有效尾部观察。
  \item 它是条件阈值，不携带冲击方向的概率模型。
  \item 报告必须能从冻结 fixture、oracle 和命令重建。
\end{enumerate}

\chapter{事件与订单簿提示}
\begin{enumerate}[leftmargin=*]
  \item 一个条件于 $\mathcal F_{t^-}$，一个对历史平均。
  \item 对 $n\log\lambda-\lambda T$ 求导。
  \item 没有积分强度项，增大强度只会提高事件点项。
  \item 固定 $\mu,\beta$，让分枝比接近 1 但不越界。
  \item 每次只改变被成交订单的剩余数量，并重新检查不变量。
  \item Coinbase 字段给出的是流动性提供方。
  \item 季节、单位、初始历史、补偿子与样本外残差都应出现。
  \item 把“成交记录”与“未成交委托状态”分开。
\end{enumerate}

\chapter{执行与做市控制提示}
\begin{enumerate}[leftmargin=*]
  \item 临时项只乘本期量，永久项累积到后续价格。
  \item 逐期写成交价，再乘成交数量。
  \item 只有 $R=0$ 的终端状态可行。
  \item 枚举 $0\ldots x_{\max}$ 的笛卡尔积并筛选总和。
  \item 多头库存应让两边报价同时向下移动。
  \item 第 $t$ 次报价只能使用第 $t$ 次成交前的状态。
  \item IS 的基准、现金费用、剩余量和停止状态都不可缺。
  \item 固定同一随机数流比较不同控制器。
\end{enumerate}

\MFQTransition{事件时间仿真提示}
\begin{enumerate}[leftmargin=*]
  \item 两种策略路径的协方差越高，差值方差通常越低。
  \item 点差收入与持仓面对下一次价格变化的损益分开写。
  \item 路径为重采样单位，不能把单一路径内事件当 IID 样本。
  \item 用累计强度把事件时刻映射到近似单位指数间隔。
  \item 在改变状态前验证顺序、唯一编号和最优买卖价。
  \item 补库存、成交、费用、冲击、停止和路径分布。
  \item 成交只说明别人发生了交易，不显示你的未成交委托位置。
  \item 先完成确定性小簿，再增加随机性和真实输入。
\end{enumerate}


\chapter{共享符号表}
\begin{center}
\begin{tabular}{lll}
\toprule
符号 & 含义 & 首次正式教学 \\
\midrule
$r_t$ & 时点 $t$ 的简单收益率 & 上册第 1 章 \\
$r_i$ & 资产 $i$ 的简单收益率 & 上册第 1 章 \\
$w_i$ & 资产 $i$ 的组合权重 & 上册第 1 章 \\
$R_p$ & 组合 $p$ 的简单收益率 & 上册第 1 章 \\
$\mathbb{E}[X]$ & 随机变量 $X$ 的期望 & 上册第 7 章 \\
\bottomrule
\end{tabular}
\end{center}

机器可读的权威注册位于课程清单引用的 notation registry；本表是供读者阅读的出版视图。


\chapter{共享术语表}
\begin{description}
  \item[学习单元（learning unit）] 围绕一个可观察学习产出组织的最小自学闭环。
  \item[独立 oracle（independent oracle）] 不由被测实现计算得出的期望结果。
  \item[命题（proposition）] 具有确定真假值、需要证明或反驳的陈述。
  \item[必要条件（necessary condition）] 结论成立时必须满足的条件。
  \item[充分条件（sufficient condition）] 一旦满足即可推出结论的条件。
  \item[反例（counterexample）] 满足命题假设但使结论失败的合法对象。
  \item[量纲（dimension）] 数量所属的单位类别，用于判断运算是否有意义。
  \item[度量空间（metric space）] 配备满足分离、对称和三角不等式之距离函数的集合。
  \item[Cauchy 序列（Cauchy sequence）] 尾部任意两项都能任意接近的序列。
  \item[完备性（completeness）] 每个 Cauchy 序列都收敛到空间内一点的性质。
  \item[一致收敛（uniform convergence）] 同一索引阈值同时控制定义域所有点误差的收敛。
  \item[压缩映射（contraction mapping）] 以严格小于 1 的统一因子缩短任意两点距离的映射。
  \item[$\sigma$ 代数（sigma algebra）] 包含全集并对补集和可数并封闭的事件集合族。
  \item[测度（measure）] 满足可数可加性的非负集合函数。
  \item[可测函数（measurable function）] 目标可测集的逆像仍是可测事件的函数。
  \item[简单函数（simple function）] 只取有限多个值的可测函数。
  \item[Lebesgue 积分（Lebesgue integral）] 由简单函数从下逼近构造的积分。
  \item[几乎处处（almost everywhere）] 除一个零测集外性质均成立。
  \item[$L^p$ 空间（Lp space）] $p$ 次绝对可积函数按几乎处处相等取等价类得到的完备范数空间。
  \item[本质上确界（essential supremum）] 忽略零测集后成立的最小几乎处处上界。
  \item[Hölder 不等式（Holder inequality）] 用共轭 $L^p$ 范数控制乘积积分的不等式。
  \item[Minkowski 不等式（Minkowski inequality）] 保证 $L^p$ 范数满足三角不等式的积分不等式。
  \item[Tonelli 定理（Tonelli theorem）] 对非负可测函数交换积分次序的定理。
  \item[Fubini 定理（Fubini theorem）] 在绝对可积条件下交换积分次序的定理。
  \item[乘积测度（product measure）] 从可测矩形上的测度乘积扩张得到的测度。
  \item[多元可微（multivariable differentiability）] 函数增量具有统一线性近似且相对余项消失。
  \item[梯度（gradient）] 表示标量函数一阶线性近似的偏导向量。
  \item[Jacobian] 表示向量值函数一阶线性近似的偏导矩阵。
  \item[Hessian] 描述标量函数局部二阶曲率的方阵。
  \item[矩阵微分（matrix differential）] 用线性微分和形状账本推导导数的记法。
  \item[正交投影（orthogonal projection）] 使残差位于目标子空间正交补的线性映射。
  \item[最小二乘（least squares）] 最小化残差平方和的估计方法。
  \item[谱分解（spectral decomposition）] 实对称矩阵的正交特征分解。
  \item[奇异值分解（singular value decomposition）] 任意矩阵的左右正交方向与奇异值分解。
  \item[条件数（condition number）] 问题对输入相对扰动的最坏敏感度。
  \item[概率空间（probability space）] 样本空间、事件 $\sigma$ 代数与概率测度组成的三元组。
  \item[随机变量（random variable）] 从概率空间到目标可测空间的可测映射。
  \item[推前分布（pushforward distribution）] 随机变量把底层概率测度推到目标空间得到的分布。
  \item[联合分布（joint distribution）] 随机向量在乘积空间上的分布。
  \item[边缘分布（marginal distribution）] 对联合分布的其余坐标积分或求和得到的分布。
  \item[重尾分布（heavy-tailed distribution）] 尾部衰减较慢并可能使某些矩不存在的分布。
  \item[条件期望（conditional expectation）] 给定子 $\sigma$ 代数时保持其事件上积分的可测随机变量。
  \item[塔式性质（tower property）] 对嵌套信息依次条件化等于直接对较粗信息条件化。
  \item[独立性（independence）] 联合概率对所有可测集合分解为边缘概率乘积。
  \item[条件独立（conditional independence）] 给定指定信息后联合条件分布分解为条件边缘乘积。
  \item[概率核（probability kernel）] 对状态可测、对事件为概率测度的条件分布映射。
  \item[几乎处处收敛（almost sure convergence）] 除零概率状态外逐点收敛。
  \item[依概率收敛（convergence in probability）] 偏差超过任意正阈值的概率趋零。
  \item[$L^p$ 收敛（Lp convergence）] 差的 $p$ 次绝对矩趋零。
  \item[分布收敛（convergence in distribution）] 分布函数在极限分布连续点收敛。
  \item[大数定律（law of large numbers）] 在明确矩与依赖条件下样本平均趋近总体均值。
  \item[中心极限定理（central limit theorem）] 标准化和在正则条件下趋近正态分布。
  \item[集中不等式（concentration inequality）] 控制有限样本偏离概率的非渐近上界。
  \item[估计量（estimator）] 把样本映射为参数估计的随机规则。
  \item[一致性（consistency）] 估计量依概率收敛到目标参数的性质。
  \item[渐近分布（asymptotic distribution）] 估计误差适当中心化与缩放后的极限分布。
  \item[异方差稳健标准误（heteroskedasticity-robust standard error）] 用三明治协方差估计且不要求条件误差方差恒定的标准误。
  \item[族错误率（family-wise error rate）] 一族真零假设中至少错误拒绝一个的概率。
  \item[错误发现率（false discovery rate）] 拒绝集合中错误发现比例的期望控制目标。
  \item[随机过程（stochastic process）] 同一概率空间上由索引参数化的随机变量族。
  \item[滤过（filtration）] 随时间递增、表示可用信息增长的 $\sigma$ 代数族。
  \item[Markov 性（Markov property）] 给定当前状态后未来条件分布不再依赖更早历史。
  \item[鞅（martingale）] 下一时点条件期望等于当前值的可积适应过程。
  \item[Poisson 过程（Poisson process）] 具有独立平稳 Poisson 增量的计数过程。
  \item[Brownian 运动（Brownian motion）] 具有连续路径及独立平稳高斯增量的过程。
  \item[二次变差（quadratic variation）] 沿细分割对路径增量平方求和的极限。
  \item[桥接模块（bridge module）] 根据诊断结果精确指回上册先修单元的短模块。
  \item[Capstone] 由干净复现、独立基线、失败案例和限制声明组成的综合能力证据。
\end{description}


\chapter{下册主题索引}
% Generated from curriculum/manifest.json; do not edit by hand.
\begin{longtable}{p{0.30\textwidth}p{0.60\textwidth}}
\toprule
主题 & 正式教学单元 \\
\midrule
\endfirsthead
\toprule
主题 & 正式教学单元 \\
\midrule
\endhead
Hawkes 过程 & 事件与订单簿：从条件强度到队列成交 \\
Kalman 滤波 & 估计与状态推断：滤波、切换与变点 \\
Monte Carlo & 定价与校准：从闭式、树和 PDE 到参数化曲面 \\
不变量 & Brainteaser：把陌生题压缩成可证明结构 \\
临时冲击 & 执行与做市：价格冲击与库存反馈 \\
二次变差 & 随机分析与无套利：从二次变差到风险中性测度；定价与校准：从闭式、树和 PDE 到参数化曲面 \\
交易成本 & 多因子研究：组合收益、归因与交易成本；统计套利研究：从预测到净收益；Alpha 决策：从预测损失到交易收益；离散对冲：Greeks、交易成本与误差分布；组合优化：从观点后验到稳健实施 \\
价格时间优先 & 事件与订单簿：从条件强度到队列成交 \\
协整 & 动态与长期关系：从单位根到 OU \\
协方差矩阵 & 风险估计：从协方差到风险预算 \\
可复现研究包 & 多因子研究：组合收益、归因与交易成本；统计套利研究：从预测到净收益；Alpha 决策：从预测损失到交易收益；离散对冲：Greeks、交易成本与误差分布；尾部与压力风险：从有效样本到反向压力 \\
因子风险模型 & 风险估计：从协方差到风险预算 \\
多重检验 & 多因子模型：预测回归与风险价格 \\
实现差额 & 执行与做市：价格冲击与库存反馈；事件时间仿真与综合练习 \\
容量 & 多因子研究：组合收益、归因与交易成本 \\
换手率 & 多因子研究：组合收益、归因与交易成本；组合优化：从观点后验到稳健实施 \\
数值稳定性 & 多因子估计：正则化、数值实现与外部数据 \\
数据泄漏 & 多因子模型：预测回归与风险价格；多因子估计：正则化、数值实现与外部数据；表格统计学习：从决策树到梯度提升；验证与监控：嵌套时序、校准、解释与漂移；Alpha 决策：从预测损失到交易收益 \\
时点数据 & 多因子模型：预测回归与风险价格；多因子估计：正则化、数值实现与外部数据 \\
条件强度 & 事件与订单簿：从条件强度到队列成交；事件时间仿真与综合练习 \\
样本外 & 表格统计学习：从决策树到梯度提升 \\
滚动向前验证 & 动态与长期关系：从单位根到 OU；统计套利研究：从预测到净收益；验证与监控：嵌套时序、校准、解释与漂移 \\
状态空间模型 & 估计与状态推断：滤波、切换与变点 \\
鞅 & 随机分析与无套利：从二次变差到风险中性测度 \\
预期损失 & 尾部与压力风险：从有效样本到反向压力 \\
风险价值 & 尾部与压力风险：从有效样本到反向压力 \\
\bottomrule
\end{longtable}


\chapter{补充习题来源}
% Generated by tools/build_course_assets.py; edit curriculum/interview-problem-ledger.json instead.
补充习题参考 Xinfeng Zhou 的《A Practical Guide to Quantitative Finance Interviews》（2008）。下表列出原题名及本书位置。题目按本书的符号与教学顺序重新表述，解答独立编写；题源用于进一步阅读与对照。
\begin{longtable}{p{0.14\textwidth}p{0.38\textwidth}p{0.34\textwidth}}
编号 & 原书题名 & 本书位置 \\ \hline
GB-2-01 & Screwy pirates & 下册《Brainteaser：把陌生题压缩成可证明结构》 \\
GB-2-02 & Tiger and sheep & 下册《Brainteaser：把陌生题压缩成可证明结构》 \\
GB-2-03 & River crossing & 下册《Brainteaser：把陌生题压缩成可证明结构》 \\
GB-2-04 & Birthday problem & 下册《Brainteaser：把陌生题压缩成可证明结构》 \\
GB-2-05 & Card game & 下册《Brainteaser：把陌生题压缩成可证明结构》 \\
GB-2-06 & Burning ropes & 下册《Brainteaser：把陌生题压缩成可证明结构》 \\
GB-2-07 & Defective ball & 下册《Brainteaser：把陌生题压缩成可证明结构》 \\
GB-2-08 & Trailing zeros & 下册《Brainteaser：把陌生题压缩成可证明结构》 \\
GB-2-09 & Horse race & 下册《Brainteaser：把陌生题压缩成可证明结构》 \\
GB-2-10 & Infinite sequence & 下册《Brainteaser：把陌生题压缩成可证明结构》 \\
GB-2-11 & Box packing & 下册《Brainteaser：把陌生题压缩成可证明结构》 \\
GB-2-12 & Calendar cubes & 下册《Brainteaser：把陌生题压缩成可证明结构》 \\
GB-2-13 & Door to offer & 下册《Brainteaser：把陌生题压缩成可证明结构》 \\
GB-2-14 & Message delivery & 下册《Brainteaser：把陌生题压缩成可证明结构》 \\
GB-2-15 & Last ball & 下册《Brainteaser：把陌生题压缩成可证明结构》 \\
GB-2-16 & Light switches & 下册《Brainteaser：把陌生题压缩成可证明结构》 \\
GB-2-17 & Quant salary & 下册《Brainteaser：把陌生题压缩成可证明结构》 \\
GB-2-18 & Coin piles & 下册《Brainteaser：把陌生题压缩成可证明结构》 \\
GB-2-19 & Mislabeled bags & 下册《Brainteaser：把陌生题压缩成可证明结构》 \\
GB-2-20 & Wise men & 下册《Brainteaser：把陌生题压缩成可证明结构》 \\
GB-2-21 & Clock pieces & 下册《Brainteaser：把陌生题压缩成可证明结构》 \\
GB-2-22 & Missing integers & 下册《Brainteaser：把陌生题压缩成可证明结构》 \\
GB-2-23 & Counterfeit coins I & 下册《Brainteaser：把陌生题压缩成可证明结构》 \\
GB-2-24 & Matching socks & 下册《Brainteaser：把陌生题压缩成可证明结构》 \\
GB-2-25 & Handshakes & 下册《Brainteaser：把陌生题压缩成可证明结构》 \\
GB-2-26 & Have we met before? & 下册《Brainteaser：把陌生题压缩成可证明结构》 \\
GB-2-27 & Ants on a square & 下册《Brainteaser：把陌生题压缩成可证明结构》 \\
GB-2-28 & Counterfeit coins II & 下册《Brainteaser：把陌生题压缩成可证明结构》 \\
GB-2-29 & Prisoner problem & 下册《Brainteaser：把陌生题压缩成可证明结构》 \\
GB-2-30 & Division by 9 & 下册《Brainteaser：把陌生题压缩成可证明结构》 \\
GB-2-31 & Chameleon colors & 下册《Brainteaser：把陌生题压缩成可证明结构》 \\
GB-2-32 & Coin split problem & 下册《Brainteaser：把陌生题压缩成可证明结构》 \\
GB-2-33 & Chocolate bar problem & 下册《Brainteaser：把陌生题压缩成可证明结构》 \\
GB-2-34 & Race track & 下册《Brainteaser：把陌生题压缩成可证明结构》 \\
GB-2-35 & Irrational number & 下册《Brainteaser：把陌生题压缩成可证明结构》 \\
GB-2-36 & Rainbow hats & 下册《Brainteaser：把陌生题压缩成可证明结构》 \\
GB-3-02 & Maximum and minimum & 上册第 5 章《多元微积分与矩阵微分》 \\
GB-3-08 & Taylor series & 上册第 5 章《多元微积分与矩阵微分》 \\
GB-3-09 & Newton's method & 上册第 14 章《浮点数、数值线代与病态问题》 \\
GB-3-10 & Lagrange multipliers & 上册第 13 章《凸优化、KKT 与对偶》 \\
GB-3-16 & QR decomposition & 上册第 6 章《线性代数、分解、投影与条件数》 \\
GB-3-17 & Determinant, eigenvalue and eigenvector & 上册第 6 章《线性代数、分解、投影与条件数》 \\
GB-3-18 & Positive semidefinite and positive definite matrices & 上册第 6 章《线性代数、分解、投影与条件数》 \\
GB-3-19 & LU decomposition and Cholesky decomposition & 上册第 6 章《线性代数、分解、投影与条件数》 \\
GB-4-01 & Coin toss game I & 上册第 7 章《概率空间、随机变量与分布》 \\
GB-4-03 & Drunk passenger & 上册第 7 章《概率空间、随机变量与分布》 \\
GB-4-05 & Poker hands & 上册第 7 章《概率空间、随机变量与分布》 \\
GB-4-10 & Birthday problem II & 上册第 7 章《概率空间、随机变量与分布》 \\
GB-4-15 & Unfair coin & 上册第 8 章《条件期望、独立性与概率核》 \\
GB-4-16 & Fair coin from an unfair coin & 上册第 8 章《条件期望、独立性与概率核》 \\
GB-4-20 & Monty Hall & 上册第 8 章《条件期望、独立性与概率核》 \\
GB-4-24 & Russian roulette & 上册第 8 章《条件期望、独立性与概率核》 \\
GB-4-25 & Aces & 上册第 8 章《条件期望、独立性与概率核》 \\
GB-4-28 & Cars on road & 上册第 8 章《条件期望、独立性与概率核》 \\
GB-4-29 & Meeting probability & 上册第 7 章《概率空间、随机变量与分布》 \\
GB-4-31 & Poisson process property & 上册第 7 章《概率空间、随机变量与分布》 \\
GB-4-32 & Normal moments & 上册第 9 章《概率收敛、极限定理与集中不等式》 \\
GB-4-37 & Sum of random variables & 上册第 7 章《概率空间、随机变量与分布》 \\
GB-4-38 & Coupon collection & 上册第 7 章《概率空间、随机变量与分布》 \\
GB-4-39 & Joint default & 上册第 7 章《概率空间、随机变量与分布》 \\
GB-4-40 & Expected maximum and minimum & 上册第 9 章《概率收敛、极限定理与集中不等式》 \\
GB-4-41 & Correlation maximum and minimum & 上册第 9 章《概率收敛、极限定理与集中不等式》 \\
GB-5-01 & Gambler's ruin II & 上册第 11 章《随机过程、Markov、鞅、Poisson 与 Brownian》 \\
GB-5-03 & Coin triplets & 上册第 11 章《随机过程、Markov、鞅、Poisson 与 Brownian》 \\
GB-5-05 & Drunk man & 上册第 11 章《随机过程、Markov、鞅、Poisson 与 Brownian》 \\
GB-5-07 & Ticket line & 上册第 11 章《随机过程、Markov、鞅、Poisson 与 Brownian》 \\
GB-5-11 & World series & 上册第 11 章《随机过程、Markov、鞅、Poisson 与 Brownian》 \\
GB-5-15 & Stopping and first passage & 下册《随机分析与无套利：从二次变差到风险中性测度》 \\
GB-6-02 & Put-call parity & 下册《定价与校准：从闭式、树和 PDE 到参数化曲面》 \\
GB-6-14 & Portfolio optimization & 下册《组合优化：从观点后验到稳健实施》 \\
GB-6-15 & Value at Risk & 下册《尾部与压力风险：从有效样本到反向压力》 \\
GB-7-14 & Monte Carlo & 上册第 14 章《浮点数、数值线代与病态问题》 \\
GB-7-15 & Finite difference & 上册第 14 章《浮点数、数值线代与病态问题》 \\
\end{longtable}


\chapter{参考文献}
\printbibliography[heading=none]

\chapter{许可说明}
原创代码采用 MIT License；原创书稿与图形采用 CC BY-NC-SA 4.0；ElegantBook 类文件沿用 LPPL；真实数据切片将在引入时分别记录来源与许可。

\end{document}
